\documentclass[10pt,a4paper]{article}
\usepackage{DDTA_METHODOLOGY_GUIDE_STYLE_R1}
\usepackage{paracol}
\geometry{a4paper,landscape,top=15mm,bottom=15mm,left=17mm,right=17mm,headheight=15pt}
\setlength{\parskip}{0.38em}
\titlespacing*{\section}{0pt}{1.2em}{0.45em}
\titlespacing*{\subsection}{0pt}{0.9em}{0.35em}
\titlespacing*{\subsubsection}{0pt}{0.7em}{0.25em}

\hypersetup{
  pdftitle={DDTA DermaTriage - Parallel Case Study R25 - FR03 split review},
  pdfauthor={Documentation-Driven Threat Analysis research project}
}

% -----------------------------------------------------------------------------
% PAGE-INTEGRITY HASHES
% Generated over the canonical LaTeX payload between PAGE-CONTENT markers.
% Hash lines themselves are excluded. On the integrity-index page, hashes of
% preceding pages are resolved before hashing, while the page's own hash is
% represented by the canonical token <SELF> to avoid self-reference.
% -----------------------------------------------------------------------------
\newcommand{\PageMDfive}[1]{\csname PageMDfive#1\endcsname}
%<PAGE-HASH-DEFS>
\expandafter\def\csname PageMDfive1\endcsname{377b35624c4b40d544cde72bcf00740c}
\expandafter\def\csname PageMDfive2\endcsname{a3275772d6cd2ce37195f7d8163144e8}
\expandafter\def\csname PageMDfive3\endcsname{46cbd47a5fe8e2c86417634fdddea7c8}
\expandafter\def\csname PageMDfive4\endcsname{d464848d4f632cc58be0991cacf1b647}
\expandafter\def\csname PageMDfive5\endcsname{b20178fc4e6cdc140cf16360634f5344}
\expandafter\def\csname PageMDfive6\endcsname{a3198ac668083dc67b750dff9bd44175}
\expandafter\def\csname PageMDfive7\endcsname{e50995f0de66f6d8f847ce9ddfa6912f}
\expandafter\def\csname PageMDfive8\endcsname{452c3e6e4e672ba2f0d437f6363ee29c}
\expandafter\def\csname PageMDfive9\endcsname{bb03073eb7e38e5511cf24446218dd47}
\expandafter\def\csname PageMDfive10\endcsname{cba1a94e83a61109b957b847a25c73cd}
\expandafter\def\csname PageMDfive11\endcsname{9d2c86231c54c3d67fc1a9a7c8ccbad7}
\expandafter\def\csname PageMDfive12\endcsname{986ebd679b75abfb4aba7daf6535585b}
\expandafter\def\csname PageMDfive13\endcsname{8c9bec5d00a83f1943ff65f5809fa586}
\expandafter\def\csname PageMDfive14\endcsname{9414cbe89b9c69e3d18aed9f06cfec3e}
\expandafter\def\csname PageMDfive15\endcsname{db9d71a6eeaa423ad8f10b085637530b}
\expandafter\def\csname PageMDfive16\endcsname{b79abb8c1930e6e407ffcd6795722076}
\expandafter\def\csname PageMDfive17\endcsname{e9da460bcff5c8d61a5361f9897fcaae}
\expandafter\def\csname PageMDfive18\endcsname{60b077964dc202b0e8ae17b0e9cce7d0}
\expandafter\def\csname PageMDfive19\endcsname{a54afe35ed1503c6eda41c4482b13966}
\expandafter\def\csname PageMDfive20\endcsname{5ec2d19412461c60854ca7cef75c9f94}
\expandafter\def\csname PageMDfive21\endcsname{2a7e92a9a7609d1466e545ebf818e94a}
\expandafter\def\csname PageMDfive22\endcsname{9cc601ec85fa3f86c36ea04b5fa377fa}
\expandafter\def\csname PageMDfive23\endcsname{0c4840b8de97233c6f01ca44a133875f}
\expandafter\def\csname PageMDfive24\endcsname{690b289527212e55343c18f392c6fb41}
\expandafter\def\csname PageMDfive25\endcsname{dc1c9a8641558fca5003a43e5c5d78fc}
\expandafter\def\csname PageMDfive26\endcsname{b1d93470f3d0174451208c9bd19cea66}
\expandafter\def\csname PageMDfive27\endcsname{4cc450cd755e4afafca7c0ba2a932a3a}
\expandafter\def\csname PageMDfive28\endcsname{0bc6d19506b654c3e6ef3aebd3d68a34}
\expandafter\def\csname PageMDfive29\endcsname{e87f73ed8a77121cd40e5f04b3f87658}
\expandafter\def\csname PageMDfive30\endcsname{e4232081ef5e3573bf9355c5c2394851}
\expandafter\def\csname PageMDfive31\endcsname{7ccf212862038cd02fccb29b590e7ce6}
\expandafter\def\csname PageMDfive32\endcsname{7a62b581d6ec0a4f598d0c0edbcf24fe}
\expandafter\def\csname PageMDfive33\endcsname{1226aa35875fdc1a1db3573018cb07ea}
\expandafter\def\csname PageMDfive34\endcsname{43f138821276abfcc9fc2aa76d4cc622}
\expandafter\def\csname PageMDfive35\endcsname{2d870f81a5f24c387896ef85f94eacbd}
\expandafter\def\csname PageMDfive36\endcsname{6a4807e4d18d4a97f3b02d62ad4f24e0}
\expandafter\def\csname PageMDfive37\endcsname{475d59cb9fc8bc3342897fab12ed179c}
\expandafter\def\csname PageMDfive38\endcsname{1596523e5d2a3554727d058d351afb63}
\expandafter\def\csname PageMDfive39\endcsname{9faf8d41344390a05878e91d5d4ac7bd}
\expandafter\def\csname PageMDfive40\endcsname{81091a5b8d2f1477e8e49d8ebfd5bfc8}
\expandafter\def\csname PageMDfive41\endcsname{3fb4ab29517018401c110b38685507d7}
\expandafter\def\csname PageMDfive42\endcsname{df70ba8f961369f44bdd2b8e064a53bd}
\expandafter\def\csname PageMDfive43\endcsname{a22a5ac76953eb09ad50e9b2bbe35c12}
%</PAGE-HASH-DEFS>

\newcommand{\PageHashFooter}[1]{%
  \fancyfoot[R]{\sffamily\scriptsize\color{DDTAGray}MD5 contenuto: \texttt{#1}}%
}

\newcommand{\IntegrityRow}[3]{#1 & #2 & {\footnotesize\ttfamily #3} \\
}

\newtcolorbox{ddtaquestions}[1][]{%
  breakable,enhanced,arc=2pt,boxrule=0.65pt,
  colframe=DDTAGray!65,colback=DDTALightGray,
  left=8pt,right=8pt,top=7pt,bottom=7pt,
  fonttitle=\sffamily\bfseries,title={Domande / chiarimenti richiesti},#1}

% BA authoring-state convention used during the progressive analysis.
% Typography is a state marker, not a semantic type.
\newtcolorbox{ddtareferentidentified}[1][]{%
  breakable,enhanced,arc=2pt,boxrule=0.75pt,
  colframe=DDTAAmber,colback=DDTALightAmber,
  left=7pt,right=7pt,top=5pt,bottom=5pt,
  fonttitle=\sffamily\bfseries,title={BAReferent identificati},#1}
\newtcolorbox{ddtareferentaccepted}[1][]{%
  breakable,enhanced,arc=2pt,boxrule=0.75pt,
  colframe=DDTAGreen,colback=DDTALightGreen,
  left=7pt,right=7pt,top=5pt,bottom=5pt,
  fonttitle=\sffamily\bfseries,title={BAReferent accettati},#1}
\newtcolorbox{ddtapropositionidentified}[1][]{%
  breakable,enhanced,arc=2pt,boxrule=0.75pt,
  colframe=DDTAAmber,colback=DDTALightAmber,
  left=7pt,right=7pt,top=5pt,bottom=5pt,
  fonttitle=\sffamily\bfseries,title={BAProposition identificate},#1}
\newtcolorbox{ddtapropositionaccepted}[1][]{%
  breakable,enhanced,arc=2pt,boxrule=0.75pt,
  colframe=DDTAGreen,colback=DDTALightGreen,
  left=7pt,right=7pt,top=5pt,bottom=5pt,
  fonttitle=\sffamily\bfseries,title={BAProposition accettate},#1}
\newtcolorbox{ddtaopenrelations}[1][]{%
  breakable,enhanced,arc=2pt,boxrule=0.65pt,
  colframe=DDTAGray!70,colback=DDTALightGray,
  left=7pt,right=7pt,top=5pt,bottom=5pt,
  fonttitle=\sffamily\bfseries,title={Relazioni BA aperte},#1}
\newcommand{\BACandidate}[1]{\emph{#1}}
\newcommand{\BAAccepted}[1]{\textbf{#1}}
\newcommand{\BANew}{\texttt{(+)}}
\newcommand{\BAFirstAccepted}{\texttt{(b)}}


\begin{document}
\selectlanguage{italian}
\fancyhead[L]{\sffamily\small\color{DDTAGray}DDTA - DermaTriage Case Study}
\fancyhead[R]{\sffamily\small\color{DDTABlue}Documentation + Base Analysis}
\fancyfoot[L]{\sffamily\scriptsize\color{DDTAGray}Documentation 2/3 - Base Analysis 1/3}

%<PAGE-DESC:1>Frontespizio e struttura del case study
%<PAGE-CONTENT:1>
\begin{titlepage}
\thispagestyle{empty}
\vspace*{10mm}
{\sffamily\bfseries\color{DDTANavy}\fontsize{15}{18}\selectfont Documentation-Driven Threat Analysis}\par
\vspace{5mm}
{\sffamily\bfseries\color{DDTANavy}\fontsize{29}{33}\selectfont DermaTriage Case Study}\par
\vspace{2mm}
{\sffamily\color{DDTABlue}\fontsize{18}{22}\selectfont Documentation + Base Analysis}\par

\vspace{12mm}
\begin{tcolorbox}[enhanced,arc=3pt,boxrule=0pt,colback=DDTANavy,left=13pt,right=13pt,top=12pt,bottom=12pt]
{\color{white}\sffamily\bfseries\large Documentazione gerarchica DDTA del progetto DermaTriage}\par
\vspace{2mm}
{\color{white!88}\small La documentazione di progetto occupa circa due terzi della pagina; la Base Analysis occupa il terzo laterale e viene avviata progressivamente distinguendo elementi candidati da elementi stabili / accettati. La documentazione e' ordinata gerarchicamente: MacroRequirement, quindi le sue Decision e, immediatamente sotto ogni Decision, i relativi FunctionalRequirement.}
\end{tcolorbox}

\vspace{5mm}
\begin{ddtaopen}[title={RIQUADRO TEMPORANEO - WORKING SET E RIFERIMENTI - DA CANCELLARE A CHIUSURA DEL CASE STUDY}]
{\scriptsize
\textbf{Attivi:} \nolinkurl{DDTA_DOCUMENTATION_AUTHORING_GUIDE_R7_REBUILD_R8.tex}; \nolinkurl{DDTA_BASE_ANALYSIS_GUIDE_REBUILD_R7.tex}; questo \nolinkurl{DDTA_DERMATRIAGE_PARALLEL_CASE_STUDY_R25_FR03_SPLIT_REVIEW_R1.tex}. Working plan: \nolinkurl{DDTA_R25_BASE_ANALYSIS_GUIDE_REBUILD_WORK_PLAN_R5.md}.\par
\textbf{Riferimenti BA:} \nolinkurl{DDTA_BASE_ANALYSIS_OPERATIONAL_GUIDE_R3.tex} + BA0--BA5 e \nolinkurl{DDTA_BASE_ANALYSIS_GUIDE_REBUILD_R7.tex}. \textbf{Fase corrente:} R25 revisiona \texttt{FR-MR01-03-03} separando, in review, preparazione della base storica e retrieval runtime. I label \texttt{03-03A/B} sono provvisori; i nomi BA restano nella sola colonna Base Analysis. Il necessity test su \texttt{correlate} resta tracciato nella pagina temporanea rossa.\par
{\tiny\textbf{Vincolo anti-regressione:} usare solo i campi documentali previsti dalla guida corrente e, per la BA, solo campi, metadata, operatori e ruoli gia' ammessi da R3 + BA0--BA5. Non aggiungere campi, stati o riquadri semantici non concordati: se un campo sembra necessario, fermarsi e discuterlo prima di inserirlo; le note di review restano nelle strutture gia' previste o nella working analysis. Preservare invariati i blocchi \texttt{PAGE-CONTENT} non interessati e rigenerare gli MD5 con lo script repository.}
}
\end{ddtaopen}

\vfill
\begin{ddtacore}[title={Struttura del case study}]
La lettura segue sempre il containment documentale: MR $\rightarrow$ Decision $\rightarrow$ FR. Non vengono raccolte prima tutte le Decision e poi tutti gli FR. Ogni FR descrive direttamente il funzionamento e l'implementazione corrente di DermaTriage; i riquadri grigi della documentazione sono riservati a domande e richieste di chiarimento. La colonna Base Analysis usa sempre quattro riquadri: BAReferent identificati, BAReferent accettati, BAProposition identificate e BAProposition accettate. Un eventuale riquadro grigio aggiuntivo, Relazioni BA aperte, conserva soltanto relazioni gia' riconosciute ma non ancora sufficientemente strutturate come BAProposition. Corsivo = candidato, grassetto = stabile / accettato, \texttt{(+)} = prima introduzione dell'identita' BA, \texttt{(b)} = prima promozione in grassetto. \texttt{??} identifica esclusivamente uno slot interno ancora irrisolto di una proposition gia' strutturata. \texttt{NON ANALIZZATO} indica una fase non ancora eseguita; -- indica analisi eseguita senza elementi di quel tipo. Il ragionamento di studio resta esterno al case study; il registro BA finale raccoglie la tracciabilita' complessiva.
\end{ddtacore}

\vspace{5mm}
{\sffamily\scriptsize\color{DDTAGray}MD5 contenuto pagina 1: \texttt{\PageMDfive{1}}}\par
\end{titlepage}
%</PAGE-CONTENT:1>
\clearpage
%<PAGE-DESC:2>Contesto generale del progetto - Project Problem Framing
%<PAGE-CONTENT:2>
\PageHashFooter{\PageMDfive{2}}
\section{Contesto generale del progetto}

\setlength{\columnsep}{7mm}
\setlength{\columnseprule}{0.35pt}
\columnratio{0.67}
\begin{paracol}{2}

% LEFT 2/3 - DOCUMENTATION
{\sffamily\large\bfseries\color{DDTANavy}Documentazione DDTA}\par
\vspace{0.5em}

\begin{ddtacore}[title={Project Problem Framing}]
Il problema che DermaTriage deve affrontare e' il supporto al triage precoce di casi dermatologici relativi a lesioni potenzialmente oncologiche. Il progetto parte dalle informazioni gia' disponibili sul caso per determinarne l'urgenza e una priorita' operativa di presa in carico, con l'obiettivo di favorire un instradamento specialistico tempestivo.
\end{ddtacore}

\switchcolumn

{\sffamily\large\bfseries\color{DDTABlue}Base Analysis}\par
\vspace{0.5em}

\begin{ddtareferentidentified}
\scriptsize
\BACandidate{DermaTriage} \BANew{}; \BACandidate{DermatologicalCase} \BANew{}; \BACandidate{CaseUrgency} \BANew{}; \BACandidate{OperationalTriagePriority} \BANew{}; \BACandidate{DermatologicalTriageEvaluation} \BANew{}; \BACandidate{CaseInformation} \BANew{}.\par
\end{ddtareferentidentified}

\begin{ddtareferentaccepted}
\scriptsize
--
\end{ddtareferentaccepted}

\begin{ddtapropositionidentified}
\scriptsize
\texttt{NON ANALIZZATO}
\end{ddtapropositionidentified}

\begin{ddtapropositionaccepted}
\scriptsize
\texttt{NON ANALIZZATO}
\end{ddtapropositionaccepted}

\end{paracol}
%</PAGE-CONTENT:2>
\clearpage
%<PAGE-DESC:3>MR-01 - Valutazione di triage del caso dermatologico
%<PAGE-CONTENT:3>
\PageHashFooter{\PageMDfive{3}}
\section{MR-01 - Valutazione di triage del caso dermatologico}

\setlength{\columnsep}{7mm}
\setlength{\columnseprule}{0.35pt}
\columnratio{0.67}
\begin{paracol}{2}

{\sffamily\large\bfseries\color{DDTANavy}Documentazione DDTA}\par
\vspace{0.5em}

\begin{ddtacore}[title={MacroRequirement}]
\textbf{Intent}\par
Determinare, a partire dalle informazioni disponibili sul caso dermatologico, l'urgenza del caso e la relativa priorita' operativa di triage.

\textbf{Context}\par
La valutazione di triage utilizza le informazioni disponibili sul caso dermatologico. La disponibilita' e la tipologia delle evidenze possono variare tra i casi.

\textbf{Stakeholders}\par
Paziente.

\textbf{Scope}\par
\textbf{IN:} determinazione dell'urgenza del caso e della relativa priorita' operativa di triage a partire dalle informazioni di caso governate dal progetto.\par
\textbf{OUT:} indicazione della destinazione specialistica; validazione o correzione medica dell'esito; adattamento successivo del comportamento del sistema sulla base della revisione clinica.

\textbf{Assumptions / Constraints}\par
--

\textbf{dependsOn}\par
\texttt{None}
\end{ddtacore}

\begin{ddtaquestions}
\small
\begin{itemize}
  \item HIGH / MEDIUM / LOW e' una convenzione governata autonomamente o un vocabolario downstream usato coerentemente?
  \item Quando il triage viene eseguito senza immagine, quali risultati deve produrre il percorso \emph{symptom-only}? Deve limitarsi alla determinazione dell'urgenza oppure deve produrre anche altri risultati previsti dal percorso image-based? In caso affermativo, quali?
  \item I percorsi diretto e B4-integrato sono due strategie che il progetto intende preservare oppure due interfacce della realizzazione corrente?
  \item Il recupero di casi storici ha governance autonoma oppure e' comportamento interno della pipeline image-based?
\end{itemize}
\end{ddtaquestions}

\switchcolumn

{\sffamily\large\bfseries\color{DDTABlue}Base Analysis}\par
\vspace{0.5em}

\begin{ddtareferentidentified}
\scriptsize
\BACandidate{Patient} \BANew{}; \BACandidate{CaseInformation} \BANew{}.\par
\end{ddtareferentidentified}

\begin{ddtareferentaccepted}
\scriptsize
\BAAccepted{DermatologicalCase} \BAFirstAccepted{}; \BAAccepted{DermatologicalTriageEvaluation} \BAFirstAccepted{}; \BAAccepted{CaseUrgency} \BAFirstAccepted{}; \BAAccepted{OperationalTriagePriority} \BAFirstAccepted{}.\par
\end{ddtareferentaccepted}

\begin{ddtapropositionidentified}
\scriptsize
{\itshape
\texttt{BAP-MR01-01}\par
\texttt{semanticOperatorKey = produce}\par
\hspace*{1em}\texttt{actor -> DermatologicalTriageEvaluation}\par
\hspace*{1em}\texttt{input -> CaseInformation}\par
\hspace*{1em}\texttt{result -> CaseUrgency}\par
\hspace*{1em}\texttt{result -> OperationalTriagePriority}\par
\texttt{polarity = affirmative}\par
}
\end{ddtapropositionidentified}

\begin{ddtapropositionaccepted}
\scriptsize
--
\end{ddtapropositionaccepted}

\begin{ddtaopenrelations}
\scriptsize
\texttt{CaseUrgency <relazione da determinare> DermatologicalCase}\par
\texttt{OperationalTriagePriority <relazione da determinare> DermatologicalCase}\par
\end{ddtaopenrelations}

\end{paracol}
%</PAGE-CONTENT:3>
\clearpage
%<PAGE-DESC:4>MR-01 > DEC-MR01-01 - Valutazione di urgenza in assenza di immagine
%<PAGE-CONTENT:4>
\PageHashFooter{\PageMDfive{4}}
\subsection{DEC-MR01-01 - Valutazione di urgenza in assenza di immagine}
{\sffamily\scriptsize\color{DDTAGray}Gerarchia: \texttt{MR-01} $\rightarrow$ \texttt{DEC-MR01-01}}\par

\setlength{\columnsep}{7mm}
\setlength{\columnseprule}{0.35pt}
\columnratio{0.67}
\begin{paracol}{2}

{\sffamily\large\bfseries\color{DDTANavy}Documentazione DDTA}\par
\vspace{0.5em}

\begin{ddtacore}[title={Decision}]
\textbf{Parent MR}\par
\texttt{MR-01}

\textbf{Context}\par
Un caso dermatologico puo' avere un'immagine della lesione, oppure no.

\textbf{Decision}\par
DermaTriage supporta la valutazione di urgenza anche per i casi privi di immagine della lesione, adottando in tale situazione una strategia basata sulle informazioni sintomatologiche disponibili sul caso.

\textbf{Consequences}\par
L'immagine non costituisce una precondizione assoluta per la determinazione dell'urgenza. In sua assenza, l'urgenza puo' comunque essere determinata utilizzando le informazioni sintomatologiche disponibili sul caso.
\end{ddtacore}

\begin{ddtaquestions}
\small
\begin{itemize}
  \item Che cosa costituisce esattamente un'immagine assente o non disponibile?
  \item Quali informazioni sintomatologiche minime devono essere disponibili?
  \item Quando l'immagine della lesione non e' disponibile, il progetto governa un comportamento / percorso \emph{symptom-only} distinto oppure un unico comportamento di valutazione nel quale l'immagine costituisce un input opzionale?
\end{itemize}
\end{ddtaquestions}

\switchcolumn

{\sffamily\large\bfseries\color{DDTABlue}Base Analysis}\par
\vspace{0.5em}

\begin{ddtareferentidentified}
\scriptsize
\BACandidate{LesionImage} \BANew{}; \BACandidate{SymptomInformation} \BANew{}.\par
\end{ddtareferentidentified}

\begin{ddtareferentaccepted}
\scriptsize
\BAAccepted{DermaTriage} \BAFirstAccepted{}; \BAAccepted{DermatologicalTriageEvaluation}.\par
\end{ddtareferentaccepted}

\begin{ddtapropositionidentified}
\scriptsize
{\itshape
\texttt{BAP-DEC01-01}\par
\texttt{semanticOperatorKey = produce}\par
\hspace*{1em}\texttt{actor -> DermaTriage}\par
\hspace*{1em}\texttt{input -> SymptomInformation}\par
\hspace*{1em}\texttt{result -> CaseUrgency}\par
\hspace*{1em}\texttt{scopedModifier}\par
\hspace*{2em}\texttt{condition -> LesionImage unavailable}\par
\texttt{polarity = affirmative}\par
}
\end{ddtapropositionidentified}

\begin{ddtapropositionaccepted}
\scriptsize
--
\end{ddtapropositionaccepted}

\begin{ddtaopenrelations}
\scriptsize
\texttt{SymptomInformation <relazione da determinare> CaseInformation}\par
\texttt{LesionImage <relazione da determinare> CaseInformation}\par
\end{ddtaopenrelations}

\end{paracol}
%</PAGE-CONTENT:4>
\clearpage
%<PAGE-DESC:5>MR-01 > DEC-MR01-01 > FR-MR01-01-01 - Determinazione dell'urgenza su base sintomatologica
%<PAGE-CONTENT:5>
\PageHashFooter{\PageMDfive{5}}
\subsubsection{FR-MR01-01-01 - Determinazione dell'urgenza su base sintomatologica in assenza di immagine}
{\sffamily\scriptsize\color{DDTAGray}Gerarchia: \texttt{MR-01} $\rightarrow$ \texttt{DEC-MR01-01} $\rightarrow$ \texttt{FR-MR01-01-01}}\par

\setlength{\columnsep}{7mm}
\setlength{\columnseprule}{0.35pt}
\columnratio{0.67}
\begin{paracol}{2}

{\sffamily\large\bfseries\color{DDTANavy}Documentazione DDTA}\par
\vspace{0.5em}

\begin{ddtacore}[title={FunctionalRequirement}]
\begin{tabularx}{\linewidth}{L{38mm}Y}
\toprule
\textbf{ID} & \texttt{FR-MR01-01-01} \\
\textbf{Lifecycle} & \texttt{candidate} \\
\textbf{Authority} & \texttt{CANDIDATE\_NON\_CURRENT} \\
\textbf{Type} & \texttt{FunctionalRequirement} \\
\textbf{Title} & Determinazione dell'urgenza su base sintomatologica in assenza di immagine \\
\textbf{Parent Decision} & \texttt{DEC-MR01-01} \\
\textbf{Owning MR (derived)} & \texttt{MR-01} \\
\bottomrule
\end{tabularx}

\textbf{Functional obligation}\par
Quando l'immagine della lesione non e' disponibile, DermaTriage determina l'urgenza del caso utilizzando le informazioni sintomatologiche disponibili sullo stesso caso. Nel workflow B4-integrated la documentazione identifica tra tali informazioni campi quali \texttt{itching}, \texttt{bleeding}, \texttt{growth/change} e \texttt{pain}.

\textbf{Functional clauses / inherited \texttt{normativeClause [1..*]}}\par
Quando l'immagine della lesione non e' disponibile, DermaTriage MUST determinare l'urgenza del caso utilizzando le informazioni sintomatologiche disponibili associate allo stesso caso.

\textbf{Riferimenti strutturati / evidenze governate}\par
--

\textbf{SpecializedRequirement relation}\par
Nessuna relazione stabilita.
\end{ddtacore}

\begin{ddtaquestions}
\small
\begin{itemize}
  \item Quali informazioni sintomatologiche sono supportate e quali, se presenti, sono obbligatorie per poter determinare l'urgenza?
  \item Come vengono gestiti campi sintomatologici mancanti, non validi o parziali?
  \item Qual e' la regola completa di symptom scoring che produce l'urgenza?
\end{itemize}
\end{ddtaquestions}

\switchcolumn

{\sffamily\large\bfseries\color{DDTABlue}Base Analysis}\par
\vspace{0.5em}

\begin{ddtareferentidentified}
\scriptsize
--
\end{ddtareferentidentified}

\begin{ddtareferentaccepted}
\scriptsize
\BAAccepted{DermaTriage}; \BAAccepted{DermatologicalCase}; \BAAccepted{DermatologicalTriageEvaluation}; \BAAccepted{CaseUrgency}; \BAAccepted{LesionImage} \BAFirstAccepted{}; \BAAccepted{SymptomInformation} \BAFirstAccepted{}.\par
\end{ddtareferentaccepted}

\begin{ddtapropositionidentified}
\scriptsize
\texttt{BAP-FR01-01-01}\par
\texttt{semanticOperatorKey = produce}\par
\hspace*{1em}\texttt{actor -> DermaTriage}\par
\hspace*{1em}\texttt{input -> SymptomInformation}\par
\hspace*{1em}\texttt{result -> CaseUrgency}\par
\hspace*{1em}\texttt{scopedModifier}\par
\hspace*{2em}\texttt{condition -> LesionImage unavailable}\par
\texttt{polarity = affirmative}
\end{ddtapropositionidentified}

\begin{ddtapropositionaccepted}
\scriptsize
\textbf{BAP-FR01-01-01}.
\end{ddtapropositionaccepted}

\end{paracol}
%</PAGE-CONTENT:5>
\clearpage
%<PAGE-DESC:6>MR-01 > DEC-MR01-03 - Percorso image-based a quattro stadi coordinati
%<PAGE-CONTENT:6>
\PageHashFooter{\PageMDfive{6}}
\subsection{DEC-MR01-03 - Percorso image-based a quattro stadi coordinati}
{\sffamily\scriptsize\color{DDTAGray}Gerarchia: \texttt{MR-01} $\rightarrow$ \texttt{DEC-MR01-03}}\par


\setlength{\columnsep}{7mm}
\setlength{\columnseprule}{0.35pt}
\columnratio{0.67}
\begin{paracol}{2}

{\sffamily\large\bfseries\color{DDTANavy}Documentazione DDTA}\par
\vspace{0.5em}

\begin{ddtacore}[title={Decision}]
\textbf{Parent MR}\par
\texttt{MR-01}

\textbf{Context}\par
Quando e' disponibile un'immagine della lesione, DermaTriage usa un percorso analitico image-based articolato in quattro stadi che producono e combinano risultati diversi prima della sintesi finale di triage.

\textbf{Decision}\par
DermaTriage realizza il percorso image-based mediante quattro stadi analitici coordinati. Stage 1 e Stage 2 elaborano entrambi l'immagine della lesione e producono rispettivamente una classificazione di urgenza con confidence e una descrizione clinica. La descrizione clinica alimenta Stage 3, che recupera casi storici simili. Stage 4 combina gli output degli Stage 1, 2 e 3 con le informazioni sintomatologiche disponibili per produrre la sintesi finale di triage.

\textbf{Consequences}\par
Il percorso image-based richiede la disponibilita' dell'immagine della lesione e resta distinto dal comportamento symptom-only applicabile quando l'immagine non e' disponibile. I quattro stadi mantengono responsabilita' analitiche distinte che convergono nella sintesi finale: Stage 3 usa la descrizione clinica prodotta da Stage 2, mentre Stage 1 e Stage 2 non dipendono reciprocamente dai rispettivi output. La documentazione corrente non stabilisce che tali elaborazioni debbano essere eseguite in parallelo. Le tecnologie e i modelli dei singoli stadi possono cambiare o evolvere senza per forza modificare questa struttura.
\end{ddtacore}

\begin{ddtaquestions}
\small
\begin{itemize}
  \item Il numero esatto di quattro stadi e la loro composizione sono un commitment stabile oppure una struttura della realizzazione corrente?
  \item L'ordine Stage 1--4 governa anche l'ordine temporale di esecuzione oppure soltanto l'organizzazione logica del percorso, fermo il dataflow esplicitamente documentato?
  \item Stage 1 e Stage 2 possono essere eseguiti indipendentemente o in concorrenza quando la stessa immagine e' disponibile, oppure esiste un sequencing runtime non esplicitato?
\end{itemize}
\end{ddtaquestions}


\switchcolumn

{\sffamily\large\bfseries\color{DDTABlue}Base Analysis}\par
\vspace{0.5em}

\begin{ddtareferentidentified}
\scriptsize
\BACandidate{ImageUrgencyAnalysis} \BANew{}; \BACandidate{ClinicalDescriptionGeneration} \BANew{}; \BACandidate{HistoricalCaseRetrieval} \BANew{}; \BACandidate{MultiSourceTriageSynthesis} \BANew{}.\par
\end{ddtareferentidentified}

\begin{ddtareferentaccepted}
\scriptsize
\BAAccepted{DermaTriage}; \BAAccepted{DermatologicalTriageEvaluation}; \BAAccepted{LesionImage}.\par
\end{ddtareferentaccepted}

\begin{ddtapropositionidentified}
\scriptsize
--
\end{ddtapropositionidentified}

\begin{ddtapropositionaccepted}
\scriptsize
--
\end{ddtapropositionaccepted}


\end{paracol}
%</PAGE-CONTENT:6>
\clearpage
%<PAGE-DESC:7>MR-01 > DEC-MR01-03 > FR-MR01-03-01 - Stage 1: urgenza e confidence da immagine
%<PAGE-CONTENT:7>
\PageHashFooter{\PageMDfive{7}}
\subsubsection{FR-MR01-03-01 - Stage 1: urgenza e confidence da immagine}
{\sffamily\scriptsize\color{DDTAGray}Gerarchia: \texttt{MR-01} $\rightarrow$ \texttt{DEC-MR01-03} $\rightarrow$ \texttt{FR-MR01-03-01}}\par


\setlength{\columnsep}{7mm}
\setlength{\columnseprule}{0.35pt}
\columnratio{0.67}
\begin{paracol}{2}

{\sffamily\large\bfseries\color{DDTANavy}Documentazione DDTA}\par
\vspace{0.5em}

\begin{ddtacore}[title={FunctionalRequirement}]
\begin{tabularx}{\linewidth}{L{38mm}Y}
\toprule
\textbf{ID} & \texttt{FR-MR01-03-01} \\
\textbf{Lifecycle} & \texttt{candidate} \\
\textbf{Authority} & \texttt{CANDIDATE\_NON\_CURRENT} \\
\textbf{Type} & \texttt{FunctionalRequirement} \\
\textbf{Title} & Stage 1: urgenza e confidence da immagine \\
\textbf{Parent Decision} & \texttt{DEC-MR01-03} \\
\textbf{Owning MR (derived)} & \texttt{MR-01} \\
\bottomrule
\end{tabularx}

\textbf{Functional obligation}\par
Nel percorso image-based, lo Stage 1 usa un modello \texttt{EfficientNet-B4} fine-tuned per elaborare immagini RGB a \texttt{380x380} e produrre un valore di urgenza \texttt{HIGH / MEDIUM / LOW} con confidence associata. La classificazione \texttt{HIGH} usa una soglia pari a \texttt{0.25}. Il risultato resta distinto dalla successiva priorita' operativa P1--P4.

\textbf{Functional clauses / inherited \texttt{normativeClause [1..*]}}\par
\begin{enumerate}
  \item Lo Stage 1 MUST elaborare l'immagine della lesione mediante il modello \texttt{EfficientNet-B4} fine-tuned su input RGB \texttt{380x380}.
  \item Lo Stage 1 MUST produrre un valore di urgenza appartenente a \texttt{HIGH / MEDIUM / LOW} e una confidence associata.
  \item La classificazione \texttt{HIGH} MUST applicare la soglia \texttt{0.25} nella configurazione documentata.
\end{enumerate}

\textbf{Riferimenti strutturati / evidenze governate}\par
--

\textbf{SpecializedRequirement relation}\par
Nessuna relazione stabilita.
\end{ddtacore}

\begin{ddtaquestions}
\small
\begin{itemize}
  \item La soglia 0.25 e' configurabile a runtime, un default oppure un valore governato e stabile?
  \item Qual e' la regola completa che distingue MEDIUM da LOW?
  \item Come vengono gestite immagini assenti, non valide o non processabili nel percorso image-based?
\end{itemize}
\end{ddtaquestions}




\switchcolumn

{\sffamily\large\bfseries\color{DDTABlue}Base Analysis}\par
\vspace{0.5em}

\begin{ddtareferentidentified}
\scriptsize
\BACandidate{EfficientNet-B4} \BANew{}; \BACandidate{ImageUrgencyClassification} \BANew{}.\par
\end{ddtareferentidentified}

\begin{ddtareferentaccepted}
\scriptsize
\BAAccepted{DermaTriage}; \BAAccepted{LesionImage}; \BAAccepted{ImageUrgencyAnalysis} \BAFirstAccepted{}; \BAAccepted{EfficientNet-B4} \BAFirstAccepted{}; \BAAccepted{ImageUrgencyClassification} \BAFirstAccepted{}.\par
\end{ddtareferentaccepted}

\begin{ddtapropositionidentified}
\tiny
\texttt{BAP-FR03-01-01}\par
\texttt{semanticOperatorKey = produce}\par
\hspace*{1em}\texttt{actor -> ImageUrgencyAnalysis}\par
\hspace*{1em}\texttt{input -> LesionImage}\par
\hspace*{1em}\texttt{result -> ImageUrgencyClassification}\par
\hspace*{1em}\texttt{scopedModifier}\par
\hspace*{2em}\texttt{condition -> LesionImage available}\par
\texttt{polarity = affirmative}\par
\medskip
\texttt{BAP-FR03-01-02}\par
\texttt{semanticOperatorKey = realize}\par
\hspace*{1em}\texttt{abstract -> ImageUrgencyAnalysis}\par
\hspace*{1em}\texttt{realization -> EfficientNet-B4}\par
\texttt{polarity = affirmative}
\end{ddtapropositionidentified}

\begin{ddtapropositionaccepted}
\tiny
\textbf{BAP-FR03-01-01}; \textbf{BAP-FR03-01-02}.
\end{ddtapropositionaccepted}


\end{paracol}
%</PAGE-CONTENT:7>
\clearpage
%<PAGE-DESC:8>MR-01 > DEC-MR01-03 > FR-MR01-03-02 - Stage 2: descrizione clinica strutturata
%<PAGE-CONTENT:8>
\PageHashFooter{\PageMDfive{8}}
\subsubsection{FR-MR01-03-02 - Stage 2: descrizione clinica strutturata}
{\sffamily\scriptsize\color{DDTAGray}Gerarchia: \texttt{MR-01} $\rightarrow$ \texttt{DEC-MR01-03} $\rightarrow$ \texttt{FR-MR01-03-02}}\par


\setlength{\columnsep}{7mm}
\setlength{\columnseprule}{0.35pt}
\columnratio{0.67}
\begin{paracol}{2}

{\sffamily\large\bfseries\color{DDTANavy}Documentazione DDTA}\par
\vspace{0.5em}

\begin{ddtacore}[title={FunctionalRequirement}]
\begin{tabularx}{\linewidth}{L{38mm}Y}
\toprule
\textbf{ID} & \texttt{FR-MR01-03-02} \\
\textbf{Lifecycle} & \texttt{candidate} \\
\textbf{Authority} & \texttt{CANDIDATE\_NON\_CURRENT} \\
\textbf{Type} & \texttt{FunctionalRequirement} \\
\textbf{Title} & Stage 2: descrizione clinica strutturata \\
\textbf{Parent Decision} & \texttt{DEC-MR01-03} \\
\textbf{Owning MR (derived)} & \texttt{MR-01} \\
\bottomrule
\end{tabularx}

\textbf{Functional obligation}\par
Nel percorso image-based, lo Stage 2 usa \texttt{Qwen2-VL-7B-Instruct} come vision-language model per derivare dall'immagine della lesione una descrizione clinica testuale strutturata in cinque parti: \texttt{shape}, \texttt{borders}, \texttt{color}, \texttt{texture} e \texttt{suspicious features}. La descrizione prodotta e' resa disponibile allo Stage 3 per il retrieval dei casi storici e resta disponibile allo Stage 4 come output dello Stage 2.

\textbf{Functional clauses / inherited \texttt{normativeClause [1..*]}}\par
\begin{enumerate}
  \item Lo Stage 2 MUST usare \texttt{Qwen2-VL-7B-Instruct} per elaborare l'immagine della lesione.
  \item Lo Stage 2 MUST produrre una descrizione clinica testuale strutturata in cinque parti: \texttt{shape}, \texttt{borders}, \texttt{color}, \texttt{texture} e \texttt{suspicious features}.
  \item Lo Stage 2 MUST rendere la descrizione prodotta disponibile allo Stage 3 per il retrieval e allo Stage 4 come proprio output analitico.
\end{enumerate}

\textbf{Riferimenti strutturati / evidenze governate}\par
--

\textbf{SpecializedRequirement relation}\par
Nessuna relazione stabilita.
\end{ddtacore}

\begin{ddtaquestions}
\small
\begin{itemize}
  \item Ordine, lingua, label e separatori delle cinque parti sono obbligatori o modificabili?
  \item Qual e' il comportamento richiesto se la descrizione non puo' essere prodotta completamente?
  \item Limiti di lunghezza, encoding/charset e regole per caratteri speciali sono governati oppure restano non specificati?
\end{itemize}
\end{ddtaquestions}




\switchcolumn

{\sffamily\large\bfseries\color{DDTABlue}Base Analysis}\par
\vspace{0.5em}

\begin{ddtareferentidentified}
\scriptsize
\BACandidate{Qwen2-VL-7B-Instruct} \BANew{}; \BACandidate{ClinicalDescription} \BANew{}; \BACandidate{ClinicalDescriptionContentContract} \BANew{}.\par
\end{ddtareferentidentified}

\begin{ddtareferentaccepted}
\scriptsize
\BAAccepted{DermaTriage}; \BAAccepted{LesionImage}; \BAAccepted{ClinicalDescriptionGeneration} \BAFirstAccepted{}; \BAAccepted{Qwen2-VL-7B-Instruct} \BAFirstAccepted{}; \BAAccepted{ClinicalDescription} \BAFirstAccepted{}; \BAAccepted{ClinicalDescriptionContentContract} \BAFirstAccepted{}.\par
\end{ddtareferentaccepted}

\begin{ddtapropositionidentified}
\fontsize{5.3}{5.8}\selectfont
\texttt{BAP-FR03-02-01}\par
\texttt{semanticOperatorKey = produce}\par
\hspace*{1em}\texttt{actor -> ClinicalDescriptionGeneration}\par
\hspace*{1em}\texttt{input -> LesionImage}\par
\hspace*{1em}\texttt{result -> ClinicalDescription}\par
\hspace*{1em}\texttt{scopedModifier}\par
\hspace*{2em}\texttt{condition -> LesionImage available}\par
\texttt{polarity = affirmative}\par
\smallskip
\texttt{BAP-FR03-02-02}\par
\texttt{semanticOperatorKey = realize}\par
\hspace*{1em}\texttt{abstract -> ClinicalDescriptionGeneration}\par
\hspace*{1em}\texttt{realization -> Qwen2-VL-7B-Instruct}\par
\texttt{polarity = affirmative}\par
\smallskip
\texttt{BAP-FR03-02-03}\par
\texttt{semanticOperatorKey = constrain}\par
\hspace*{1em}\texttt{constraintTarget -> ClinicalDescription}\par
\hspace*{1em}\texttt{constraintValue -> ClinicalDescriptionContentContract}\par
\texttt{polarity = affirmative}\par
\smallskip
\texttt{BAP-FR03-02-04}\par
\texttt{semanticOperatorKey = constrain}\par
\hspace*{1em}\texttt{constraintTarget -> ClinicalDescriptionContentContract}\par
\hspace*{1em}\texttt{constraintValue: property -> representation; presence -> required; vocabulary -> [TEXTUAL]}\par
\hspace*{1em}\texttt{constraintValue: property -> partCount; presence -> required; vocabulary -> [5]}\par
\hspace*{1em}\texttt{constraintValue: property -> shape; presence -> required}\par
\hspace*{1em}\texttt{constraintValue: property -> borders; presence -> required}\par
\hspace*{1em}\texttt{constraintValue: property -> color; presence -> required}\par
\hspace*{1em}\texttt{constraintValue: property -> texture; presence -> required}\par
\hspace*{1em}\texttt{constraintValue: property -> suspiciousFeatures; presence -> required}\par
\texttt{polarity = affirmative}
\end{ddtapropositionidentified}

\begin{ddtapropositionaccepted}
\fontsize{5.5}{6.0}\selectfont
\textbf{BAP-FR03-02-01}; \textbf{BAP-FR03-02-02}; \textbf{BAP-FR03-02-03}; \textbf{BAP-FR03-02-04}.
\end{ddtapropositionaccepted}

\end{paracol}
%</PAGE-CONTENT:8>
\clearpage
%<PAGE-DESC:9>MR-01 > DEC-MR01-03 > FR-MR01-03-03A [review label] - Preparazione della base storica per il retrieval
%<PAGE-CONTENT:9>
\PageHashFooter{\PageMDfive{9}}
\subsubsection{FR-MR01-03-03A [review label] - Preparazione della base storica per il retrieval}
{\sffamily\scriptsize\color{DDTAGray}Gerarchia candidata: \texttt{MR-01} $\rightarrow$ \texttt{DEC-MR01-03} $\rightarrow$ \texttt{FR-MR01-03-03A}; label provvisorio per la review dello split.}\par

\setlength{\columnsep}{7mm}
\setlength{\columnseprule}{0.35pt}
\columnratio{0.67}
\begin{paracol}{2}

{\sffamily\large\bfseries\color{DDTANavy}Documentazione DDTA}\par
\vspace{0.4em}

\begin{ddtacore}[title={FunctionalRequirement - candidato di split}]
\begin{tabularx}{\linewidth}{L{38mm}Y}
\toprule
\textbf{ID} & \texttt{FR-MR01-03-03A} (review label) \\
\textbf{Lifecycle} & \texttt{candidate} \\
\textbf{Authority} & \texttt{CANDIDATE\_NON\_CURRENT} \\
\textbf{Type} & \texttt{FunctionalRequirement} \\
\textbf{Title} & Preparazione della base storica per il retrieval \\
\textbf{Parent Decision} & \texttt{DEC-MR01-03} (candidate parent da confermare) \\
\textbf{Owning MR (derived)} & \texttt{MR-01} \\
\bottomrule
\end{tabularx}

\textbf{Functional obligation}\par
Per rendere disponibile allo Stage 3 la base storica usata dal retrieval, le descrizioni dei casi dermatologici storici contenute nel campo \texttt{image\_description} vengono trasformate in rappresentazioni vettoriali e indicizzate in un archivio interrogabile. Nella realizzazione corrente, gli embedding sono prodotti con \texttt{sentence-transformers/all-MiniLM-L6-v2} e indicizzati in \texttt{ChromaDB}, che svolge il ruolo di archivio vettoriale.

\textbf{Functional clauses / inherited \texttt{normativeClause [1..*]}}\par
\begin{enumerate}
  \item La base storica utilizzata dal retrieval MUST essere preparata a partire dalle descrizioni dei casi dermatologici storici usate per il confronto di similarita'.
  \item Le rappresentazioni preparate MUST essere indicizzate in una forma interrogabile dallo Stage 3.
\end{enumerate}

\textbf{Riferimenti strutturati / evidenze governate}\par
--

\textbf{SpecializedRequirement relation}\par
Nessuna relazione stabilita.
\end{ddtacore}

\begin{ddtaquestions}
\fontsize{6.0}{6.8}\selectfont
\begin{itemize}
  \item Quando deve essere costruita, aggiornata o rigenerata la base usata dal retrieval?
  \item Quale comportamento e' richiesto se la base non e' disponibile, e' incompleta o non e' aggiornata?
  \item \texttt{DEC-MR01-03} e' il semantic owner corretto oppure la preparazione della base esprime una scelta progettuale autonoma?
\end{itemize}
\end{ddtaquestions}

\switchcolumn

{\sffamily\large\bfseries\color{DDTABlue}Base Analysis}\par
\vspace{0.4em}

\begin{ddtareferentidentified}[top=3pt,bottom=3pt,left=5pt,right=5pt]
\fontsize{5.35}{5.9}\selectfont
\BACandidate{HistoricalCaseDescription} \BANew{}; \BACandidate{HistoricalCaseEmbeddingGeneration} \BANew{}; \BACandidate{HistoricalCaseEmbedding} \BANew{}; \BACandidate{HistoricalCaseVectorIndexing} \BANew{}; \BACandidate{HistoricalCaseVectorIndex} \BANew{}.
\end{ddtareferentidentified}

\begin{ddtareferentaccepted}[top=3pt,bottom=3pt,left=5pt,right=5pt]
\fontsize{5.15}{5.7}\selectfont
\BAAccepted{HistoricalCaseDescription} \BAFirstAccepted{}; \BAAccepted{HistoricalCaseEmbeddingGeneration} \BAFirstAccepted{}; \BAAccepted{HistoricalCaseEmbedding} \BAFirstAccepted{}; \BAAccepted{HistoricalCaseVectorIndexing} \BAFirstAccepted{}; \BAAccepted{HistoricalCaseVectorIndex} \BAFirstAccepted{}; \BAAccepted{all-MiniLM-L6-v2} \BAFirstAccepted{}; \BAAccepted{ChromaDB} \BAFirstAccepted{}.
\end{ddtareferentaccepted}

\begin{ddtapropositionidentified}[top=3pt,bottom=3pt,left=5pt,right=5pt]
\fontsize{4.75}{5.15}\selectfont
\texttt{BAP-FR03-03-06}\par
\texttt{semanticOperatorKey = produce}\par
\hspace*{1em}\texttt{actor -> HistoricalCaseEmbeddingGeneration}\par
\hspace*{1em}\texttt{input -> HistoricalCaseDescription}\par
\hspace*{1em}\texttt{result -> HistoricalCaseEmbedding}\par
\texttt{polarity = affirmative}\par
\smallskip
\texttt{BAP-FR03-03-07}\par
\texttt{semanticOperatorKey = realize}\par
\hspace*{1em}\texttt{abstract -> HistoricalCaseEmbeddingGeneration}\par
\hspace*{1em}\texttt{realization -> all-MiniLM-L6-v2}\par
\texttt{polarity = affirmative}\par
\smallskip
\texttt{BAP-FR03-03-08}\par
\texttt{semanticOperatorKey = produce}\par
\hspace*{1em}\texttt{actor -> HistoricalCaseVectorIndexing}\par
\hspace*{1em}\texttt{input -> HistoricalCaseEmbedding}\par
\hspace*{1em}\texttt{result -> HistoricalCaseVectorIndex}\par
\texttt{polarity = affirmative}\par
\smallskip
\texttt{BAP-FR03-03-09}\par
\texttt{semanticOperatorKey = realize}\par
\hspace*{1em}\texttt{abstract -> HistoricalCaseVectorIndexing}\par
\hspace*{1em}\texttt{realization -> ChromaDB}\par
\texttt{polarity = affirmative}
\end{ddtapropositionidentified}

\begin{ddtapropositionaccepted}[top=3pt,bottom=3pt,left=5pt,right=5pt]
\fontsize{5.0}{5.5}\selectfont
\textbf{BAP-FR03-03-06}; \textbf{BAP-FR03-03-07}; \textbf{BAP-FR03-03-08}; \textbf{BAP-FR03-03-09}.
\end{ddtapropositionaccepted}

\end{paracol}
%</PAGE-CONTENT:9>
\clearpage
%<PAGE-DESC:10>MR-01 > DEC-MR01-03 > FR-MR01-03-03B [review label] - Stage 3: recupero di casi storici simili
%<PAGE-CONTENT:10>
\PageHashFooter{\PageMDfive{10}}
\subsubsection{FR-MR01-03-03B [review label] - Stage 3: recupero di casi storici simili}
{\sffamily\scriptsize\color{DDTAGray}Gerarchia candidata: \texttt{MR-01} $\rightarrow$ \texttt{DEC-MR01-03} $\rightarrow$ \texttt{FR-MR01-03-03B}; label provvisorio per la review dello split.}\par

\setlength{\columnsep}{7mm}
\setlength{\columnseprule}{0.35pt}
\columnratio{0.67}
\begin{paracol}{2}

{\sffamily\large\bfseries\color{DDTANavy}Documentazione DDTA}\par
\vspace{0.4em}

\begin{ddtacore}[title={FunctionalRequirement - candidato di split}]
\begin{tabularx}{\linewidth}{L{38mm}Y}
\toprule
\textbf{ID} & \texttt{FR-MR01-03-03B} (review label) \\
\textbf{Lifecycle} & \texttt{candidate} \\
\textbf{Authority} & \texttt{CANDIDATE\_NON\_CURRENT} \\
\textbf{Type} & \texttt{FunctionalRequirement} \\
\textbf{Title} & Stage 3: recupero di casi storici simili \\
\textbf{Parent Decision} & \texttt{DEC-MR01-03} \\
\textbf{Owning MR (derived)} & \texttt{MR-01} \\
\bottomrule
\end{tabularx}

\textbf{Functional obligation}\par
Nel percorso image-based, lo Stage 3 recupera casi dermatologici storici simili al caso corrente usando come query la descrizione clinica prodotta dallo Stage 2 e interrogando la base storica predisposta per il retrieval. I casi recuperati vengono resi disponibili allo Stage 4 come input della sintesi multi-source.

Nella realizzazione corrente, la descrizione clinica del caso corrente viene rappresentata mediante \texttt{sentence-transformers/all-MiniLM-L6-v2}; \texttt{ChromaDB} confronta tale rappresentazione con l'archivio storico mediante cosine similarity e recupera i cinque casi con maggiore similarita'. Per ciascun caso recuperato il risultato espone \texttt{patient\_id}, \texttt{diagnosis}, \texttt{urgency\_level} e \texttt{similarity\_score}.

\textbf{Functional clauses / inherited \texttt{normativeClause [1..*]}}\par
\begin{enumerate}
  \item Lo Stage 3 MUST usare la descrizione clinica prodotta dallo Stage 2 come input del retrieval.
  \item Lo Stage 3 MUST recuperare casi dermatologici storici simili interrogando la base storica predisposta per il retrieval.
  \item Lo Stage 3 MUST rendere i casi recuperati disponibili allo Stage 4 come input della sintesi multi-source.
\end{enumerate}

\textbf{Riferimenti strutturati / evidenze governate}\par
--

\textbf{SpecializedRequirement relation}\par
Nessuna relazione stabilita.
\end{ddtacore}

\begin{ddtaquestions}
\fontsize{6.0}{6.8}\selectfont
\begin{itemize}
  \item Il valore top-5 e' un requisito stabile oppure una configurazione modificabile?
  \item La cosine similarity e' una proprieta' governata del retrieval oppure appartiene alla realizzazione corrente?
  \item Che cosa deve accadere quando sono disponibili meno di cinque casi utilizzabili?
  \item I campi \texttt{patient\_id}, \texttt{diagnosis}, \texttt{urgency\_level} e \texttt{similarity\_score} costituiscono il contratto completo dell'output verso Stage 4?
\end{itemize}
\end{ddtaquestions}

\switchcolumn

{\sffamily\large\bfseries\color{DDTABlue}Base Analysis}\par
\vspace{0.4em}

\begin{ddtareferentidentified}[top=3pt,bottom=3pt,left=5pt,right=5pt]
\fontsize{5.35}{5.9}\selectfont
\BACandidate{ClinicalDescriptionEmbeddingGeneration} \BANew{}; \BACandidate{ClinicalDescriptionEmbedding} \BANew{}; \BACandidate{HistoricalCaseSimilarityRetrieval} \BANew{}; \BACandidate{RetrievedHistoricalCases} \BANew{}.
\end{ddtareferentidentified}

\begin{ddtareferentaccepted}[top=3pt,bottom=3pt,left=5pt,right=5pt]
\fontsize{4.95}{5.45}\selectfont
\BAAccepted{ClinicalDescriptionGeneration}; \BAAccepted{ClinicalDescription}; \BAAccepted{HistoricalCaseRetrieval} \BAFirstAccepted{}; \BAAccepted{ClinicalDescriptionEmbeddingGeneration} \BAFirstAccepted{}; \BAAccepted{ClinicalDescriptionEmbedding} \BAFirstAccepted{}; \BAAccepted{HistoricalCaseVectorIndex}; \BAAccepted{HistoricalCaseSimilarityRetrieval} \BAFirstAccepted{}; \BAAccepted{RetrievedHistoricalCases} \BAFirstAccepted{}; \BAAccepted{all-MiniLM-L6-v2}; \BAAccepted{ChromaDB}; \BAAccepted{MultiSourceTriageSynthesis} \BAFirstAccepted{}.
\end{ddtareferentaccepted}

\begin{ddtapropositionidentified}[top=3pt,bottom=3pt,left=5pt,right=5pt]
\fontsize{4.25}{4.65}\selectfont
\texttt{BAP-FR03-03-04}\par
\texttt{semanticOperatorKey = produce}\par
\hspace*{1em}\texttt{actor -> ClinicalDescriptionEmbeddingGeneration}\par
\hspace*{1em}\texttt{input -> ClinicalDescription}\par
\hspace*{1em}\texttt{result -> ClinicalDescriptionEmbedding}\par
\hspace*{1em}\texttt{scopedModifier}\par
\hspace*{2em}\texttt{condition -> LesionImage available}\par
\texttt{polarity = affirmative}\par
\smallskip
\texttt{BAP-FR03-03-05}\par
\texttt{semanticOperatorKey = realize}\par
\hspace*{1em}\texttt{abstract -> ClinicalDescriptionEmbeddingGeneration}\par
\hspace*{1em}\texttt{realization -> all-MiniLM-L6-v2}\par
\texttt{polarity = affirmative}\par
\smallskip
\texttt{BAP-FR03-03-10}\par
\texttt{semanticOperatorKey = produce}\par
\hspace*{1em}\texttt{actor -> HistoricalCaseSimilarityRetrieval}\par
\hspace*{1em}\texttt{input -> ClinicalDescriptionEmbedding}\par
\hspace*{1em}\texttt{input -> HistoricalCaseVectorIndex}\par
\hspace*{1em}\texttt{result -> RetrievedHistoricalCases}\par
\hspace*{1em}\texttt{scopedModifier}\par
\hspace*{2em}\texttt{condition -> LesionImage available}\par
\texttt{polarity = affirmative}\par
\smallskip
\texttt{BAP-FR03-03-11}\par
\texttt{semanticOperatorKey = realize}\par
\hspace*{1em}\texttt{abstract -> HistoricalCaseSimilarityRetrieval}\par
\hspace*{1em}\texttt{realization -> ChromaDB}\par
\texttt{polarity = affirmative}\par
\smallskip
\texttt{BAP-FR03-03-12}\par
\texttt{semanticOperatorKey = transfer}\par
\hspace*{1em}\texttt{source -> ClinicalDescriptionGeneration}\par
\hspace*{1em}\texttt{destination -> ClinicalDescriptionEmbeddingGeneration}\par
\hspace*{1em}\texttt{content -> ClinicalDescription}\par
\hspace*{1em}\texttt{scopedModifier}\par
\hspace*{2em}\texttt{condition -> LesionImage available}\par
\texttt{polarity = affirmative}\par
\smallskip
\texttt{BAP-FR03-03-13}\par
\texttt{semanticOperatorKey = transfer}\par
\hspace*{1em}\texttt{source -> HistoricalCaseSimilarityRetrieval}\par
\hspace*{1em}\texttt{destination -> MultiSourceTriageSynthesis}\par
\hspace*{1em}\texttt{content -> RetrievedHistoricalCases}\par
\hspace*{1em}\texttt{scopedModifier}\par
\hspace*{2em}\texttt{condition -> LesionImage available}\par
\texttt{polarity = affirmative}
\end{ddtapropositionidentified}

\begin{ddtapropositionaccepted}[top=3pt,bottom=3pt,left=5pt,right=5pt]
\fontsize{4.75}{5.2}\selectfont
\textbf{BAP-FR03-03-04}; \textbf{BAP-FR03-03-05}; \textbf{BAP-FR03-03-10}; \textbf{BAP-FR03-03-11}; \textbf{BAP-FR03-03-12}; \textbf{BAP-FR03-03-13}.
\end{ddtapropositionaccepted}

\end{paracol}
%</PAGE-CONTENT:10>
\clearpage
%<PAGE-DESC:11>MR-01 > DEC-MR01-03 > FR-MR01-03-04 - Stage 4: sintesi multi-source
%<PAGE-CONTENT:11>
\PageHashFooter{\PageMDfive{11}}
\subsubsection{FR-MR01-03-04 - Stage 4: sintesi multi-source}
{\sffamily\scriptsize\color{DDTAGray}Gerarchia: \texttt{MR-01} $\rightarrow$ \texttt{DEC-MR01-03} $\rightarrow$ \texttt{FR-MR01-03-04}}\par


\setlength{\columnsep}{7mm}
\setlength{\columnseprule}{0.35pt}
\columnratio{0.67}
\begin{paracol}{2}

{\sffamily\large\bfseries\color{DDTANavy}Documentazione DDTA}\par
\vspace{0.5em}

\begin{ddtacore}[title={FunctionalRequirement}]
\begin{tabularx}{\linewidth}{L{38mm}Y}
\toprule
\textbf{ID} & \texttt{FR-MR01-03-04} \\
\textbf{Lifecycle} & \texttt{candidate} \\
\textbf{Authority} & \texttt{CANDIDATE\_NON\_CURRENT} \\
\textbf{Type} & \texttt{FunctionalRequirement} \\
\textbf{Title} & Stage 4: sintesi multi-source \\
\textbf{Parent Decision} & \texttt{DEC-MR01-03} \\
\textbf{Owning MR (derived)} & \texttt{MR-01} \\
\bottomrule
\end{tabularx}

\textbf{Functional obligation}\par
Lo Stage 4 usa \texttt{BioMistral-7B} per combinare gli output degli Stage 1, 2 e 3 con le informazioni sintomatologiche disponibili. L'output e' una sintesi AI di triage in JSON che contiene almeno \texttt{urgency}, \texttt{confidence}, \texttt{reasoning} e \texttt{pathology}. I valori \texttt{urgency} e \texttt{confidence} sono resi disponibili al successivo passaggio di derivazione della priorita' operativa, realizzato nell'architettura corrente dall'\emph{Adaptation Layer}, per determinare la priorita' P1--P4; la documentazione corrente non stabilisce che \texttt{reasoning} e \texttt{pathology} partecipino a tale mapping.

\textbf{Functional clauses / inherited \texttt{normativeClause [1..*]}}\par
\begin{enumerate}
  \item Lo Stage 4 MUST usare \texttt{BioMistral-7B} per eseguire la sintesi multi-source.
  \item La sintesi MUST combinare gli output disponibili degli Stage 1, 2 e 3 con le informazioni sintomatologiche del caso.
  \item Lo Stage 4 MUST produrre una sintesi AI di triage in rappresentazione JSON contenente almeno \texttt{urgency}, \texttt{confidence}, \texttt{reasoning} e \texttt{pathology}; i valori \texttt{urgency} e \texttt{confidence} MUST essere resi disponibili al successivo passaggio di derivazione della priorita' operativa.
\end{enumerate}

\textbf{Riferimenti strutturati / evidenze governate}\par
\texttt{DEC-MR01-04}; \texttt{FR-MR01-04-01}

\textbf{SpecializedRequirement relation}\par
Nessuna relazione stabilita.
\end{ddtacore}

\begin{ddtaquestions}
\scriptsize
\begin{itemize}
  \setlength{\itemsep}{1pt}
  \item Qual e' il contratto JSON governato della \texttt{AITriageSynthesis}: quali sono i nomi canonici dei campi relativi a urgency e pathology e deve essere presente anche \texttt{recommended\_action}?
  \item Qual e' il significato governato del contenuto relativo alla pathology e quali usi sono autorizzati?
  \item Come deve comportarsi la sintesi se uno degli input degli stadi precedenti e' assente o incompleto?
\end{itemize}
\end{ddtaquestions}



\switchcolumn

{\sffamily\large\bfseries\color{DDTABlue}Base Analysis}\par
\vspace{0.5em}

\begin{ddtareferentidentified}[top=3pt,bottom=3pt,left=5pt,right=5pt]
\fontsize{5.6}{6.1}\selectfont
\BACandidate{BioMistral-7B} \BANew{}; \BACandidate{TriageConfidence} \BANew{}; \BACandidate{AITriageSynthesis} \BANew{}; \BACandidate{AITriageSynthesisContentContract} \BANew{}.
\end{ddtareferentidentified}

\begin{ddtareferentaccepted}[top=3pt,bottom=3pt,left=5pt,right=5pt]
\fontsize{5.3}{5.8}\selectfont
\BAAccepted{DermaTriage}; \BAAccepted{DermatologicalCase}; \BAAccepted{ImageUrgencyClassification}; \BAAccepted{ClinicalDescription}; \BAAccepted{RetrievedHistoricalCases}; \BAAccepted{SymptomInformation}; \BAAccepted{CaseUrgency}; \BAAccepted{MultiSourceTriageSynthesis}; \BAAccepted{BioMistral-7B} \BAFirstAccepted{}; \BAAccepted{TriageConfidence} \BAFirstAccepted{}; \BAAccepted{AITriageSynthesis} \BAFirstAccepted{}; \BAAccepted{AITriageSynthesisContentContract} \BAFirstAccepted{}.
\end{ddtareferentaccepted}

\begin{ddtapropositionidentified}[top=3pt,bottom=3pt,left=5pt,right=5pt]
\fontsize{4.0}{4.35}\selectfont
\texttt{BAP-FR03-04-01}\par
\texttt{semanticOperatorKey = produce}\par
\hspace*{1em}\texttt{actor -> MultiSourceTriageSynthesis}\par
\hspace*{1em}\texttt{input -> ImageUrgencyClassification}\par
\hspace*{1em}\texttt{input -> ClinicalDescription}\par
\hspace*{1em}\texttt{input -> RetrievedHistoricalCases}\par
\hspace*{1em}\texttt{input -> SymptomInformation}\par
\hspace*{1em}\texttt{result -> AITriageSynthesis}\par
\hspace*{1em}\texttt{result -> CaseUrgency}\par
\hspace*{1em}\texttt{result -> TriageConfidence}\par
\hspace*{1em}\texttt{scopedModifier}\par
\hspace*{2em}\texttt{condition -> LesionImage available}\par
\texttt{polarity = affirmative}\par
\smallskip
\texttt{BAP-FR03-04-02}\par
\texttt{semanticOperatorKey = realize}\par
\hspace*{1em}\texttt{abstract -> MultiSourceTriageSynthesis}\par
\hspace*{1em}\texttt{realization -> BioMistral-7B}\par
\texttt{polarity = affirmative}\par
\smallskip
\texttt{BAP-FR03-04-03}\par
\texttt{semanticOperatorKey = constrain}\par
\hspace*{1em}\texttt{constraintTarget -> AITriageSynthesis}\par
\hspace*{1em}\texttt{constraintValue -> AITriageSynthesisContentContract}\par
\texttt{polarity = affirmative}\par
\smallskip
\texttt{BAP-FR03-04-04}\par
\texttt{semanticOperatorKey = constrain}\par
\hspace*{1em}\texttt{constraintTarget -> AITriageSynthesisContentContract}\par
\hspace*{1em}\texttt{constraintValue -> representation required in [JSON]}\par
\hspace*{1em}\texttt{constraintValue -> urgency required}\par
\hspace*{1em}\texttt{constraintValue -> confidence required}\par
\hspace*{1em}\texttt{constraintValue -> reasoning required}\par
\hspace*{1em}\texttt{constraintValue -> pathology required}\par
\texttt{polarity = affirmative}
\end{ddtapropositionidentified}

\begin{ddtapropositionaccepted}[top=3pt,bottom=3pt,left=5pt,right=5pt]
\fontsize{5.2}{5.7}\selectfont
\textbf{BAP-FR03-04-01}; \textbf{BAP-FR03-04-02}; \textbf{BAP-FR03-04-03}; \textbf{BAP-FR03-04-04}.
\end{ddtapropositionaccepted}


\end{paracol}
%</PAGE-CONTENT:11>
\clearpage
%<PAGE-DESC:12>MR-01 > DEC-MR01-04 - Separazione tra valutazione dell'urgenza e priorita operativa
%<PAGE-CONTENT:12>
\PageHashFooter{\PageMDfive{12}}
\subsection{DEC-MR01-04 - Separazione tra valutazione dell'urgenza e priorita' operativa}
{\sffamily\scriptsize\color{DDTAGray}Gerarchia: \texttt{MR-01} $\rightarrow$ \texttt{DEC-MR01-04}}\par


\setlength{\columnsep}{7mm}
\setlength{\columnseprule}{0.35pt}
\columnratio{0.67}
\begin{paracol}{2}

{\sffamily\large\bfseries\color{DDTANavy}Documentazione DDTA}\par
\vspace{0.5em}

\begin{ddtacore}[title={Decision}]
\textbf{Parent MR}\par
\texttt{MR-01}

\textbf{Context}\par
Nel percorso basato sull'immagine, la pipeline analitica produce una \texttt{AITriageSynthesis} che rende disponibili, tra gli altri contenuti governati, la valutazione analitica di urgenza e la confidence associata. La priorita' operativa per la presa in carico viene determinata in un passaggio successivo e richiede un dominio di valori governato.

\textbf{Decision}\par
Nel percorso basato sull'immagine, DermaTriage mantiene separata la sintesi analitica di triage dalla successiva determinazione della priorita' operativa. La derivazione operativa utilizza la valutazione di urgenza e la confidence applicabile rese disponibili dalla \texttt{AITriageSynthesis}; la priorita' risultante e' rappresentata mediante la scala P1-P4.

\textbf{Consequences}\par
La \texttt{AITriageSynthesis} resta un risultato della pipeline analitica e la priorita' operativa viene determinata successivamente, assumendo uno dei valori P1, P2, P3 o P4. La regola che seleziona il livello e' governata dal relativo FunctionalRequirement. Questa Decision non stabilisce se il passaggio successivo riceva l'intera \texttt{AITriageSynthesis} oppure soltanto i valori di urgenza e confidence necessari al mapping.
\end{ddtacore}

\begin{ddtaquestions}
\small
\begin{itemize}
  \item La separazione tra \texttt{AITriageSynthesis} e derivazione della priorita' e' un commitment architetturale stabile oppure una struttura della realizzazione corrente?
  \item Il passaggio verso la derivazione operativa trasferisce l'intera \texttt{AITriageSynthesis} oppure soltanto \texttt{CaseUrgency} e \texttt{TriageConfidence}?
  \item La stessa derivazione P1-P4 si applica anche alla \texttt{CaseUrgency} prodotta dal percorso symptom-only?
\end{itemize}
\end{ddtaquestions}


\switchcolumn

{\sffamily\large\bfseries\color{DDTABlue}Base Analysis}\par
\vspace{0.5em}

\begin{ddtareferentidentified}
\scriptsize
--
\end{ddtareferentidentified}

\begin{ddtareferentaccepted}
\scriptsize
\BAAccepted{DermaTriage}; \BAAccepted{AITriageSynthesis}; \BAAccepted{CaseUrgency}; \BAAccepted{TriageConfidence}; \BAAccepted{OperationalTriagePriority}.\par
\end{ddtareferentaccepted}

\begin{ddtapropositionidentified}
\scriptsize
\emph{\texttt{BAP-DEC04-01}}\par
\emph{\texttt{semanticOperatorKey = constrain}}\par
\emph{\hspace*{1em}\texttt{constraintTarget -> OperationalTriagePriority}}\par
\emph{\hspace*{1em}\texttt{constraintValue -> \{P1, P2, P3, P4\}}}\par
\emph{\texttt{polarity = affirmative}}
\end{ddtapropositionidentified}

\begin{ddtapropositionaccepted}
\scriptsize
--
\end{ddtapropositionaccepted}

\end{paracol}
%</PAGE-CONTENT:12>
\clearpage
%<PAGE-DESC:13>MR-01 > DEC-MR01-04 > FR-MR01-04-01 - Mapping urgenza/confidence verso P1-P4
%<PAGE-CONTENT:13>
\PageHashFooter{\PageMDfive{13}}
\subsubsection{FR-MR01-04-01 - Mapping urgenza/confidence verso P1-P4}
{\sffamily\scriptsize\color{DDTAGray}Gerarchia: \texttt{MR-01} $\rightarrow$ \texttt{DEC-MR01-04} $\rightarrow$ \texttt{FR-MR01-04-01}}\par

\setlength{\columnsep}{7mm}
\setlength{\columnseprule}{0.35pt}
\columnratio{0.67}
\begin{paracol}{2}

{\sffamily\large\bfseries\color{DDTANavy}Documentazione DDTA}\par
\vspace{0.3em}
\small

\begin{ddtacore}[title={FunctionalRequirement}]
\begin{tabularx}{\linewidth}{L{38mm}Y}
\toprule
\textbf{ID} & \texttt{FR-MR01-04-01} \\
\textbf{Lifecycle} & \texttt{candidate} \\
\textbf{Authority} & \texttt{CANDIDATE\_NON\_CURRENT} \\
\textbf{Type} & \texttt{FunctionalRequirement} \\
\textbf{Title} & Mapping urgenza/confidence verso P1--P4 \\
\textbf{Parent Decision} & \texttt{DEC-MR01-04} \\
\textbf{Owning MR (derived)} & \texttt{MR-01} \\
\bottomrule
\end{tabularx}

\textbf{Functional obligation}\par
Nel percorso basato sull'immagine, DermaTriage deriva la priorita' operativa P1--P4 dalla valutazione analitica di urgenza e dalla confidence applicabile rese disponibili dalla sintesi finale di triage. Per i casi \texttt{HIGH}, la soglia stretta \texttt{confidence > 0.85} seleziona \texttt{P1}; gli altri casi \texttt{HIGH} selezionano \texttt{P2}; \texttt{MEDIUM} seleziona \texttt{P3}; \texttt{LOW} seleziona \texttt{P4}. Nella realizzazione corrente, tale derivazione e' svolta dall'\emph{Adaptation Layer}, implementato nel modulo \texttt{adaptation\_layer.py}; la funzione \texttt{map\_urgency\_to\_p\_scale()} applica il mapping.

\textbf{Functional clauses / inherited \texttt{normativeClause [1..*]}}\par
\begin{tabularx}{\linewidth}{L{52mm}Y}
\toprule
\textbf{Condizione} & \textbf{Priorita' richiesta} \\
\midrule
\texttt{HIGH} + \texttt{confidence > 0.85} & DermaTriage MUST produrre \texttt{P1}. \\
\texttt{HIGH}, negli altri casi applicabili & DermaTriage MUST produrre \texttt{P2}. \\
\texttt{MEDIUM} & DermaTriage MUST produrre \texttt{P3}. \\
\texttt{LOW} & DermaTriage MUST produrre \texttt{P4}. \\
\bottomrule
\end{tabularx}

\textbf{Riferimenti strutturati / evidenze governate}\par
\emph{Adaptation Layer}; \texttt{adaptation\_layer.py}; \texttt{map\_urgency\_to\_p\_scale()}.

\textbf{SpecializedRequirement relation}\par
Nessuna relazione stabilita.
\end{ddtacore}

\begin{ddtaquestions}
\small
\begin{itemize}
  \item Quando l'urgenza proviene dal percorso symptom-only, si applica lo stesso mapping verso P1-P4 oppure una regola diversa?
  \item La soglia \texttt{confidence > 0.85} e' una regola governata del comportamento oppure un parametro configurabile?
  \item I valori temporali associati a P1-P4 hanno natura normativa di SLA? Quale trigger e quale owner li governano?
  \item Come viene trattata una confidence assente, non valida o fuori dal dominio atteso?
\end{itemize}
\end{ddtaquestions}

\switchcolumn

{\sffamily\large\bfseries\color{DDTABlue}Base Analysis}\par
\vspace{0.4em}

\begin{ddtareferentidentified}
\scriptsize
--
\end{ddtareferentidentified}

\begin{ddtareferentaccepted}
\scriptsize
\BAAccepted{DermaTriage}; \BAAccepted{CaseUrgency}; \BAAccepted{OperationalTriagePriority}; \BAAccepted{TriageConfidence}; \BAAccepted{OperationalPriorityDerivation} \BAFirstAccepted{}; \BAAccepted{AdaptationLayer} \BAFirstAccepted{}.\par
\end{ddtareferentaccepted}

\begin{ddtapropositionidentified}[top=3pt,bottom=3pt,left=5pt,right=5pt]
\fontsize{4.55}{4.95}\selectfont
\texttt{--}\par
\emph{Nessuna BAProposition aggiuntiva: la result-selection structure R7 qualifica localmente l'owner \texttt{produce} accettato sotto.}
\end{ddtapropositionidentified}

\begin{ddtapropositionaccepted}[top=3pt,bottom=3pt,left=5pt,right=5pt]
\fontsize{4.12}{4.48}\selectfont
\texttt{BAP-FR04-01-01}\par
\texttt{semanticOperatorKey = produce}\par
\hspace*{1em}\texttt{actor -> OperationalPriorityDerivation}\par
\hspace*{1em}\texttt{input -> CaseUrgency}\par
\hspace*{1em}\texttt{input -> TriageConfidence}\par
\hspace*{1em}\texttt{result -> OperationalTriagePriority}\par
\hspace*{1em}\texttt{scopedModifier}\par
\hspace*{2em}\texttt{condition -> LesionImage available}\par
\hspace*{1em}\texttt{operatorStructure [CANDIDATE R7 / NOT ADMITTED]}\par
\hspace*{2em}\texttt{decisionRule}\par
\hspace*{3em}\texttt{branch}\par
\hspace*{4em}\texttt{when -> CaseUrgency eq HIGH AND TriageConfidence gt 0.85}\par
\hspace*{4em}\texttt{resultAssignment -> OperationalTriagePriority = P1}\par
\hspace*{3em}\texttt{branch}\par
\hspace*{4em}\texttt{when -> other applicable cases within CaseUrgency eq HIGH}\par
\hspace*{4em}\texttt{resultAssignment -> OperationalTriagePriority = P2}\par
\hspace*{3em}\texttt{branch: CaseUrgency eq MEDIUM -> P3}\par
\hspace*{3em}\texttt{branch: CaseUrgency eq LOW -> P4}\par
\texttt{polarity = affirmative}\par
\emph{Acceptance boundary: la proposition \texttt{produce} e' accettata; la struttura annidata e' mostrata dentro il suo owner ma resta candidate R7 e non e' ancora parte della semantica BA ammessa. \texttt{eq} e' normalizzazione nominale; \texttt{gt} e' la sola ordered comparison candidate. Il residuale HIGH non equivale a \texttt{confidence <= 0.85}, non e' un ELSE globale e non copre automaticamente confidence missing/invalid/out-of-domain.}\par
\smallskip
\texttt{BAP-FR04-01-02}\par
\texttt{semanticOperatorKey = realize}\par
\hspace*{1em}\texttt{abstract -> OperationalPriorityDerivation}\par
\hspace*{1em}\texttt{realization -> AdaptationLayer}\par
\texttt{polarity = affirmative}
\end{ddtapropositionaccepted}

\end{paracol}
%</PAGE-CONTENT:13>
\clearpage
%<PAGE-DESC:14>MR-02 - Indicazione della destinazione specialistica
%<PAGE-CONTENT:14>
\PageHashFooter{\PageMDfive{14}}
\section{MR-02 - Indicazione della destinazione specialistica}

\setlength{\columnsep}{7mm}
\setlength{\columnseprule}{0.35pt}
\columnratio{0.67}
\begin{paracol}{2}

{\sffamily\large\bfseries\color{DDTANavy}Documentazione DDTA}\par
\vspace{0.5em}

\begin{ddtacore}[title={MacroRequirement}]
\textbf{Intent}\par
Determinare un'indicazione di destinazione specialistica per il caso dermatologico, a supporto del successivo instradamento verso l'assistenza appropriata.

\textbf{Context}\par
DermaTriage associa al risultato del processo di triage un'informazione orientata alla destinazione specialistica. La documentazione disponibile non stabilisce il vocabolario delle destinazioni, la regola di selezione ne' la responsabilita' di assegnazione o prenotazione.

\textbf{Stakeholders}\par
--

\textbf{Scope}\par
\textbf{IN:} produzione di un'indicazione di destinazione specialistica per il caso.\par
\textbf{OUT:} prenotazione; assegnazione effettiva; tassonomia completa delle destinazioni; regola completa di selezione e fallback.

\textbf{Assumptions / Constraints}\par
--

\textbf{dependsOn}\par
\texttt{None}
\end{ddtacore}

\begin{ddtacheck}[title={STOP AT MR}]
Il significato macro e' sufficientemente distinguibile, ma la documentazione disponibile non governa commitment ulteriori sufficienti per una decomposizione fedele. Il branch si arresta intenzionalmente a MacroRequirement.
\end{ddtacheck}

\begin{ddtaquestions}
\small
\begin{itemize}
  \item Qual e' il vocabolario governato delle destinazioni specialistiche?
  \item E' documentata una regola di selezione della destinazione?
  \item Chi possiede l'assegnazione o la prenotazione effettiva?
\end{itemize}
\end{ddtaquestions}


\switchcolumn

{\sffamily\large\bfseries\color{DDTABlue}Base Analysis}\par
\vspace{0.5em}

\begin{ddtareferentidentified}
\scriptsize
\BACandidate{SpecialistDestinationIndication} \BANew{}.
\end{ddtareferentidentified}

\begin{ddtareferentaccepted}
\scriptsize
\BAAccepted{DermaTriage}; \BAAccepted{SpecialistDestinationIndication} \BAFirstAccepted{}.
\end{ddtareferentaccepted}

\begin{ddtapropositionidentified}
\scriptsize
\texttt{BAP-MR02-01}\par
\texttt{semanticOperatorKey = produce}\par
\hspace*{1em}\texttt{actor -> DermaTriage}\par
\hspace*{1em}\texttt{result -> SpecialistDestinationIndication}\par
\texttt{polarity = affirmative}
\end{ddtapropositionidentified}

\begin{ddtapropositionaccepted}
\scriptsize
\textbf{BAP-MR02-01}.\par
\emph{Acceptance boundary: \texttt{SpecialistDestinationIndication} rappresenta l'indicazione specialistica governata a livello macro. La proposition non introduce input non documentati e non implica selezione di uno specifico professionista, assegnazione, prenotazione, routing effettivo o presa in carico.}
\end{ddtapropositionaccepted}

\begin{ddtaopen}[title={Pressure di projection - composizione / system boundary [DEFERRED]}]
\fontsize{5.45}{6.05}\selectfont
\texttt{DermaTriage} denota il software/sistema complessivo, non un singolo componente interno. Per questo \texttt{BAP-MR02-01} conserva \texttt{actor -> DermaTriage}: non viene sostituito con un componente interno soltanto per semplificare la projection. La rappresentazione grafica dovra' distinguere il whole system dagli elementi interni. La convenzione di lavoro preferita e' un \textbf{rettangolo esterno tratteggiato} etichettato \texttt{DermaTriage}, con al suo interno gli elementi la cui appartenenza al sistema sia documentata. Il nesting grafico non costituisce ancora una proposition \texttt{contains}/\texttt{partOf} ne' ammette un nuovo operatore BA. Il caso resta come pressione per riesaminare una possibile struttura riusabile di composition/membership dopo il completamento dell'analisi della documentazione con i costrutti R7 gia' disponibili.
\end{ddtaopen}

\end{paracol}
%</PAGE-CONTENT:14>
\clearpage
%<PAGE-DESC:15>MR-03 - Gestione della validazione clinica degli esiti
%<PAGE-CONTENT:15>
\PageHashFooter{\PageMDfive{15}}
\section{MR-03 - Gestione della validazione clinica degli esiti}

\setlength{\columnsep}{7mm}
\setlength{\columnseprule}{0.35pt}
\columnratio{0.67}
\begin{paracol}{2}

{\sffamily\large\bfseries\color{DDTANavy}Documentazione DDTA}\par
\vspace{0.35em}

\begin{ddtacore}[title={MacroRequirement}]
\fontsize{6.15}{6.75}\selectfont
\begin{tabularx}{\linewidth}{L{38mm}Y}
\toprule
\textbf{ID} & \texttt{MR-03} \\
\textbf{Lifecycle} & \texttt{current} \\
\textbf{Authority} & \texttt{CURRENT\_GOVERNED} \\
\textbf{Type} & \texttt{MacroRequirement} \\
\textbf{Title} & Gestione della validazione clinica degli esiti \\
\bottomrule
\end{tabularx}

\textbf{Intent}\par
Supportare la registrazione e l'utilizzo della validazione o correzione effettuata da un professionista sanitario sugli esiti prodotti da DermaTriage, mantenendo distinguibile il risultato della revisione dall'esito originario.

\textbf{Context}\par
Il giudizio clinico appartiene al professionista sanitario. DermaTriage possiede la responsabilita' di gestire, registrare e correlare l'esito della revisione al risultato originario del sistema, senza sostituirsi all'autorita' clinica. Il progetto supporta sia operazioni amministrative DermaTriage sia scambio di validazioni con il sistema esterno B4.

\textbf{Stakeholders}\par
Professionisti sanitari che effettuano la revisione; responsabili del servizio DermaTriage che utilizzano il risultato della revisione; processi downstream di adattamento che consumano evidence derivata dalla revisione.

\textbf{Scope}\par
\textbf{IN:} registrazione correlata del risultato della revisione; distinzione tra conferma e correzione; mantenimento della distinguibilita' tra esito originario e revisione.\par
\textbf{OUT:} decisione clinica stessa; contenuto completo del payload di correzione; policy completa di history/retention/finality; semantica completa del disaccordo usato downstream.

\textbf{Assumptions / Constraints}\par
La revisione clinica e' trattata come input proveniente da un professionista sanitario; il progetto non assume che conferma, correzione e disaccordo siano sinonimi.

\textbf{dependsOn}\par
Nessuna dipendenza macro canonica e' stabilita per \texttt{MR-03} nel baseline corrente.
\end{ddtacore}

\switchcolumn

{\sffamily\large\bfseries\color{DDTABlue}Base Analysis}\par
\vspace{0.35em}

\begin{ddtareferentidentified}[top=4pt,bottom=4pt,left=5pt,right=5pt]
\fontsize{5.25}{5.8}\selectfont
\BACandidate{ClinicalReviewResult} \BANew{}; \BACandidate{OriginalDermaTriageOutcome} \BANew{}; \BACandidate{HealthcareProfessional} \BANew{}.
\end{ddtareferentidentified}

\begin{ddtareferentaccepted}[top=4pt,bottom=4pt,left=5pt,right=5pt]
\fontsize{5.15}{5.65}\selectfont
\BAAccepted{DermaTriage}; \BAAccepted{ClinicalReviewResult} \BAFirstAccepted{}; \BAAccepted{OriginalDermaTriageOutcome} \BAFirstAccepted{}; \BAAccepted{HealthcareProfessional} \BAFirstAccepted{}.
\end{ddtareferentaccepted}

\begin{ddtapropositionidentified}[top=4pt,bottom=4pt,left=5pt,right=5pt]
\fontsize{4.95}{5.45}\selectfont
{\itshape
\texttt{BAP-MR03-01}\par
\texttt{semanticOperatorKey = transfer}\par
\hspace*{1em}\texttt{source -> HealthcareProfessional}\par
\hspace*{1em}\texttt{destination -> DermaTriage}\par
\hspace*{1em}\texttt{content -> ClinicalReviewResult}\par
\texttt{polarity = affirmative}\par
}
\end{ddtapropositionidentified}

\begin{ddtapropositionaccepted}[top=4pt,bottom=4pt,left=5pt,right=5pt]
\fontsize{5.15}{5.65}\selectfont
--\par
\emph{Il transfer resta candidato: MR-03 stabilisce l'origine clinica della revisione, ma non ancora il delivery path tecnico concreto.}
\end{ddtapropositionaccepted}

\begin{ddtaopenrelations}[top=4pt,bottom=4pt,left=5pt,right=5pt]
\fontsize{4.7}{5.2}\selectfont
\texttt{DermaTriage <record / persistence da decomporre> ClinicalReviewResult}\par
\texttt{OriginalDermaTriageOutcome <distinguishable from> ClinicalReviewResult}\par
La same-case association fra revisione ed esito originario resta un concern documentale/data-contract: non viene materializzata come BAProposition autonoma finche' non emerge una relazione indipendentemente targetable. Nessun nuovo operatore viene introdotto in questo checkpoint.
\end{ddtaopenrelations}

\end{paracol}
%</PAGE-CONTENT:15>
\clearpage
%<PAGE-DESC:16>MR-03 > DEC-03 - Separazione tra esito originario e revisione clinica
%<PAGE-CONTENT:16>
\PageHashFooter{\PageMDfive{16}}
\subsection{DEC-03 - Separazione tra esito originario e revisione clinica}
{\sffamily\scriptsize\color{DDTAGray}Gerarchia: \texttt{MR-03} $\rightarrow$ \texttt{DEC-03}}\par

\setlength{\columnsep}{7mm}
\setlength{\columnseprule}{0.35pt}
\columnratio{0.67}
\begin{paracol}{2}

{\sffamily\large\bfseries\color{DDTANavy}Documentazione DDTA}\par
\vspace{0.45em}

\begin{ddtacore}[title={Decision}]
\fontsize{6.3}{6.95}\selectfont
\begin{tabularx}{\linewidth}{L{38mm}Y}
\toprule
\textbf{ID} & \texttt{DEC-03} \\
\textbf{Lifecycle} & \texttt{current} \\
\textbf{Authority} & \texttt{CURRENT\_GOVERNED} \\
\textbf{Type} & \texttt{Decision} \\
\textbf{Title} & Separazione tra esito originario e revisione clinica \\
\textbf{Parent MacroRequirement} & \texttt{MR-03} \\
\textbf{Decision kind} & \texttt{SELECTABLE} \\
\bottomrule
\end{tabularx}

\textbf{Context}\par
Una revisione clinica puo' confermare o correggere un esito prodotto da DermaTriage. Sovrascrivere semanticamente l'esito originario renderebbe indistinguibili il comportamento del sistema e il successivo giudizio clinico.

\textbf{Decision}\par
DermaTriage mantiene l'esito originario distinguibile dal successivo risultato della revisione clinica.

\textbf{Consequences}\par
L'esito originario e il risultato della revisione clinica restano distinguibili. Una conferma dell'esito originario resta distinguibile da una sua correzione.
\end{ddtacore}

\switchcolumn

{\sffamily\large\bfseries\color{DDTABlue}Base Analysis}\par
\vspace{0.45em}

\begin{ddtareferentidentified}
\scriptsize
--
\end{ddtareferentidentified}

\begin{ddtareferentaccepted}
\scriptsize
\BAAccepted{DermaTriage}; \BAAccepted{OriginalDermaTriageOutcome}; \BAAccepted{ClinicalReviewResult}.
\end{ddtareferentaccepted}

\begin{ddtapropositionidentified}
\scriptsize
--
\end{ddtapropositionidentified}

\begin{ddtapropositionaccepted}
\scriptsize
--
\end{ddtapropositionaccepted}

\begin{ddtaopenrelations}
\fontsize{5.15}{5.7}\selectfont
\texttt{OriginalDermaTriageOutcome <distinguishableFrom> ClinicalReviewResult}\par
La distinzione e' un fatto governato, ma non viene forzata in \texttt{constrain}: la semantica e' relazionale tra due identita' e R7 non dispone ancora di un operatore dedicato. \texttt{DermaTriage} preserva tale distinzione.
\end{ddtaopenrelations}

\end{paracol}
%</PAGE-CONTENT:16>
\clearpage
%<PAGE-DESC:17>MR-03 > DEC-03 > FR-03 - Registrazione correlata della revisione clinica
%<PAGE-CONTENT:17>
\PageHashFooter{\PageMDfive{17}}
\subsubsection{FR-03 - Registrazione correlata della revisione clinica}
{\sffamily\scriptsize\color{DDTAGray}Gerarchia: \texttt{MR-03} $\rightarrow$ \texttt{DEC-03} $\rightarrow$ \texttt{FR-03}}\par

\setlength{\columnsep}{7mm}
\setlength{\columnseprule}{0.35pt}
\columnratio{0.67}
\begin{paracol}{2}

{\sffamily\large\bfseries\color{DDTANavy}Documentazione DDTA}\par
\vspace{0.4em}

\begin{ddtacore}[title={FunctionalRequirement}]
\fontsize{6.15}{6.75}\selectfont
\begin{tabularx}{\linewidth}{L{38mm}Y}
\toprule
\textbf{ID} & \texttt{FR-03} \\
\textbf{Lifecycle} & \texttt{current} \\
\textbf{Authority} & \texttt{CURRENT\_GOVERNED} \\
\textbf{Type} & \texttt{FunctionalRequirement} \\
\textbf{Title} & Registrazione correlata della revisione clinica \\
\textbf{Parent Decision} & \texttt{DEC-03} \\
\textbf{Owning MR (derived)} & \texttt{MR-03} \\
\textbf{SpecializedRequirement relation} & 0 active instances \\
\bottomrule
\end{tabularx}

\textbf{Functional obligation}\par
Registrare il risultato della revisione clinica mantenendone la correlazione con l'esito DermaTriage originario.

\textbf{Functional clauses / inherited \texttt{normativeClause [1..*]}}\par
Quando viene fornita una validazione o correzione clinica relativa a un esito DermaTriage, DermaTriage MUST registrare il risultato della revisione associandolo all'esito originario cui la revisione si riferisce.

\textbf{SPO references (L2 support)}\par
\begin{lstlisting}
DermaTriage -- record --> ClinicalReviewResult
ClinicalReviewResult -- associated with --> OriginalDermaTriageOutcome
\end{lstlisting}
\end{ddtacore}

\switchcolumn

{\sffamily\large\bfseries\color{DDTABlue}Base Analysis}\par
\vspace{0.4em}

\begin{ddtareferentidentified}
\scriptsize
--
\end{ddtareferentidentified}

\begin{ddtareferentaccepted}
\scriptsize
\BAAccepted{DermaTriage}; \BAAccepted{ClinicalReviewResult}; \BAAccepted{OriginalDermaTriageOutcome}.
\end{ddtareferentaccepted}

\begin{ddtapropositionidentified}
\scriptsize
--
\end{ddtapropositionidentified}

\begin{ddtapropositionaccepted}
\scriptsize
--
\end{ddtapropositionaccepted}

\begin{ddtaopenrelations}
\fontsize{5.0}{5.5}\selectfont
\texttt{DermaTriage <record / persistence> ClinicalReviewResult}\par
\medskip
\textbf{Registration pressure.} \texttt{record} non viene assimilato a \texttt{produce}. Un eventuale \texttt{transfer} verso storage verra' introdotto solo quando la destination di storage sara' governata; le operazioni interne dello storage non sono anticipate.\par
\textbf{Case-binding.} L'associazione della review allo specifico esito originario resta informazione governata, ma viene trattata come concern di identity/data contract e non come BAProposition autonoma. Verra' riaperta solo se evidence successiva mostra una relazione indipendentemente targetable.
\end{ddtaopenrelations}

\end{paracol}
%</PAGE-CONTENT:17>
\clearpage
%<PAGE-DESC:18>MR-03 > DEC-03 > FR-12 - Distinzione tra conferma e correzione della revisione clinica
%<PAGE-CONTENT:18>
\PageHashFooter{\PageMDfive{18}}
\subsubsection{FR-12 - Distinzione tra conferma e correzione della revisione clinica}
{\sffamily\scriptsize\color{DDTAGray}Gerarchia: \texttt{MR-03} $\rightarrow$ \texttt{DEC-03} $\rightarrow$ \texttt{FR-12}}\par

\setlength{\columnsep}{7mm}
\setlength{\columnseprule}{0.35pt}
\columnratio{0.67}
\begin{paracol}{2}

{\sffamily\large\bfseries\color{DDTANavy}Documentazione DDTA}\par
\vspace{0.4em}

\begin{ddtacore}[title={FunctionalRequirement}]
\fontsize{6.15}{6.75}\selectfont
\begin{tabularx}{\linewidth}{L{38mm}Y}
\toprule
\textbf{ID} & \texttt{FR-12} \\
\textbf{Lifecycle} & \texttt{current} \\
\textbf{Authority} & \texttt{CURRENT\_GOVERNED} \\
\textbf{Type} & \texttt{FunctionalRequirement} \\
\textbf{Title} & Distinzione tra conferma e correzione della revisione clinica \\
\textbf{Parent Decision} & \texttt{DEC-03} \\
\textbf{Owning MR (derived)} & \texttt{MR-03} \\
\textbf{SpecializedRequirement relation} & 0 active instances \\
\bottomrule
\end{tabularx}

\textbf{Functional obligation}\par
Rendere distinguibile se il risultato della revisione clinica conferma oppure corregge l'esito originario.

\textbf{Functional clauses / inherited \texttt{normativeClause [1..*]}}\par
Per ogni risultato di revisione clinica registrato, DermaTriage MUST mantenere distinguibile una conferma dell'esito originario da una correzione dell'esito originario.

\textbf{SPO references (L2 support)}\par
\begin{lstlisting}
ClinicalReviewResult -- classify disposition as --> ConfirmationOrCorrection
DermaTriage -- preserve distinction of --> ClinicalReviewDisposition
\end{lstlisting}
\end{ddtacore}

\switchcolumn

{\sffamily\large\bfseries\color{DDTABlue}Base Analysis}\par
\vspace{0.4em}

\begin{ddtareferentidentified}
\scriptsize
\BACandidate{ClinicalReviewDisposition} \BANew{}.
\end{ddtareferentidentified}

\begin{ddtareferentaccepted}
\fontsize{5.2}{5.7}\selectfont
\BAAccepted{DermaTriage}; \BAAccepted{ClinicalReviewResult}; \BAAccepted{OriginalDermaTriageOutcome}; \BAAccepted{ClinicalReviewDisposition} \BAFirstAccepted{}.
\end{ddtareferentaccepted}

\begin{ddtapropositionidentified}
\fontsize{5.0}{5.5}\selectfont
\texttt{BAP-FR12-01}\par
\texttt{semanticOperatorKey = constrain}\par
\hspace*{1em}\texttt{constraintTarget -> ClinicalReviewResult}\par
\hspace*{1em}\texttt{constraintValue}\par
\hspace*{2em}\texttt{property -> clinicalReviewDisposition}\par
\hspace*{2em}\texttt{vocabulary -> [confirmation, correction]}\par
\texttt{polarity = affirmative}
\end{ddtapropositionidentified}

\begin{ddtapropositionaccepted}
\scriptsize
\textbf{BAP-FR12-01}.
\end{ddtapropositionaccepted}

\begin{ddtaopenrelations}
\fontsize{5.05}{5.55}\selectfont
La proposition governa il contratto di distinguibilita', non la regola che seleziona runtime un valore. \texttt{disagreement} non viene aggiunto al vocabulary: MR-03 vieta di assumerlo sinonimo di conferma/correzione, mentre FR-12 governa soltanto questi due valori.
\end{ddtaopenrelations}

\end{paracol}
%</PAGE-CONTENT:18>
\clearpage
%<PAGE-DESC:19>MR-03 > DEC-15 - Scambio della validazione clinica con B4
%<PAGE-CONTENT:19>
\PageHashFooter{\PageMDfive{19}}
\subsection{DEC-15 - Scambio della validazione clinica con B4}
{\sffamily\scriptsize\color{DDTAGray}Gerarchia: \texttt{MR-03} $\rightarrow$ \texttt{DEC-15}}\par

\setlength{\columnsep}{7mm}
\setlength{\columnseprule}{0.35pt}
\columnratio{0.67}
\begin{paracol}{2}

{\sffamily\large\bfseries\color{DDTANavy}Documentazione DDTA}\par
\vspace{0.4em}

\begin{ddtacore}[title={Decision}]
\fontsize{6.1}{6.7}\selectfont
\begin{tabularx}{\linewidth}{L{38mm}Y}
\toprule
\textbf{ID} & \texttt{DEC-15} \\
\textbf{Lifecycle} & \texttt{current} \\
\textbf{Authority} & \texttt{CURRENT\_GOVERNED} \\
\textbf{Type} & \texttt{Decision} \\
\textbf{Title} & Scambio della validazione clinica con B4 \\
\textbf{Parent MacroRequirement} & \texttt{MR-03} \\
\textbf{Decision kind} & \texttt{SELECTABLE} \\
\bottomrule
\end{tabularx}

\textbf{Context}\par
Nel workflow integrato, la review clinica deve rimanere correlata alla consultation e all'esito di triage scambiati con B4.

\textbf{Decision}\par
DermaTriage usa le API B4 per il diagnostic-output write-back, il medical-validation write-back e il retrieval dei validated outcomes, mantenendo l'identita' della consultation/case durante lo scambio.

\textbf{Consequences}\par
B4 resta un sistema esterno consumato; la responsabilita' di registrare e correlare la review resta di DermaTriage.
\end{ddtacore}

\switchcolumn

{\sffamily\large\bfseries\color{DDTABlue}Base Analysis}\par
\vspace{0.4em}

\begin{ddtareferentidentified}[top=3pt,bottom=3pt,left=5pt,right=5pt]
\fontsize{4.9}{5.35}\selectfont
\BACandidate{B4} \BANew{}; \BACandidate{B4TriageOutcome} \BANew{}; \BACandidate{ClinicianValidation} \BANew{}; \BACandidate{ValidatedOutcome} \BANew{}.
\end{ddtareferentidentified}

\begin{ddtareferentaccepted}[top=3pt,bottom=3pt,left=5pt,right=5pt]
\fontsize{5.0}{5.45}\selectfont
\BAAccepted{DermaTriage}; \BAAccepted{B4} \BAFirstAccepted{}.
\end{ddtareferentaccepted}

\begin{ddtapropositionidentified}[top=3pt,bottom=3pt,left=5pt,right=5pt]
\fontsize{4.15}{4.55}\selectfont
{\itshape
\texttt{BAP-DEC15-01}\par
\texttt{semanticOperatorKey = transfer}\par
\hspace*{1em}\texttt{source -> DermaTriage}\par
\hspace*{1em}\texttt{destination -> B4}\par
\hspace*{1em}\texttt{content -> B4TriageOutcome}\par
\texttt{polarity = affirmative}\par
\medskip
\texttt{BAP-DEC15-02}\par
\texttt{semanticOperatorKey = transfer}\par
\hspace*{1em}\texttt{source -> DermaTriage}\par
\hspace*{1em}\texttt{destination -> B4}\par
\hspace*{1em}\texttt{content -> ClinicianValidation}\par
\texttt{polarity = affirmative}\par
\medskip
\texttt{BAP-DEC15-03}\par
\texttt{semanticOperatorKey = transfer}\par
\hspace*{1em}\texttt{source -> B4}\par
\hspace*{1em}\texttt{destination -> DermaTriage}\par
\hspace*{1em}\texttt{content -> ValidatedOutcome}\par
\texttt{polarity = affirmative}
}
\end{ddtapropositionidentified}

\begin{ddtapropositionaccepted}[top=3pt,bottom=3pt,left=5pt,right=5pt]
\fontsize{4.8}{5.25}\selectfont
--\par
\emph{Le tre conveyance sono candidate a livello Decision; FR-23 consolida source, destination e content.}
\end{ddtapropositionaccepted}

\begin{ddtaopenrelations}[top=3pt,bottom=3pt,left=5pt,right=5pt]
\fontsize{4.6}{5.05}\selectfont
La preservazione della stessa consultation/case identity e' mantenuta come qualifier documentale/data-contract e non diventa una proposition di correlazione autonoma. L'uso delle API B4 resta evidence del meccanismo di scambio; non viene introdotto un \texttt{realize} senza una distinta identita' di behavior/componente.
\end{ddtaopenrelations}

\end{paracol}
%</PAGE-CONTENT:19>
\clearpage
%<PAGE-DESC:20>MR-03 > DEC-15 > FR-23 - Write-back e retrieval della validazione B4
%<PAGE-CONTENT:20>
\PageHashFooter{\PageMDfive{20}}
\subsubsection{FR-23 - Write-back e retrieval della validazione B4}
{\sffamily\scriptsize\color{DDTAGray}Gerarchia: \texttt{MR-03} $\rightarrow$ \texttt{DEC-15} $\rightarrow$ \texttt{FR-23}}\par

\setlength{\columnsep}{7mm}
\setlength{\columnseprule}{0.35pt}
\columnratio{0.67}
\begin{paracol}{2}

{\sffamily\large\bfseries\color{DDTANavy}Documentazione DDTA}\par
\vspace{0.35em}

\begin{ddtacore}[title={FunctionalRequirement}]
\fontsize{5.9}{6.45}\selectfont
\begin{tabularx}{\linewidth}{L{38mm}Y}
\toprule
\textbf{ID} & \texttt{FR-23} \\
\textbf{Lifecycle} & \texttt{current} \\
\textbf{Authority} & \texttt{CURRENT\_GOVERNED} \\
\textbf{Type} & \texttt{FunctionalRequirement} \\
\textbf{Title} & Write-back e retrieval della validazione B4 \\
\textbf{Parent Decision} & \texttt{DEC-15} \\
\textbf{Owning MR (derived)} & \texttt{MR-03} \\
\textbf{SpecializedRequirement relation} & 0 active instances \\
\bottomrule
\end{tabularx}

\textbf{Functional obligation}\par
Scambiare con B4 l'esito diagnostico/di triage e la successiva validazione mantenendo la correlazione con lo stesso caso.

\textbf{Functional clauses / inherited \texttt{normativeClause [1..*]}}\par
Nel workflow B4-integrated, DermaTriage MUST mantenere la stessa consultation/case identity durante diagnostic-output write-back, medical-validation write-back e retrieval dei validated outcomes, cosi' che la revisione resti associabile all'esito originario cui si riferisce.

\textbf{SPO references (L2 support)}\par
\begin{lstlisting}
DermaTriage -- writes back --> B4TriageOutcome
ClinicianValidation -- written back to --> B4
B4 -- provides --> ValidatedOutcome
ValidatedOutcome -- correlatedWith --> OriginalDermaTriageOutcome
\end{lstlisting}
\end{ddtacore}

\switchcolumn

{\sffamily\large\bfseries\color{DDTABlue}Base Analysis}\par
\vspace{0.35em}

\begin{ddtareferentidentified}[top=3pt,bottom=3pt,left=5pt,right=5pt]
\fontsize{4.9}{5.35}\selectfont
--
\end{ddtareferentidentified}

\begin{ddtareferentaccepted}[top=3pt,bottom=3pt,left=5pt,right=5pt]
\fontsize{4.7}{5.15}\selectfont
\BAAccepted{DermaTriage}; \BAAccepted{B4}; \BAAccepted{B4TriageOutcome} \BAFirstAccepted{}; \BAAccepted{ClinicianValidation} \BAFirstAccepted{}; \BAAccepted{ValidatedOutcome} \BAFirstAccepted{}; \BAAccepted{OriginalDermaTriageOutcome}.
\end{ddtareferentaccepted}

\begin{ddtapropositionidentified}[top=3pt,bottom=3pt,left=5pt,right=5pt]
\fontsize{4.15}{4.55}\selectfont
\texttt{BAP-FR23-01}\par
\texttt{semanticOperatorKey = transfer}\par
\hspace*{1em}\texttt{source -> DermaTriage}\par
\hspace*{1em}\texttt{destination -> B4}\par
\hspace*{1em}\texttt{content -> B4TriageOutcome}\par
\texttt{polarity = affirmative}\par
\medskip
\texttt{BAP-FR23-02}\par
\texttt{semanticOperatorKey = transfer}\par
\hspace*{1em}\texttt{source -> DermaTriage}\par
\hspace*{1em}\texttt{destination -> B4}\par
\hspace*{1em}\texttt{content -> ClinicianValidation}\par
\texttt{polarity = affirmative}\par
\medskip
\texttt{BAP-FR23-03}\par
\texttt{semanticOperatorKey = transfer}\par
\hspace*{1em}\texttt{source -> B4}\par
\hspace*{1em}\texttt{destination -> DermaTriage}\par
\hspace*{1em}\texttt{content -> ValidatedOutcome}\par
\texttt{polarity = affirmative}
\end{ddtapropositionidentified}

\begin{ddtapropositionaccepted}[top=3pt,bottom=3pt,left=5pt,right=5pt]
\fontsize{4.8}{5.25}\selectfont
\textbf{BAP-FR23-01}; \textbf{BAP-FR23-02}; \textbf{BAP-FR23-03}.
\end{ddtapropositionaccepted}

\begin{ddtaopenrelations}[top=3pt,bottom=3pt,left=5pt,right=5pt]
\fontsize{4.45}{4.9}\selectfont
La same-case identity resta un concern di contratto/binding, non una BAProposition autonoma. Non viene inoltre assunta equivalenza semantica fra \texttt{B4TriageOutcome} e \texttt{OriginalDermaTriageOutcome}, ne' fra \texttt{ClinicianValidation}/\texttt{ValidatedOutcome} e \texttt{ClinicalReviewResult}: tali identity relation restano \texttt{NOT SPECIFIED} e nessun mapping viene inventato.
\end{ddtaopenrelations}

\end{paracol}
%</PAGE-CONTENT:20>
\clearpage
%<PAGE-DESC:21>MR-03 > DEC-16 - Autenticazione delle interfacce di review
%<PAGE-CONTENT:21>
\PageHashFooter{\PageMDfive{21}}
\subsection{DEC-16 - Autenticazione delle interfacce di review}
{\sffamily\scriptsize\color{DDTAGray}Gerarchia: \texttt{MR-03} $\rightarrow$ \texttt{DEC-16}}\par

\setlength{\columnsep}{7mm}
\setlength{\columnseprule}{0.35pt}
\columnratio{0.67}
\begin{paracol}{2}

{\sffamily\large\bfseries\color{DDTANavy}Documentazione DDTA}\par
\vspace{0.35em}

\begin{ddtacore}[title={Decision}]
\fontsize{5.9}{6.45}\selectfont
\begin{tabularx}{\linewidth}{L{38mm}Y}
\toprule
\textbf{ID} & \texttt{DEC-16} \\
\textbf{Lifecycle} & \texttt{current} \\
\textbf{Authority} & \texttt{CURRENT\_GOVERNED} \\
\textbf{Type} & \texttt{Decision} \\
\textbf{Title} & Autenticazione delle interfacce di review \\
\textbf{Parent MacroRequirement} & \texttt{MR-03} \\
\textbf{Decision kind} & \texttt{SELECTABLE} \\
\bottomrule
\end{tabularx}

\textbf{Context}\par
Le operazioni amministrative di review/correction e l'interazione con B4 attraversano boundary differenti e usano meccanismi di autenticazione distinti.

\textbf{Decision}\par
Le operazioni amministrative DermaTriage usano \texttt{X-API-Key}; l'interazione con B4 usa bearer JWT, con acquisizione e refresh del token gestiti dal client B4.

\textbf{Consequences}\par
Le operazioni amministrative DermaTriage e l'interazione con B4 usano meccanismi di autenticazione distinti.
\end{ddtacore}

\switchcolumn

{\sffamily\large\bfseries\color{DDTABlue}Base Analysis}\par
\vspace{0.35em}

\begin{ddtareferentidentified}[top=3pt,bottom=3pt,left=5pt,right=5pt]
\fontsize{4.8}{5.25}\selectfont
\BACandidate{ProtectedDermaTriageOperation} \BANew{}; \BACandidate{B4CallAuthentication} \BANew{}.
\end{ddtareferentidentified}

\begin{ddtareferentaccepted}[top=3pt,bottom=3pt,left=5pt,right=5pt]
\fontsize{5.0}{5.45}\selectfont
\BAAccepted{DermaTriage}; \BAAccepted{B4}.
\end{ddtareferentaccepted}

\begin{ddtapropositionidentified}[top=3pt,bottom=3pt,left=5pt,right=5pt]
\fontsize{4.15}{4.55}\selectfont
{\itshape
\texttt{BAP-DEC16-01}\par
\texttt{semanticOperatorKey = constrain}\par
\hspace*{1em}\texttt{constraintTarget -> ProtectedDermaTriageOperation}\par
\hspace*{1em}\texttt{constraintValue}\par
\hspace*{2em}\texttt{property -> authenticationMechanism}\par
\hspace*{2em}\texttt{vocabulary -> [X-API-Key]}\par
\texttt{polarity = affirmative}\par
\medskip
\texttt{BAP-DEC16-02}\par
\texttt{semanticOperatorKey = constrain}\par
\hspace*{1em}\texttt{constraintTarget -> B4CallAuthentication}\par
\hspace*{1em}\texttt{constraintValue}\par
\hspace*{2em}\texttt{property -> authenticationMechanism}\par
\hspace*{2em}\texttt{vocabulary -> [bearer JWT]}\par
\texttt{polarity = affirmative}
}
\end{ddtapropositionidentified}

\begin{ddtapropositionaccepted}[top=3pt,bottom=3pt,left=5pt,right=5pt]
\fontsize{4.7}{5.15}\selectfont
--\par
\emph{Le due restrizioni restano candidate a livello Decision. FR-24 consolida il target delle operazioni amministrative protette; FR-25 conserva la funzione generale di autenticazione B4 e la sua realizzazione concreta.}
\end{ddtapropositionaccepted}

\end{paracol}
%</PAGE-CONTENT:21>
\clearpage
%<PAGE-DESC:22>MR-03 > DEC-16 > FR-24 - X-API-Key per operazioni amministrative DermaTriage
%<PAGE-CONTENT:22>
\PageHashFooter{\PageMDfive{22}}
\subsubsection{FR-24 - X-API-Key per operazioni amministrative DermaTriage}
{\sffamily\scriptsize\color{DDTAGray}Gerarchia: \texttt{MR-03} $\rightarrow$ \texttt{DEC-16} $\rightarrow$ \texttt{FR-24}}\par

\setlength{\columnsep}{7mm}
\setlength{\columnseprule}{0.35pt}
\columnratio{0.67}
\begin{paracol}{2}

{\sffamily\large\bfseries\color{DDTANavy}Documentazione DDTA}\par
\vspace{0.35em}

\begin{ddtacore}[title={FunctionalRequirement}]
\fontsize{5.8}{6.35}\selectfont
\begin{tabularx}{\linewidth}{L{38mm}Y}
\toprule
\textbf{ID} & \texttt{FR-24} \\
\textbf{Lifecycle} & \texttt{current} \\
\textbf{Authority} & \texttt{CURRENT\_GOVERNED} \\
\textbf{Type} & \texttt{FunctionalRequirement} \\
\textbf{Title} & X-API-Key per operazioni amministrative DermaTriage \\
\textbf{Parent Decision} & \texttt{DEC-16} \\
\textbf{Owning MR (derived)} & \texttt{MR-03} \\
\textbf{SpecializedRequirement relation} & 0 active instances \\
\bottomrule
\end{tabularx}

\textbf{Functional obligation}\par
Richiedere la API key prevista per le operazioni amministrative protette.

\textbf{Functional clauses / inherited \texttt{normativeClause [1..*]}}\par
Le operazioni amministrative DermaTriage di review, validation/correction, prompt management e classifier retraining esposte come protette MUST richiedere il meccanismo \texttt{X-API-Key} previsto dall'interfaccia.

\textbf{SPO references (L2 support)}\par
\begin{lstlisting}
AdministrativeClient -- presents --> X-API-Key
ProtectedDermaTriageOperation -- requires --> X-API-Key
\end{lstlisting}
\end{ddtacore}

\switchcolumn

{\sffamily\large\bfseries\color{DDTABlue}Base Analysis}\par
\vspace{0.35em}

\begin{ddtareferentidentified}[top=3pt,bottom=3pt,left=5pt,right=5pt]
\fontsize{4.9}{5.35}\selectfont
--
\end{ddtareferentidentified}

\begin{ddtareferentaccepted}[top=3pt,bottom=3pt,left=5pt,right=5pt]
\fontsize{4.8}{5.25}\selectfont
\BAAccepted{DermaTriage}; \BAAccepted{ProtectedDermaTriageOperation} \BAFirstAccepted{}.
\end{ddtareferentaccepted}

\begin{ddtapropositionidentified}[top=3pt,bottom=3pt,left=5pt,right=5pt]
\fontsize{4.25}{4.65}\selectfont
\texttt{BAP-FR24-01}\par
\texttt{semanticOperatorKey = constrain}\par
\hspace*{1em}\texttt{constraintTarget -> ProtectedDermaTriageOperation}\par
\hspace*{1em}\texttt{constraintValue}\par
\hspace*{2em}\texttt{property -> authenticationMechanism}\par
\hspace*{2em}\texttt{presence -> required}\par
\hspace*{2em}\texttt{vocabulary -> [X-API-Key]}\par
\texttt{polarity = affirmative}
\end{ddtapropositionidentified}

\begin{ddtapropositionaccepted}[top=3pt,bottom=3pt,left=5pt,right=5pt]
\fontsize{4.8}{5.25}\selectfont
\textbf{BAP-FR24-01}.
\end{ddtapropositionaccepted}

\begin{ddtaopenrelations}[top=3pt,bottom=3pt,left=5pt,right=5pt]
\fontsize{4.55}{5.0}\selectfont
\texttt{AdministrativeClient -- presents --> X-API-Key} resta evidenza documentale del meccanismo di autenticazione e non viene materializzato come transfer autonomo: FR-24 non richiede una distinta identita' BA del client amministrativo per esprimere il vincolo governato.
\end{ddtaopenrelations}

\end{paracol}
%</PAGE-CONTENT:22>
\clearpage
%<PAGE-DESC:23>MR-03 > DEC-16 > FR-25 - Bearer JWT per interazione B4
%<PAGE-CONTENT:23>
\PageHashFooter{\PageMDfive{23}}
\subsubsection{FR-25 - Bearer JWT per interazione B4}
{\sffamily\scriptsize\color{DDTAGray}Gerarchia: \texttt{MR-03} $\rightarrow$ \texttt{DEC-16} $\rightarrow$ \texttt{FR-25}}\par

\setlength{\columnsep}{7mm}
\setlength{\columnseprule}{0.35pt}
\columnratio{0.67}
\begin{paracol}{2}

{\sffamily\large\bfseries\color{DDTANavy}Documentazione DDTA}\par
\vspace{0.35em}

\begin{ddtacore}[title={FunctionalRequirement}]
\fontsize{5.75}{6.3}\selectfont
\begin{tabularx}{\linewidth}{L{38mm}Y}
\toprule
\textbf{ID} & \texttt{FR-25} \\
\textbf{Lifecycle} & \texttt{current} \\
\textbf{Authority} & \texttt{CURRENT\_GOVERNED} \\
\textbf{Type} & \texttt{FunctionalRequirement} \\
\textbf{Title} & Bearer JWT per interazione B4 \\
\textbf{Parent Decision} & \texttt{DEC-16} \\
\textbf{Owning MR (derived)} & \texttt{MR-03} \\
\textbf{SpecializedRequirement relation} & 0 active instances \\
\bottomrule
\end{tabularx}

\textbf{Functional obligation}\par
Autenticare le chiamate DermaTriage verso B4 mediante il flusso bearer JWT documentato.

\textbf{Functional clauses / inherited \texttt{normativeClause [1..*]}}\par
Il client B4 di DermaTriage MUST acquisire e rinnovare il token previsto dal flusso di autenticazione B4 e MUST usare il bearer JWT nelle chiamate B4 che richiedono autenticazione.

\textbf{SPO references (L2 support)}\par
\begin{lstlisting}
DermaTriageB4Client -- obtains/refreshes --> B4BearerJWT
DermaTriageB4Client -- presents --> B4BearerJWT
\end{lstlisting}
\end{ddtacore}

\switchcolumn

{\sffamily\large\bfseries\color{DDTABlue}Base Analysis}\par
\vspace{0.35em}

\begin{ddtareferentidentified}[top=3pt,bottom=3pt,left=5pt,right=5pt]
\fontsize{4.65}{5.1}\selectfont
\BACandidate{DermaTriageB4Client} \BANew{}; \BACandidate{B4BearerJWT} \BANew{}.
\end{ddtareferentidentified}

\begin{ddtareferentaccepted}[top=3pt,bottom=3pt,left=5pt,right=5pt]
\fontsize{4.55}{5.0}\selectfont
\BAAccepted{DermaTriage}; \BAAccepted{B4}; \BAAccepted{B4CallAuthentication} \BAFirstAccepted{}; \BAAccepted{DermaTriageB4Client} \BAFirstAccepted{}; \BAAccepted{B4BearerJWT} \BAFirstAccepted{}.
\end{ddtareferentaccepted}

\begin{ddtapropositionidentified}[top=3pt,bottom=3pt,left=5pt,right=5pt]
\fontsize{3.95}{4.35}\selectfont
\texttt{BAP-FR25-01}\par
\texttt{semanticOperatorKey = realize}\par
\hspace*{1em}\texttt{abstract -> B4CallAuthentication}\par
\hspace*{1em}\texttt{realization -> DermaTriageB4Client}\par
\texttt{polarity = affirmative}\par
\medskip
\texttt{BAP-FR25-02}\par
\texttt{semanticOperatorKey = constrain}\par
\hspace*{1em}\texttt{constraintTarget -> B4CallAuthentication}\par
\hspace*{1em}\texttt{constraintValue}\par
\hspace*{2em}\texttt{property -> authenticationMechanism}\par
\hspace*{2em}\texttt{vocabulary -> [bearer JWT]}\par
\texttt{polarity = affirmative}\par
\medskip
{\itshape
\texttt{BAP-FR25-03}\par
\texttt{semanticOperatorKey = transfer}\par
\hspace*{1em}\texttt{source -> ??}\par
\hspace*{1em}\texttt{destination -> DermaTriageB4Client}\par
\hspace*{1em}\texttt{content -> B4BearerJWT}\par
\texttt{polarity = affirmative}
}
\end{ddtapropositionidentified}

\begin{ddtapropositionaccepted}[top=3pt,bottom=3pt,left=5pt,right=5pt]
\fontsize{4.55}{5.0}\selectfont
\textbf{BAP-FR25-01}; \textbf{BAP-FR25-02}.\par
\emph{BAP-FR25-03 resta candidate: il transfer del token e' riconosciuto, ma la source non e' specificata.}
\end{ddtapropositionaccepted}

\begin{ddtaopenrelations}[top=3pt,bottom=3pt,left=5pt,right=5pt]
\fontsize{4.15}{4.55}\selectfont
\texttt{source -> ??} registra un gap documentale, non un'entita' inferita: authorization server/token endpoint restano \texttt{NOT SPECIFIED}. Lo SPO \texttt{presents} e' mantenuto come evidenza dell'uso della credenziale nel comportamento \texttt{B4CallAuthentication} e non genera automaticamente un ulteriore transfer.
\end{ddtaopenrelations}

\end{paracol}
%</PAGE-CONTENT:23>
\clearpage
%<PAGE-DESC:24>MR-04 - Adattamento controllato sulla base della revisione clinica
%<PAGE-CONTENT:24>
\PageHashFooter{\PageMDfive{24}}
\section{MR-04 - Adattamento controllato sulla base della revisione clinica}

\setlength{\columnsep}{7mm}
\setlength{\columnseprule}{0.35pt}
\columnratio{0.67}
\begin{paracol}{2}

{\sffamily\large\bfseries\color{DDTANavy}Documentazione DDTA}\par
\vspace{0.35em}

\begin{ddtacore}[title={MacroRequirement}]
\fontsize{5.65}{6.15}\selectfont
\begin{tabularx}{\linewidth}{L{38mm}Y}
\toprule
\textbf{ID} & \texttt{MR-04} \\
\textbf{Lifecycle} & \texttt{current} \\
\textbf{Authority} & \texttt{CURRENT\_GOVERNED} \\
\textbf{Type} & \texttt{MacroRequirement} \\
\textbf{Title} & Adattamento controllato sulla base della revisione clinica \\
\bottomrule
\end{tabularx}

\textbf{Intent}\par
Consentire l'evoluzione controllata del comportamento di DermaTriage utilizzando evidence derivata dalle validazioni e correzioni cliniche accumulate, evitando che modifiche degradanti vengano accettate senza controllo.

\textbf{Context}\par
DermaTriage utilizza risultati di revisione clinica per supportare l'evoluzione di capability differenti, in particolare il percorso di evoluzione dei prompt usati da Qwen2-VL e BioMistral e il percorso di adattamento di EfficientNet-B4. I due percorsi hanno condizioni di attivazione, evidence qualificante e risultati di lifecycle distinti.

\textbf{Stakeholders}\par
Responsabili dell'evoluzione e manutenzione di DermaTriage; professionisti sanitari le cui revisioni forniscono evidence clinica; utilizzatori e soggetti influenzati dalla qualita' del triage dopo l'adattamento.

\textbf{Scope}\par
\textbf{IN:} attivazione su accumulo di evidence; separazione dei percorsi di adattamento; selezione dell'evidence recente; qualificazione del disaccordo clinico; derivazione del target di supervisione; valutazione comparativa; proprieta' qualitative di accettazione; reversibilita' post-adoption.\par
\textbf{OUT:} giudizio clinico posseduto da \texttt{MR-03}; autorita' finale di deployment non governata.

\textbf{Assumptions / Constraints}\par
Le revisioni cliniche consumate da questo branch sono governate da \texttt{MR-03}. I percorsi di adattamento possono riusare dati solo quando tali dati soddisfano indipendentemente le regole di ciascun percorso.

\textbf{dependsOn}\par
\texttt{MR-04 dependsOn MR-03}. La dipendenza esprime consumo del significato di revisione clinica posseduto da \texttt{MR-03}; non crea co-ownership dei Requirement.
\end{ddtacore}

\switchcolumn

{\sffamily\large\bfseries\color{DDTABlue}Base Analysis}\par
\vspace{0.35em}

\begin{ddtareferentidentified}[top=3pt,bottom=3pt,left=5pt,right=5pt]
\fontsize{4.8}{5.25}\selectfont
\texttt{--}
\end{ddtareferentidentified}

\begin{ddtareferentaccepted}[top=3pt,bottom=3pt,left=5pt,right=5pt]
\fontsize{4.55}{5.0}\selectfont
\BAAccepted{DermaTriage}; \BAAccepted{ClinicalReviewResult}; \BAAccepted{Qwen2-VL-7B-Instruct}; \BAAccepted{BioMistral-7B}; \BAAccepted{EfficientNet-B4}.\par
\BAAccepted{ControlledDermaTriageAdaptation} \BAFirstAccepted{}; \BAAccepted{PromptEvolutionPath} \BAFirstAccepted{}; \BAAccepted{ClassifierAdaptationPath} \BAFirstAccepted{}.
\end{ddtareferentaccepted}

\begin{ddtapropositionidentified}[top=3pt,bottom=3pt,left=5pt,right=5pt]
\fontsize{4.45}{4.9}\selectfont
\texttt{--}\par
\emph{Il MacroRequirement governa l'identita' della responsabilita' di adattamento e di due percorsi distinti, ma non stabilisce ancora una relation BA sufficientemente precisa tra la responsabilita' macro e i percorsi.}
\end{ddtapropositionidentified}

\begin{ddtapropositionaccepted}[top=3pt,bottom=3pt,left=5pt,right=5pt]
\fontsize{4.55}{5.0}\selectfont
\texttt{--}
\end{ddtapropositionaccepted}

\begin{ddtaopenrelations}[top=3pt,bottom=3pt,left=5pt,right=5pt]
\fontsize{4.15}{4.55}\selectfont
I due percorsi hanno activation condition, qualifying evidence e lifecycle result distinti. Non si introduce ancora \texttt{realize}, \texttt{contains}/\texttt{partOf} o una struttura di aggregation: la relazione fra \texttt{ControlledDermaTriageAdaptation} e i due percorsi resta aperta fino all'analisi dei branch figli. \texttt{MR-04 dependsOn MR-03} resta relazione documentale di consumo del significato clinico e non diventa automaticamente una BAProposition. L'autorita' finale di deployment resta \texttt{NOT GOVERNED}.
\end{ddtaopenrelations}

\end{paracol}
%</PAGE-CONTENT:24>
\clearpage
%<PAGE-DESC:25>MR-04 > DEC-04 - Attivazione dell'adattamento su accumulo di evidence clinica
%<PAGE-CONTENT:25>
\PageHashFooter{\PageMDfive{25}}
\subsection{DEC-04 - Attivazione dell'adattamento su accumulo di evidence clinica}
{\sffamily\scriptsize\color{DDTAGray}Gerarchia: \texttt{MR-04} $\rightarrow$ \texttt{DEC-04}}\par

\setlength{\columnsep}{7mm}
\setlength{\columnseprule}{0.35pt}
\columnratio{0.67}
\begin{paracol}{2}

{\sffamily\large\bfseries\color{DDTANavy}Documentazione DDTA}\par
\vspace{0.35em}

\begin{ddtacore}[title={Decision}]
\fontsize{5.8}{6.35}\selectfont
\begin{tabularx}{\linewidth}{L{40mm}Y}
\toprule
\textbf{ID} & \texttt{DEC-04} \\
\textbf{Lifecycle} & \texttt{current} \\
\textbf{Authority} & \texttt{CURRENT\_GOVERNED} \\
\textbf{Type} & \texttt{Decision} \\
\textbf{Title} & Attivazione dell'adattamento su accumulo di evidence clinica \\
\textbf{Parent MacroRequirement} & \texttt{MR-04} \\
\textbf{Decision kind} & \texttt{SELECTABLE} \\
\bottomrule
\end{tabularx}

\textbf{Context}\par
L'adattamento non viene attivato per ogni singola revisione clinica. I diversi percorsi devono accumulare evidence pertinente prima di avviare il proprio ciclo di evoluzione.

\textbf{Decision}\par
Ogni processo di adattamento viene attivato soltanto quando l'evidence clinica pertinente accumulata soddisfa la condizione di accumulo governata per quel processo.

\textbf{Consequences}\par
Il percorso prompt e il percorso classificatorio hanno soglie distinte: 10 correzioni pertinenti per il prompt e 50 correzioni qualificanti per il classifier. \texttt{FR-04} e \texttt{FR-05} governano separatamente i due trigger.
\end{ddtacore}

\switchcolumn

{\sffamily\large\bfseries\color{DDTABlue}Base Analysis}\par
\vspace{0.35em}

\begin{ddtareferentidentified}[top=3pt,bottom=3pt,left=5pt,right=5pt]
\fontsize{4.7}{5.15}\selectfont
\texttt{--}
\end{ddtareferentidentified}

\begin{ddtareferentaccepted}[top=3pt,bottom=3pt,left=5pt,right=5pt]
\fontsize{4.55}{5.0}\selectfont
\BAAccepted{DermaTriage}; \BAAccepted{PromptEvolutionPath}; \BAAccepted{ClassifierAdaptationPath}.
\end{ddtareferentaccepted}

\begin{ddtapropositionidentified}[top=3pt,bottom=3pt,left=5pt,right=5pt]
\fontsize{4.3}{4.75}\selectfont
\texttt{--}\par
\emph{La Decision governa un activation guard per ciascun percorso, ma non fornisce ancora una proposition owner alla quale applicare lo scoped modifier \texttt{condition}.}
\end{ddtapropositionidentified}

\begin{ddtapropositionaccepted}[top=3pt,bottom=3pt,left=5pt,right=5pt]
\fontsize{4.55}{5.0}\selectfont
\texttt{--}
\end{ddtapropositionaccepted}

\begin{ddtaopenrelations}[top=3pt,bottom=3pt,left=5pt,right=5pt]
\fontsize{4.05}{4.45}\selectfont
Governed facts: il percorso prompt e' attivato soltanto quando la propria condizione di accumulo e' soddisfatta; lo stesso vale indipendentemente per il percorso classificatorio. Le soglie concrete sono 10 e 50 e vengono possedute dai figli \texttt{FR-04}/\texttt{FR-05}. La lettura corrente e' \emph{condition in testa al branch}, non un nuovo \texttt{trigger} e non una relation \texttt{initiate} gia' giustificata. Riemerge la pressure per una struttura riusabile di aggregate behavior capace di possedere identita', membership e activation condition; nessun \texttt{contains}/\texttt{partOf} viene ammesso in questa pagina.
\end{ddtaopenrelations}

\end{paracol}
%</PAGE-CONTENT:25>
\clearpage
%<PAGE-DESC:26>MR-04 > DEC-04 > FR-04 - Attivazione del percorso di evoluzione del prompt
%<PAGE-CONTENT:26>
\PageHashFooter{\PageMDfive{26}}
\subsubsection{FR-04 - Attivazione del percorso di evoluzione del prompt}
{\sffamily\scriptsize\color{DDTAGray}Gerarchia: \texttt{MR-04} $\rightarrow$ \texttt{DEC-04} $\rightarrow$ \texttt{FR-04}}\par

\setlength{\columnsep}{7mm}
\setlength{\columnseprule}{0.35pt}
\columnratio{0.67}
\begin{paracol}{2}

{\sffamily\large\bfseries\color{DDTANavy}Documentazione DDTA}\par
\vspace{0.3em}

\begin{ddtacore}[title={FunctionalRequirement}]
\fontsize{5.55}{6.1}\selectfont
\begin{tabularx}{\linewidth}{L{39mm}Y}
\toprule
\textbf{ID} & \texttt{FR-04} \\
\textbf{Lifecycle} & \texttt{current} \\
\textbf{Authority} & \texttt{CURRENT\_GOVERNED} \\
\textbf{Type} & \texttt{FunctionalRequirement} \\
\textbf{Title} & Attivazione del percorso di evoluzione del prompt \\
\textbf{Parent Decision} & \texttt{DEC-04} \\
\textbf{Owning MR (derived)} & \texttt{MR-04} \\
\textbf{SpecializedRequirement relation} & 0 active instances \\
\bottomrule
\end{tabularx}

\textbf{Functional obligation}\par
Attivare un ciclo di evoluzione del prompt quando l'evidence di correzione clinica pertinente raggiunge la soglia applicabile al percorso prompt.

\textbf{Functional clauses / inherited \texttt{normativeClause [1..*]}}\par
Quando l'evidence di correzione clinica pertinente accumulata raggiunge 10 elementi, DermaTriage MUST attivare un ciclo di evoluzione del prompt.

\textbf{SPO references}\par
\texttt{AccumulatedPromptCorrectionEvidence -- reaches --> PromptEvolutionThreshold}\par
\texttt{DermaTriage -- activate --> PromptEvolutionCycle}
\end{ddtacore}

\switchcolumn

{\sffamily\large\bfseries\color{DDTABlue}Base Analysis}\par
\vspace{0.35em}

\begin{ddtareferentidentified}[top=3pt,bottom=3pt,left=5pt,right=5pt]
\fontsize{4.55}{5.0}\selectfont
\BACandidate{PromptEvolutionCycle} \BANew{}.
\end{ddtareferentidentified}

\begin{ddtareferentaccepted}[top=3pt,bottom=3pt,left=5pt,right=5pt]
\fontsize{4.55}{5.0}\selectfont
\BAAccepted{DermaTriage}; \BAAccepted{PromptEvolutionPath}.
\end{ddtareferentaccepted}

\begin{ddtapropositionidentified}[top=3pt,bottom=3pt,left=5pt,right=5pt]
\fontsize{4.3}{4.75}\selectfont
\texttt{--}\par
\emph{Il FunctionalRequirement specializza la condizione di attivazione ma non aggiunge una relation BA autonoma sufficientemente giustificata.}
\end{ddtapropositionidentified}

\begin{ddtapropositionaccepted}[top=3pt,bottom=3pt,left=5pt,right=5pt]
\fontsize{4.55}{5.0}\selectfont
\texttt{--}
\end{ddtapropositionaccepted}

\begin{ddtaopenrelations}[top=3pt,bottom=3pt,left=5pt,right=5pt]
\fontsize{4.0}{4.42}\selectfont
Governed fact: il ciclo prompt viene attivato quando l'evidence di correzione clinica pertinente accumulata \emph{raggiunge 10 elementi}. Non si inferisce \texttt{>= 10}, periodicita', reset, re-trigger o comportamento oltre la soglia. \texttt{PromptEvolutionThreshold} e \texttt{AccumulatedPromptCorrectionEvidence} restano nomi SPO e non vengono promossi automaticamente a BAReferent. \texttt{PromptEvolutionCycle} resta candidato finche' i documenti successivi non ne dimostrano membership/lifecycle riusabili. Il requisito non stabilisce componente, file, classe, funzione o altro binding al codice.
\end{ddtaopenrelations}

\end{paracol}
%</PAGE-CONTENT:26>
\clearpage
%<PAGE-DESC:27>MR-04 > DEC-04 > FR-05 - Attivazione del percorso di adattamento classificatorio
%<PAGE-CONTENT:27>
\PageHashFooter{\PageMDfive{27}}
\subsubsection{FR-05 - Attivazione del percorso di adattamento classificatorio}
{\sffamily\scriptsize\color{DDTAGray}Gerarchia: \texttt{MR-04} $\rightarrow$ \texttt{DEC-04} $\rightarrow$ \texttt{FR-05}}\par

\setlength{\columnsep}{7mm}
\setlength{\columnseprule}{0.35pt}
\columnratio{0.67}
\begin{paracol}{2}

{\sffamily\large\bfseries\color{DDTANavy}Documentazione DDTA}\par
\vspace{0.3em}

\begin{ddtacore}[title={FunctionalRequirement}]
\fontsize{5.45}{5.95}\selectfont
\begin{tabularx}{\linewidth}{L{39mm}Y}
\toprule
\textbf{ID} & \texttt{FR-05} \\
\textbf{Lifecycle} & \texttt{current} \\
\textbf{Authority} & \texttt{CURRENT\_GOVERNED} \\
\textbf{Type} & \texttt{FunctionalRequirement} \\
\textbf{Title} & Attivazione del percorso di adattamento classificatorio \\
\textbf{Parent Decision} & \texttt{DEC-04} \\
\textbf{Owning MR (derived)} & \texttt{MR-04} \\
\textbf{SpecializedRequirement relation} & 0 active instances \\
\bottomrule
\end{tabularx}

\textbf{Functional obligation}\par
Attivare il ciclo di adattamento della classificazione quando l'evidence qualificante accumulata raggiunge la soglia applicabile al percorso classificatorio.

\textbf{Functional clauses / inherited \texttt{normativeClause [1..*]}}\par
Quando l'evidence qualificante accumulata raggiunge 50 correzioni, DermaTriage MUST attivare il ciclo di adattamento classificatorio di EfficientNet-B4.

\textbf{SPO references}\par
\texttt{AccumulatedClassifierEvidence -- reaches --> ClassifierAdaptationThreshold}\par
\texttt{DermaTriage -- activate --> ClassifierAdaptationCycle}

\textbf{Cross-branch / consumption note}\par
\texttt{FR-05} consuma evidence di disaccordo qualificata secondo \texttt{FR-07}. Questo consumo non modifica la parent Decision di nessuno dei due Requirement.
\end{ddtacore}

\switchcolumn

{\sffamily\large\bfseries\color{DDTABlue}Base Analysis}\par
\vspace{0.35em}

\begin{ddtareferentidentified}[top=3pt,bottom=3pt,left=5pt,right=5pt]
\fontsize{4.55}{5.0}\selectfont
\BACandidate{ClassifierAdaptationCycle} \BANew{}.
\end{ddtareferentidentified}

\begin{ddtareferentaccepted}[top=3pt,bottom=3pt,left=5pt,right=5pt]
\fontsize{4.55}{5.0}\selectfont
\BAAccepted{DermaTriage}; \BAAccepted{ClassifierAdaptationPath}; \BAAccepted{EfficientNet-B4}.
\end{ddtareferentaccepted}

\begin{ddtapropositionidentified}[top=3pt,bottom=3pt,left=5pt,right=5pt]
\fontsize{4.3}{4.75}\selectfont
\texttt{--}\par
\emph{Il FunctionalRequirement specializza il trigger classificatorio e identifica il target tecnico EfficientNet-B4, ma non determina ancora una relation BA autonoma verso il modello o il codice.}
\end{ddtapropositionidentified}

\begin{ddtapropositionaccepted}[top=3pt,bottom=3pt,left=5pt,right=5pt]
\fontsize{4.55}{5.0}\selectfont
\texttt{--}
\end{ddtapropositionaccepted}

\begin{ddtaopenrelations}[top=3pt,bottom=3pt,left=5pt,right=5pt]
\fontsize{3.92}{4.34}\selectfont
Governed facts: il ciclo classificatorio viene attivato quando l'evidence qualificante accumulata raggiunge 50 correzioni; il ciclo riguarda l'adattamento di \texttt{EfficientNet-B4}; l'evidence qualificante consumata da \texttt{FR-05} e' quella di disaccordo qualificata secondo \texttt{FR-07}. \texttt{ClassifierAdaptationThreshold} e \texttt{AccumulatedClassifierEvidence} non vengono promossi automaticamente a BAReferent. Non si inferiscono \texttt{>= 50}, reset/re-trigger, routine di training, file, classe, funzione o altro binding al codice. La relazione formale fra \texttt{ClassifierAdaptationCycle} ed \texttt{EfficientNet-B4} resta aperta fino ai documenti che governano il retraining.
\end{ddtaopenrelations}

\end{paracol}
%</PAGE-CONTENT:27>
\clearpage
%<PAGE-DESC:28>MR-04 > DEC-05 - Separazione dei percorsi di adattamento per capability
%<PAGE-CONTENT:28>
\PageHashFooter{\PageMDfive{28}}
\subsection{DEC-05 - Separazione dei percorsi di adattamento per capability}
{\sffamily\scriptsize\color{DDTAGray}Gerarchia: \texttt{MR-04} $\rightarrow$ \texttt{DEC-05}}\par

\setlength{\columnsep}{7mm}
\setlength{\columnseprule}{0.35pt}
\columnratio{0.67}
\begin{paracol}{2}

{\sffamily\large\bfseries\color{DDTANavy}Documentazione DDTA}\par
\vspace{0.35em}

\begin{ddtacore}[title={Decision}]
\fontsize{5.55}{6.05}\selectfont
\begin{tabularx}{\linewidth}{L{40mm}Y}
\toprule
\textbf{ID} & \texttt{DEC-05} \\
\textbf{Lifecycle} & \texttt{current} \\
\textbf{Authority} & \texttt{CURRENT\_GOVERNED} \\
\textbf{Type} & \texttt{Decision} \\
\textbf{Title} & Separazione dei percorsi di adattamento per capability \\
\textbf{Parent MacroRequirement} & \texttt{MR-04} \\
\textbf{Decision kind} & \texttt{SELECTABLE} \\
\bottomrule
\end{tabularx}

\textbf{Context}\par
L'evoluzione del prompt e l'adattamento della classificazione sono capability distinte. Condividere dati o trovarsi nello stesso progetto non rende equivalenti le loro condizioni di attivazione, le loro regole di qualificazione dell'evidence o i loro risultati di lifecycle.

\textbf{Decision}\par
DermaTriage governa i percorsi di adattamento come lifecycle distinti: condizioni di attivazione, evidence qualificante e risultati di qualificazione/accettazione non si trasferiscono automaticamente da un percorso all'altro.

\textbf{Consequences}\par
Le condizioni di attivazione, la qualificazione dell'evidence e i risultati di lifecycle restano indipendenti tra i diversi percorsi di adattamento.

\textbf{Lifecycle note}\par
\texttt{FR-11 - Indipendenza dei lifecycle di adattamento} e' \textbf{superseded}. Il suo ID resta preservato come identita' storica; la baseline corrente usa \texttt{FR-13}, \texttt{FR-14} e \texttt{FR-15} per obbligazioni indipendentemente revisionabili e assessable.
\end{ddtacore}

\switchcolumn

{\sffamily\large\bfseries\color{DDTABlue}Base Analysis}\par
\vspace{0.35em}

\begin{ddtareferentidentified}[top=3pt,bottom=3pt,left=5pt,right=5pt]
\fontsize{4.7}{5.15}\selectfont
\texttt{--}
\end{ddtareferentidentified}

\begin{ddtareferentaccepted}[top=3pt,bottom=3pt,left=5pt,right=5pt]
\fontsize{4.55}{5.0}\selectfont
\BAAccepted{DermaTriage}; \BAAccepted{PromptEvolutionPath}; \BAAccepted{ClassifierAdaptationPath}.
\end{ddtareferentaccepted}

\begin{ddtapropositionidentified}[top=3pt,bottom=3pt,left=5pt,right=5pt]
\fontsize{4.4}{4.85}\selectfont
\texttt{--}\par
\emph{La separazione dei lifecycle e la non-propagazione dei qualifier sono significati governati, ma la Decision non fornisce un semanticOperatorKey corrente che li rappresenti senza inventare vocabulary o relation aggiuntive.}
\end{ddtapropositionidentified}

\begin{ddtapropositionaccepted}[top=3pt,bottom=3pt,left=5pt,right=5pt]
\fontsize{4.55}{5.0}\selectfont
\texttt{--}
\end{ddtapropositionaccepted}

\begin{ddtaopenrelations}[top=3pt,bottom=3pt,left=5pt,right=5pt]
\fontsize{4.0}{4.4}\selectfont
Governed semantic boundary: \texttt{PromptEvolutionPath} e \texttt{ClassifierAdaptationPath} restano lifecycle distinti. Membership nello stesso progetto o futuro aggregate non implica condivisione automatica di activation condition, evidence qualification o lifecycle result. Nessun \texttt{constrain}, \texttt{transfer}, \texttt{contains}/\texttt{partOf} viene introdotto per codificare artificialmente la non-propagazione. \texttt{FR-13}/\texttt{FR-14}/\texttt{FR-15} vengono analizzati come decomposizioni normative della stessa regola; il ramo resta candidato a review documentale di consolidamento dopo la chiusura di \texttt{MR-04}.
\end{ddtaopenrelations}

\end{paracol}
%</PAGE-CONTENT:28>
\clearpage
%<PAGE-DESC:29>MR-04 > DEC-05 > FR-13 - Indipendenza delle condizioni di attivazione tra percorsi di adattamento
%<PAGE-CONTENT:29>
\PageHashFooter{\PageMDfive{29}}
\subsubsection{FR-13 - Indipendenza delle condizioni di attivazione tra percorsi di adattamento}
{\sffamily\scriptsize\color{DDTAGray}Gerarchia: \texttt{MR-04} $\rightarrow$ \texttt{DEC-05} $\rightarrow$ \texttt{FR-13}}\par

\setlength{\columnsep}{7mm}
\setlength{\columnseprule}{0.35pt}
\columnratio{0.67}
\begin{paracol}{2}

{\sffamily\large\bfseries\color{DDTANavy}Documentazione DDTA}\par
\vspace{0.3em}

\begin{ddtacore}[title={FunctionalRequirement}]
\fontsize{5.45}{5.95}\selectfont
\begin{tabularx}{\linewidth}{L{39mm}Y}
\toprule
\textbf{ID} & \texttt{FR-13} \\
\textbf{Lifecycle} & \texttt{current} \\
\textbf{Authority} & \texttt{CURRENT\_GOVERNED} \\
\textbf{Type} & \texttt{FunctionalRequirement} \\
\textbf{Title} & Indipendenza delle condizioni di attivazione tra percorsi di adattamento \\
\textbf{Parent Decision} & \texttt{DEC-05} \\
\textbf{Owning MR (derived)} & \texttt{MR-04} \\
\textbf{SpecializedRequirement relation} & 0 active instances \\
\bottomrule
\end{tabularx}

\textbf{Functional obligation}\par
Impedire che la condizione di attivazione soddisfatta da un percorso sia considerata, da sola, sufficiente ad attivare un altro percorso.

\textbf{Functional clauses / inherited \texttt{normativeClause [1..*]}}\par
Il soddisfacimento della condizione di attivazione di un percorso di adattamento MUST NOT, di per se', essere trattato come soddisfacimento della condizione di attivazione di un altro percorso di adattamento.

\textbf{SPO references (L2 support)}\par
\texttt{ActivationConditionOfPathA -- MUST NOT imply -->}\par
\texttt{ActivationConditionOfPathB}
\end{ddtacore}

\switchcolumn

{\sffamily\large\bfseries\color{DDTABlue}Base Analysis}\par
\vspace{0.35em}

\begin{ddtareferentidentified}[top=3pt,bottom=3pt,left=5pt,right=5pt]
\fontsize{4.7}{5.15}\selectfont
\texttt{--}\par
\emph{\texttt{ActivationConditionOfPathA/B} sono placeholder dello SPO e non vengono promossi a BAReferent.}
\end{ddtareferentidentified}

\begin{ddtareferentaccepted}[top=3pt,bottom=3pt,left=5pt,right=5pt]
\fontsize{4.55}{5.0}\selectfont
\BAAccepted{PromptEvolutionPath}; \BAAccepted{ClassifierAdaptationPath}.
\end{ddtareferentaccepted}

\begin{ddtapropositionidentified}[top=3pt,bottom=3pt,left=5pt,right=5pt]
\fontsize{4.45}{4.9}\selectfont
\texttt{--}
\end{ddtapropositionidentified}

\begin{ddtapropositionaccepted}[top=3pt,bottom=3pt,left=5pt,right=5pt]
\fontsize{4.55}{5.0}\selectfont
\texttt{--}
\end{ddtapropositionaccepted}

\begin{ddtaopenrelations}[top=3pt,bottom=3pt,left=5pt,right=5pt]
\fontsize{4.05}{4.45}\selectfont
Governed semantic fact: il soddisfacimento della activation condition di un adaptation path non e', da solo, sufficiente a considerare soddisfatta la condition di un altro path. Il significato e' una regola di scope/non-propagazione, non una nuova funzione BA. Non si inferisce che i path non possano condividere prerequisite o che una condition non possa contribuire a un'altra quando una diversa base documentale lo governi. Rispetto a \texttt{DEC-05}, \texttt{FR-13} aggiunge soprattutto l'obbligazione normativa atomica; resta candidato a consolidamento documentale futuro.
\end{ddtaopenrelations}

\end{paracol}
%</PAGE-CONTENT:29>
\clearpage
%<PAGE-DESC:30>MR-04 > DEC-05 > FR-14 - Indipendenza della qualificazione dell'evidence tra percorsi di adattamento
%<PAGE-CONTENT:30>
\PageHashFooter{\PageMDfive{30}}
\subsubsection{FR-14 - Indipendenza della qualificazione dell'evidence tra percorsi di adattamento}
{\sffamily\scriptsize\color{DDTAGray}Gerarchia: \texttt{MR-04} $\rightarrow$ \texttt{DEC-05} $\rightarrow$ \texttt{FR-14}}\par

\setlength{\columnsep}{7mm}
\setlength{\columnseprule}{0.35pt}
\columnratio{0.67}
\begin{paracol}{2}

{\sffamily\large\bfseries\color{DDTANavy}Documentazione DDTA}\par
\vspace{0.3em}

\begin{ddtacore}[title={FunctionalRequirement}]
\fontsize{5.45}{5.95}\selectfont
\begin{tabularx}{\linewidth}{L{39mm}Y}
\toprule
\textbf{ID} & \texttt{FR-14} \\
\textbf{Lifecycle} & \texttt{current} \\
\textbf{Authority} & \texttt{CURRENT\_GOVERNED} \\
\textbf{Type} & \texttt{FunctionalRequirement} \\
\textbf{Title} & Indipendenza della qualificazione dell'evidence tra percorsi di adattamento \\
\textbf{Parent Decision} & \texttt{DEC-05} \\
\textbf{Owning MR (derived)} & \texttt{MR-04} \\
\textbf{SpecializedRequirement relation} & 0 active instances \\
\bottomrule
\end{tabularx}

\textbf{Functional obligation}\par
Impedire che evidence qualificante per un percorso sia considerata automaticamente qualificante per un altro percorso.

\textbf{Functional clauses / inherited \texttt{normativeClause [1..*]}}\par
Evidence che qualifica per un percorso di adattamento MUST NOT, di per se', essere trattata come evidence qualificante per un altro percorso di adattamento.

\textbf{SPO references (L2 support)}\par
\texttt{EvidenceQualifiedForPathA -- MUST NOT imply -->}\par
\texttt{EvidenceQualifiedForPathB}
\end{ddtacore}

\switchcolumn

{\sffamily\large\bfseries\color{DDTABlue}Base Analysis}\par
\vspace{0.35em}

\begin{ddtareferentidentified}[top=3pt,bottom=3pt,left=5pt,right=5pt]
\fontsize{4.7}{5.15}\selectfont
\texttt{--}\par
\emph{\texttt{EvidenceQualifiedForPathA/B} sono placeholder dello SPO e non vengono promossi a BAReferent.}
\end{ddtareferentidentified}

\begin{ddtareferentaccepted}[top=3pt,bottom=3pt,left=5pt,right=5pt]
\fontsize{4.55}{5.0}\selectfont
\BAAccepted{PromptEvolutionPath}; \BAAccepted{ClassifierAdaptationPath}.
\end{ddtareferentaccepted}

\begin{ddtapropositionidentified}[top=3pt,bottom=3pt,left=5pt,right=5pt]
\fontsize{4.45}{4.9}\selectfont
\texttt{--}
\end{ddtapropositionidentified}

\begin{ddtapropositionaccepted}[top=3pt,bottom=3pt,left=5pt,right=5pt]
\fontsize{4.55}{5.0}\selectfont
\texttt{--}
\end{ddtapropositionaccepted}

\begin{ddtaopenrelations}[top=3pt,bottom=3pt,left=5pt,right=5pt]
\fontsize{4.0}{4.4}\selectfont
Governed semantic fact: evidence qualificante per un adaptation path non diventa, per questo solo fatto, evidence qualificante per l'altro. La qualification resta path-local. Non si inferisce che lo stesso dato non possa essere riusato da entrambi i path: puo' esserlo soltanto quando soddisfa indipendentemente le regole governate di ciascun percorso. Nessun generic \texttt{constrain} o relation di propagation viene inventato. Rispetto a \texttt{DEC-05}, \texttt{FR-14} e' principalmente una decomposizione normativa atomica e resta candidato a consolidamento documentale futuro.
\end{ddtaopenrelations}

\end{paracol}
%</PAGE-CONTENT:30>
\clearpage
%<PAGE-DESC:31>MR-04 > DEC-05 > FR-15 - Indipendenza dei risultati di qualificazione e accettazione tra percorsi di adattamento
%<PAGE-CONTENT:31>
\PageHashFooter{\PageMDfive{31}}
\subsubsection{FR-15 - Indipendenza dei risultati di qualificazione e accettazione tra percorsi di adattamento}
{\sffamily\scriptsize\color{DDTAGray}Gerarchia: \texttt{MR-04} $\rightarrow$ \texttt{DEC-05} $\rightarrow$ \texttt{FR-15}}\par

\setlength{\columnsep}{7mm}
\setlength{\columnseprule}{0.35pt}
\columnratio{0.67}
\begin{paracol}{2}

{\sffamily\large\bfseries\color{DDTANavy}Documentazione DDTA}\par
\vspace{0.3em}

\begin{ddtacore}[title={FunctionalRequirement}]
\fontsize{5.35}{5.85}\selectfont
\begin{tabularx}{\linewidth}{L{39mm}Y}
\toprule
\textbf{ID} & \texttt{FR-15} \\
\textbf{Lifecycle} & \texttt{current} \\
\textbf{Authority} & \texttt{CURRENT\_GOVERNED} \\
\textbf{Type} & \texttt{FunctionalRequirement} \\
\textbf{Title} & Indipendenza dei risultati di qualificazione e accettazione tra percorsi di adattamento \\
\textbf{Parent Decision} & \texttt{DEC-05} \\
\textbf{Owning MR (derived)} & \texttt{MR-04} \\
\textbf{SpecializedRequirement relation} & 0 active instances \\
\bottomrule
\end{tabularx}

\textbf{Functional obligation}\par
Impedire che un risultato di qualificazione, accettazione o progressione raggiunto da un percorso venga trattato automaticamente come risultato equivalente per un altro percorso.

\textbf{Functional clauses / inherited \texttt{normativeClause [1..*]}}\par
Un risultato di qualificazione, accettazione o progressione raggiunto da un percorso di adattamento MUST NOT, di per se', essere trattato come risultato di lifecycle equivalente per un altro percorso di adattamento.

\textbf{SPO references (L2 support)}\par
\texttt{LifecycleResultOfPathA -- MUST NOT imply -->}\par
\texttt{EquivalentLifecycleResultOfPathB}
\end{ddtacore}

\switchcolumn

{\sffamily\large\bfseries\color{DDTABlue}Base Analysis}\par
\vspace{0.35em}

\begin{ddtareferentidentified}[top=3pt,bottom=3pt,left=5pt,right=5pt]
\fontsize{4.65}{5.1}\selectfont
\texttt{--}\par
\emph{\texttt{LifecycleResultOfPathA} e \texttt{EquivalentLifecycleResultOfPathB} sono placeholder dello SPO e non vengono promossi a BAReferent.}
\end{ddtareferentidentified}

\begin{ddtareferentaccepted}[top=3pt,bottom=3pt,left=5pt,right=5pt]
\fontsize{4.55}{5.0}\selectfont
\BAAccepted{PromptEvolutionPath}; \BAAccepted{ClassifierAdaptationPath}.
\end{ddtareferentaccepted}

\begin{ddtapropositionidentified}[top=3pt,bottom=3pt,left=5pt,right=5pt]
\fontsize{4.45}{4.9}\selectfont
\texttt{--}
\end{ddtapropositionidentified}

\begin{ddtapropositionaccepted}[top=3pt,bottom=3pt,left=5pt,right=5pt]
\fontsize{4.55}{5.0}\selectfont
\texttt{--}
\end{ddtapropositionaccepted}

\begin{ddtaopenrelations}[top=3pt,bottom=3pt,left=5pt,right=5pt]
\fontsize{3.95}{4.35}\selectfont
Governed semantic fact: un risultato di qualification, acceptance o progression raggiunto da un adaptation path non e', da solo, un risultato di lifecycle equivalente per l'altro. La regola non vieta risultati uguali o criteri condivisi quando siano governati indipendentemente; vieta soltanto l'equivalenza automatica. Nessuna nuova funzione BA, tecnologia o binding implementativo emerge. Con \texttt{FR-13} e \texttt{FR-14}, \texttt{FR-15} conferma il pattern di decomposizione normativa di \texttt{DEC-05}; il sotto-ramo e' un candidato forte per review documentale di merge/consolidamento dopo la chiusura di \texttt{MR-04}.
\end{ddtaopenrelations}

\end{paracol}
%</PAGE-CONTENT:31>
\clearpage

%<PAGE-DESC:32>MR-04 > DEC-06 - Accettazione comparativa dell'adattamento della classificazione di urgenza
%<PAGE-CONTENT:32>
\PageHashFooter{\PageMDfive{32}}
\subsection{DEC-06 - Accettazione comparativa dell'adattamento della classificazione di urgenza}
{\sffamily\scriptsize\color{DDTAGray}Gerarchia: \texttt{MR-04} $\rightarrow$ \texttt{DEC-06}}\par

\setlength{\columnsep}{7mm}
\setlength{\columnseprule}{0.35pt}
\columnratio{0.67}
\begin{paracol}{2}

{\sffamily\large\bfseries\color{DDTANavy}Documentazione DDTA}\par
\vspace{0.3em}

\begin{ddtacore}[title={Decision}]
\fontsize{5.15}{5.65}\selectfont
\begin{tabularx}{\linewidth}{L{40mm}Y}
\toprule
\textbf{ID} & \texttt{DEC-06} \\
\textbf{Lifecycle} & \texttt{current} \\
\textbf{Authority} & \texttt{CURRENT\_GOVERNED} \\
\textbf{Type} & \texttt{Decision} \\
\textbf{Title} & Accettazione comparativa dell'adattamento della classificazione di urgenza \\
\textbf{Parent MacroRequirement} & \texttt{MR-04} \\
\textbf{Decision kind} & \texttt{SELECTABLE} \\
\bottomrule
\end{tabularx}

\textbf{Context}\par
Un candidato di adattamento classificatorio non deve essere adottato soltanto perche' e' stato prodotto. La reference EfficientNet-B4 model version \texttt{1.0.0} riporta overall accuracy 81.74\%, Macro F1 0.8036, AUC-ROC 0.906, HIGH sensitivity 90.24\% e false LOW rate 7.93\%; la fonte riporta anche HIGH recall 81.71\% e 134/164 casi HIGH corretti, con relazione fra i due valori non chiarita. Serve una valutazione comparativa rispetto alla reference applicabile.

\textbf{Decision}\par
Un candidato di adattamento della classificazione di urgenza deve essere valutato comparativamente rispetto alla reference applicabile e puo' essere considerato qualificato per l'adozione soltanto quando soddisfa tutti i criteri comparativi governati applicabili.

\textbf{Consequences}\par
\texttt{FR-09} governa valutazione e qualificazione comparativa. La qualificazione non equivale a deployment automatico. \texttt{DEC-07} governa proprieta' qualitative aggiuntive dei criteri di accettazione.
\end{ddtacore}

\begin{ddtaquestions}[top=3pt,bottom=3pt,left=5pt,right=5pt]
\fontsize{4.65}{5.05}\selectfont
\begin{itemize}
  \item Review documentale: i risultati misurati della reference restano per ora nel perimetro FR della tesi; la collocazione naturale futura e' un artefatto dedicato di test/experiment result evidence collegato a requisito e codice di valutazione.
  \item La formulazione ``la fonte riporta'' e' residuo del reverse-documentation test e dovra' essere eliminata nella documentazione primaria autosufficiente.
  \item La relazione fra HIGH sensitivity 90.24\% e HIGH recall 81.71\% (134/164) resta \texttt{NOT SPECIFIED}; nessuna riconciliazione viene inferita.
\end{itemize}
\end{ddtaquestions}

\switchcolumn

{\sffamily\large\bfseries\color{DDTABlue}Base Analysis}\par
\vspace{0.35em}

\begin{ddtareferentidentified}[top=3pt,bottom=3pt,left=5pt,right=5pt]
\fontsize{4.55}{4.95}\selectfont
\texttt{--}
\end{ddtareferentidentified}

\begin{ddtareferentaccepted}[top=3pt,bottom=3pt,left=5pt,right=5pt]
\fontsize{4.35}{4.75}\selectfont
\BAAccepted{EfficientNet-B4}; \BAAccepted{CandidateUrgencyClassificationAdaptation} \BAFirstAccepted{}; \BAAccepted{ApplicableReferenceVersion} \BAFirstAccepted{}.
\end{ddtareferentaccepted}

\begin{ddtapropositionidentified}[top=3pt,bottom=3pt,left=5pt,right=5pt]
\fontsize{4.25}{4.65}\selectfont
\texttt{--}\par
\emph{La Decision governa il boundary di policy ``valutazione comparativa prima della qualificazione'', ma non stabilisce ancora una decomposizione BA completa fra experiment/evaluation result e qualification result. Nessun operatore viene forzato dalla sola Decision.}
\end{ddtapropositionidentified}

\begin{ddtapropositionaccepted}[top=3pt,bottom=3pt,left=5pt,right=5pt]
\fontsize{4.45}{4.85}\selectfont
\texttt{--}
\end{ddtapropositionaccepted}

\begin{ddtaopenrelations}[top=3pt,bottom=3pt,left=5pt,right=5pt]
\fontsize{3.82}{4.2}\selectfont
Governed facts: il candidato deve essere confrontato con una reference applicabile e la qualification richiede il soddisfacimento di tutti i criteri governati applicabili; qualification e deployment restano distinti. I valori di EfficientNet-B4 \texttt{1.0.0} sono measurement/reference evidence, non diventano automaticamente BAProposition ne' acceptance criteria. Macro F1 e AUC-ROC sono riportati come metriche della reference ma \texttt{DEC-06} non li rende, da sola, criteri di accettazione. Il binding verso esperimento, artefatto di output e codice che hanno prodotto tali valori resta aperto.
\end{ddtaopenrelations}

\end{paracol}
%</PAGE-CONTENT:32>
\clearpage
%<PAGE-DESC:33>MR-04 > DEC-06 > FR-09 - Qualificazione comparativa del candidato di adattamento classificatorio
%<PAGE-CONTENT:33>
\PageHashFooter{\PageMDfive{33}}
\subsubsection{FR-09 - Qualificazione comparativa del candidato di adattamento classificatorio}
{\sffamily\scriptsize\color{DDTAGray}Gerarchia: \texttt{MR-04} $\rightarrow$ \texttt{DEC-06} $\rightarrow$ \texttt{FR-09}}\par

\setlength{\columnsep}{7mm}
\setlength{\columnseprule}{0.35pt}
\columnratio{0.67}
\begin{paracol}{2}

{\sffamily\large\bfseries\color{DDTANavy}Documentazione DDTA}\par
\vspace{0.3em}

\begin{ddtacore}[title={FunctionalRequirement}]
\fontsize{5.0}{5.45}\selectfont
\begin{tabularx}{\linewidth}{L{39mm}Y}
\toprule
\textbf{ID} & \texttt{FR-09} \\
\textbf{Lifecycle} & \texttt{current} \\
\textbf{Authority} & \texttt{CURRENT\_GOVERNED} \\
\textbf{Type} & \texttt{FunctionalRequirement} \\
\textbf{Title} & Qualificazione comparativa del candidato di adattamento classificatorio \\
\textbf{Parent Decision} & \texttt{DEC-06} \\
\textbf{Owning MR (derived)} & \texttt{MR-04} \\
\textbf{SpecializedRequirement relation} & 0 active instances; \texttt{DEC-07} preserva il significato quality applicabile senza introdurre un subtype concreto \\
\bottomrule
\end{tabularx}

\textbf{Functional obligation}\par
Valutare comparativamente un candidato classificatorio e considerarlo qualificato per l'adozione solo se soddisfa tutti i criteri applicabili.

\textbf{Functional clauses / inherited \texttt{normativeClause [1..*]}}\par
\textbf{NC-09.1.} Quando un candidato di adattamento della classificazione di urgenza viene considerato per l'adozione, DermaTriage MUST valutarlo comparativamente rispetto alla versione di riferimento applicabile usando i criteri di accettazione governati.\par
\medskip
\textbf{NC-09.2.} DermaTriage MUST considerare il candidato qualificato per l'adozione soltanto se tutti i criteri comparativi governati applicabili risultano soddisfatti.

\textbf{SPO references (L2 support)}\par
\texttt{DermaTriage -- comparatively evaluate -->}\par
\texttt{CandidateUrgencyClassificationAdaptation}\par
\texttt{CandidateUrgencyClassificationAdaptation -- compared against -->}\par
\texttt{ApplicableReferenceVersion}\par
\texttt{SatisfiedAcceptanceCriteria -- qualify --> CandidateUrgencyClassificationAdaptation}
\end{ddtacore}

\begin{ddtaquestions}[top=3pt,bottom=3pt,left=5pt,right=5pt]
\fontsize{4.45}{4.85}\selectfont
\begin{itemize}
  \item Esperimento/test, artefatto di risultato e codice che producono le metriche candidate/reference non sono ancora governati come identita' concrete.
  \item Resta da stabilire se evaluation e qualification siano un'unica production con piu' result oppure due production in catena. Un eventuale \texttt{transfer} intermedio richiede source/destination e conveyance governati, non la sola data dependency.
\end{itemize}
\end{ddtaquestions}

\switchcolumn

{\sffamily\large\bfseries\color{DDTABlue}Base Analysis}\par
\vspace{0.3em}

\begin{ddtareferentidentified}[top=3pt,bottom=3pt,left=5pt,right=5pt]
\fontsize{4.05}{4.45}\selectfont
\BACandidate{CandidateEvaluationResult} \BANew{}; \BACandidate{ReferenceEvaluationResult} \BANew{}; \BACandidate{CandidateQualificationResult} \BANew{}; \BACandidate{ApplicableComparativeAcceptanceCriteria} \BANew{}.\par
\emph{Sono identita' analitiche candidate emerse dalla decomposizione; i relativi artefatti/result boundary non sono ancora documentati in modo sufficiente per l'accettazione.}
\end{ddtareferentidentified}

\begin{ddtareferentaccepted}[top=3pt,bottom=3pt,left=5pt,right=5pt]
\fontsize{4.35}{4.75}\selectfont
\BAAccepted{DermaTriage}; \BAAccepted{CandidateUrgencyClassificationAdaptation}; \BAAccepted{ApplicableReferenceVersion}.
\end{ddtareferentaccepted}

\begin{ddtapropositionidentified}[top=3pt,bottom=3pt,left=5pt,right=5pt]
\fontsize{3.72}{4.1}\selectfont
\textbf{Comparative evaluation pressure.} Il comportamento e' governato, ma il semanticOperatorKey esatto resta aperto: non viene forzato in \texttt{produce}, \texttt{observe} o \texttt{constrain}.\par
\medskip
\textbf{Qualification-rule pressure.} L'authority BA corrente ammette \texttt{decisionRule}; una candidate minima avrebbe actor \texttt{DermaTriage}, input candidato/reference, result \texttt{CandidateQualificationResult} e branch \texttt{IF all applicable criteria satisfied THEN ... = QUALIFIED}. Non viene accettata ora perche' la documentazione non ha ancora chiuso il boundary fra evaluation result e qualification result, e R7 testa la relocation di \texttt{decisionRule} sotto un owner \texttt{produce} che qui non e' ancora stabilito.
\end{ddtapropositionidentified}

\begin{ddtapropositionaccepted}[top=3pt,bottom=3pt,left=5pt,right=5pt]
\fontsize{4.45}{4.85}\selectfont
\texttt{--}
\end{ddtapropositionaccepted}

\begin{ddtaopenrelations}[top=3pt,bottom=3pt,left=5pt,right=5pt]
\fontsize{3.52}{3.9}\selectfont
Due decomposizioni restano materialmente possibili. \textbf{A)} una sola \texttt{produce} rende disponibili \texttt{CandidateEvaluationResult} e \texttt{CandidateQualificationResult}; \textbf{B)} una prima \texttt{produce} rende disponibile \texttt{CandidateEvaluationResult} e una seconda production di qualification lo usa, con \texttt{ReferenceEvaluationResult} e i criteri applicabili, per rendere disponibile \texttt{CandidateQualificationResult}. \texttt{QUALIFIED} resta per ora un valore locale candidato del result, non un BAReferent. Nessun \texttt{ELSE} viene inferito. La sola catena result-as-input non autorizza un \texttt{transfer}; esso diventerebbe rappresentabile soltanto se documentazione/code evidence governassero endpoint distinti e una conveyance. La future review deve cercare separatamente codice che calcola metriche e codice che applica la policy di qualification.
\end{ddtaopenrelations}

\end{paracol}
%</PAGE-CONTENT:33>
\clearpage
%<PAGE-DESC:34>MR-04 > DEC-07 - Prioritizzazione asimmetrica delle metriche di qualita
%<PAGE-CONTENT:34>
\PageHashFooter{\PageMDfive{34}}
\subsection{DEC-07 - Prioritizzazione asimmetrica delle metriche di qualita'}
{\sffamily\scriptsize\color{DDTAGray}Gerarchia: \texttt{MR-04} $\rightarrow$ \texttt{DEC-07}}\par

\setlength{\columnsep}{7mm}
\setlength{\columnseprule}{0.35pt}
\columnratio{0.67}
\begin{paracol}{2}

{\sffamily\large\bfseries\color{DDTANavy}Documentazione DDTA}\par
\vspace{0.35em}

\begin{ddtacore}[title={Decision}]
\fontsize{5.45}{5.95}\selectfont
\begin{tabularx}{\linewidth}{L{40mm}Y}
\toprule
\textbf{ID} & \texttt{DEC-07} \\
\textbf{Lifecycle} & \texttt{current} \\
\textbf{Authority} & \texttt{CURRENT\_GOVERNED} \\
\textbf{Type} & \texttt{Decision} \\
\textbf{Title} & Prioritizzazione asimmetrica delle metriche di qualita' \\
\textbf{Parent MacroRequirement} & \texttt{MR-04} \\
\textbf{Decision kind} & \texttt{SELECTABLE} \\
\bottomrule
\end{tabularx}

\textbf{Context}\par
La policy di accettazione non tratta tutte le metriche come intercambiabili. La documentazione attribuisce priorita' alla non-degradazione di sensitivity e false-low, consentendo per l'accuracy soltanto una degradazione limitata.

\textbf{Decision}\par
Nel confronto con la reference applicabile, sensitivity e false-low performance non devono peggiorare, mentre l'overall accuracy puo' degradare al massimo del 5\%.

\textbf{Consequences}\par
La qualificazione comparativa tratta in modo asimmetrico le metriche: sensitivity e false-low performance non devono peggiorare, mentre l'overall accuracy puo' degradare soltanto entro la tolleranza applicabile.
\end{ddtacore}

\begin{ddtaquestions}[top=3pt,bottom=3pt,left=5pt,right=5pt]
\fontsize{4.55}{4.95}\selectfont
\begin{itemize}
  \item Il literal 5\% non specifica se la degradazione sia percentuale relativa o differenza in punti percentuali: non viene normalizzato in una formula ulteriore.
  \item \texttt{DEC-06} registra sia HIGH sensitivity 90.24\% sia HIGH recall 81.71\% (134/164), con relazione non chiarita; quale misura sia la sensitivity operativa del criterio resta da chiarire.
  \item Macro F1 e AUC-ROC non vengono promossi a criteri di accettazione: \texttt{DEC-07} non lo stabilisce.
\end{itemize}
\end{ddtaquestions}

\switchcolumn

{\sffamily\large\bfseries\color{DDTABlue}Base Analysis}\par
\vspace{0.35em}

\begin{ddtareferentidentified}[top=3pt,bottom=3pt,left=5pt,right=5pt]
\fontsize{4.5}{4.9}\selectfont
\texttt{--}
\end{ddtareferentidentified}

\begin{ddtareferentaccepted}[top=3pt,bottom=3pt,left=5pt,right=5pt]
\fontsize{4.2}{4.6}\selectfont
\BAAccepted{CandidateUrgencyClassificationAdaptation}; \BAAccepted{ApplicableReferenceVersion}; \BAAccepted{ApplicableComparativeAcceptanceCriteria} \BAFirstAccepted{}.
\end{ddtareferentaccepted}

\begin{ddtapropositionidentified}[top=3pt,bottom=3pt,left=5pt,right=5pt]
\fontsize{3.9}{4.28}\selectfont
\textbf{Candidate comparison structure, non accettata:}\par
\texttt{sensitivity: candidate no worse than reference}\par
\texttt{false-low: candidate no worse than reference}\par
\texttt{overall accuracy: degradation <= 5\%}\par
\emph{La struttura preserva il significato leggibile ma non viene elevata a syntax BA: il vocabulary ordinato/scalare necessario non e' chiuso e la semantica esatta del 5\% resta documentariamente ambigua.}
\end{ddtapropositionidentified}

\begin{ddtapropositionaccepted}[top=3pt,bottom=3pt,left=5pt,right=5pt]
\fontsize{4.45}{4.85}\selectfont
\texttt{--}
\end{ddtapropositionaccepted}

\begin{ddtaopenrelations}[top=3pt,bottom=3pt,left=5pt,right=5pt]
\fontsize{3.82}{4.2}\selectfont
\texttt{DEC-07} rende identificabile e riusabile il set \texttt{ApplicableComparativeAcceptanceCriteria} richiamato da \texttt{FR-09}, ma non introduce una nuova funzione. I tre criteri sono path-local alla qualification classificatoria e devono poter alimentare la futura decision structure senza inventare operatori di confronto. La documentazione corrente richiede almeno ordered/scalar comparison pressure (non-degradation e maximum degradation); finche' comparator semantics e measurement basis non sono governati in modo sufficiente, nessuna BAProposition \texttt{constrain} o \texttt{decisionRule} viene accettata da questa Decision isolatamente.
\end{ddtaopenrelations}

\end{paracol}
%</PAGE-CONTENT:34>
\clearpage

%<PAGE-DESC:35>MR-04 > DEC-08 - Reversibilita dell'adattamento della classificazione in caso di degrado
%<PAGE-CONTENT:35>
\PageHashFooter{\PageMDfive{35}}
\subsection{DEC-08 - Reversibilita' dell'adattamento della classificazione in caso di degrado}
{\sffamily\scriptsize\color{DDTAGray}Gerarchia: \texttt{MR-04} $\rightarrow$ \texttt{DEC-08}}\par

\setlength{\columnsep}{7mm}
\setlength{\columnseprule}{0.35pt}
\columnratio{0.67}
\begin{paracol}{2}

{\sffamily\large\bfseries\color{DDTANavy}Documentazione DDTA}\par
\vspace{0.3em}

\begin{ddtacore}[title={Decision}]
\fontsize{5.2}{5.7}\selectfont
\begin{tabularx}{\linewidth}{L{40mm}Y}
\toprule
\textbf{ID} & \texttt{DEC-08} \\
\textbf{Lifecycle} & \texttt{current} \\
\textbf{Authority} & \texttt{CURRENT\_GOVERNED} \\
\textbf{Type} & \texttt{Decision} \\
\textbf{Title} & Reversibilita' dell'adattamento della classificazione in caso di degrado \\
\textbf{Parent MacroRequirement} & \texttt{MR-04} \\
\textbf{Decision kind} & \texttt{SELECTABLE} \\
\bottomrule
\end{tabularx}

\textbf{Context}\par
Un adattamento gia' adottato puo' manifestare degrado materiale dopo l'adozione. La capacita' di recuperare una versione o stato precedente deve restare distinta dalla qualificazione pre-adoption.

\textbf{Decision}\par
Un adattamento della classificazione di urgenza gia' adottato che manifesta un degrado di accuracy superiore al 5\% rispetto alla reference post-adoption applicabile deve poter essere revocato ripristinando una versione o stato precedente accettabile.

\textbf{Consequences}\par
\texttt{FR-10} governa la revoca/ripristino. Il requisito non stabilisce che il rollback sia automatico. Il threshold post-adoption \texttt{R} resta semanticamente distinto dalla tolleranza pre-adoption \texttt{T} anche quando i literal correnti coincidono.
\end{ddtacore}

\begin{ddtaquestions}[top=3pt,bottom=3pt,left=5pt,right=5pt]
\fontsize{4.35}{4.75}\selectfont
\begin{itemize}
  \item La clausola governa un solo ramo applicabile quando \texttt{accuracy degradation > R}; nessun ramo \texttt{ELSE} e nessuna proposition ``do nothing'' vengono introdotti.
  \item Il literal 5\% non chiarisce percentuale relativa vs punti percentuali; la measurement basis resta \texttt{NOT SPECIFIED}.
  \item \texttt{R} post-adoption e \texttt{T} pre-adoption restano significati distinti anche se oggi condividono il literal 5\%.
\end{itemize}
\end{ddtaquestions}

\switchcolumn

{\sffamily\large\bfseries\color{DDTABlue}Base Analysis}\par
\vspace{0.3em}

\begin{ddtareferentidentified}[top=3pt,bottom=3pt,left=5pt,right=5pt]
\fontsize{4.15}{4.55}\selectfont
\BACandidate{PostAdoptionEvaluationResult} \BANew{}.\par
\emph{Candidate decomposition: la source governa la condizione di degrado ma non identifica ancora l'esperimento, l'actor o un result artifact concreto che la renda disponibile.}
\end{ddtareferentidentified}

\begin{ddtareferentaccepted}[top=3pt,bottom=3pt,left=5pt,right=5pt]
\fontsize{4.15}{4.55}\selectfont
\BAAccepted{ClassifierAdaptationPath}; \BAAccepted{ApplicablePostAdoptionReference} \BAFirstAccepted{}; \BAAccepted{PreviousAcceptableVersionOrState} \BAFirstAccepted{}.
\end{ddtareferentaccepted}

\begin{ddtapropositionidentified}[top=3pt,bottom=3pt,left=5pt,right=5pt]
\fontsize{3.9}{4.3}\selectfont
\textbf{Post-adoption result pressure.} Una futura \texttt{produce} puo' essere appropriata se documentazione/code evidence governeranno chi calcola la valutazione post-adoption e quale result rende disponibile; oggi actor e result boundary non sono sufficientemente specificati.\par
\medskip
\textbf{Restore applicability.} \texttt{degradation > R} limita l'applicabilita' del recovery path, ma non esiste ancora una proposition owner corrente alla quale applicare \texttt{condition} senza inventare la semantica del restore.
\end{ddtapropositionidentified}

\begin{ddtapropositionaccepted}[top=3pt,bottom=3pt,left=5pt,right=5pt]
\fontsize{4.35}{4.75}\selectfont
\texttt{--}
\end{ddtapropositionaccepted}

\begin{ddtaopenrelations}[top=3pt,bottom=3pt,left=5pt,right=5pt]
\fontsize{3.72}{4.08}\selectfont
Nessuna \texttt{transition} viene accettata: la source governa il ripristino di una versione/stato precedente ma non stabilisce una stessa identity che attraversi stati \texttt{Adopted -> Revoked}. Nessun rollback automatico viene inferito. La lettura operativa e': un modello correntemente attivo puo' risultare degradato; quando il degrado supera \texttt{R}, il recovery/restore deve essere supportato. I simboli di lavoro \texttt{ME} (modello corrente) e \texttt{MOK} (precedente accettabile) non diventano nomi BA canonici. La meccanica con cui \texttt{MOK} torna attivo resta demandata a \texttt{FR-10} e all'evidence implementativa.
\end{ddtaopenrelations}

\end{paracol}
%</PAGE-CONTENT:35>
\clearpage

%<PAGE-DESC:36>MR-04 > DEC-08 > FR-10 - Revoca di un adattamento classificatorio degradante
%<PAGE-CONTENT:36>
\PageHashFooter{\PageMDfive{36}}
\subsubsection{FR-10 - Revoca di un adattamento classificatorio degradante}
{\sffamily\scriptsize\color{DDTAGray}Gerarchia: \texttt{MR-04} $\rightarrow$ \texttt{DEC-08} $\rightarrow$ \texttt{FR-10}}\par

\setlength{\columnsep}{7mm}
\setlength{\columnseprule}{0.35pt}
\columnratio{0.67}
\begin{paracol}{2}

{\sffamily\large\bfseries\color{DDTANavy}Documentazione DDTA}\par
\vspace{0.3em}

\begin{ddtacore}[title={FunctionalRequirement}]
\fontsize{4.85}{5.3}\selectfont
\begin{tabularx}{\linewidth}{L{39mm}Y}
\toprule
\textbf{ID} & \texttt{FR-10} \\
\textbf{Lifecycle} & \texttt{current} \\
\textbf{Authority} & \texttt{CURRENT\_GOVERNED} \\
\textbf{Type} & \texttt{FunctionalRequirement} \\
\textbf{Title} & Revoca di un adattamento classificatorio degradante \\
\textbf{Parent Decision} & \texttt{DEC-08} \\
\textbf{Owning MR (derived)} & \texttt{MR-04} \\
\textbf{SpecializedRequirement relation} & 0 active instances \\
\bottomrule
\end{tabularx}

\textbf{Functional obligation}\par
Supportare la revoca di un adattamento classificatorio gia' adottato quando manifesta degrado materiale, ripristinando una versione o stato precedente accettabile.

\textbf{Functional clauses / inherited \texttt{normativeClause [1..*]}}\par
Per un adattamento della classificazione di urgenza gia' adottato che manifesti un degrado di accuracy superiore al 5\% rispetto alla reference post-adoption applicabile, DermaTriage MUST supportarne la revoca mediante il ripristino di una versione o stato precedente accettabile; le versioni ripristinabili sono mantenute in \nolinkurl{models/versions/} con tracking in \nolinkurl{db/model_versions.json} e modello attivo in \nolinkurl{models/efficientnet_b4.pth}.

\textbf{SPO references (L2 support)}\par
\texttt{MateriallyDegradingAdoptedAdaptation -- triggers support for --> Revocation}\par
\texttt{DermaTriage -- restore --> PreviousAcceptableVersionOrState}
\end{ddtacore}

\begin{ddtaquestions}[top=3pt,bottom=3pt,left=5pt,right=5pt]
\fontsize{4.02}{4.42}\selectfont
\begin{itemize}
  \item \texttt{models/versions/}, \texttt{db/model\_versions.json} e \texttt{models/efficientnet\_b4.pth} sono binding tecnici governati da preservare; non vengono sostituiti da nomi BA generici.
  \item Il restore puo' essere copy/move/overwrite/load/reference switch o altro: la meccanica di conveyance non e' specificata. Il destino del modello precedentemente attivo (conservato, archiviato, cancellato) e' \texttt{NOT SPECIFIED}.
  \item TODO cumulativo persistence/storage: riesaminare insieme \texttt{FR-03} (record/persistence), worklist \texttt{T15} (LocalDiagnosisStore) e \texttt{FR-10}, distinguendo \texttt{transfer}, write/persist riuscito e associazione at-rest. Non si riapre ora \texttt{storedIn}.
\end{itemize}
\end{ddtaquestions}

\switchcolumn

{\sffamily\large\bfseries\color{DDTABlue}Base Analysis}\par
\vspace{0.3em}

\begin{ddtareferentidentified}[top=3pt,bottom=3pt,left=5pt,right=5pt]
\fontsize{4.0}{4.38}\selectfont
\BACandidate{PostAdoptionEvaluationResult};\par
\emph{Resta candidate finche' non sono governati producer, result artifact e measurement boundary. I tre path tecnici sono preservati come implementation bindings e non vengono promossi automaticamente a BAReferent.}
\end{ddtareferentidentified}

\begin{ddtareferentaccepted}[top=3pt,bottom=3pt,left=5pt,right=5pt]
\fontsize{4.2}{4.6}\selectfont
\BAAccepted{DermaTriage}; \BAAccepted{ApplicablePostAdoptionReference}; \BAAccepted{PreviousAcceptableVersionOrState}.
\end{ddtareferentaccepted}

\begin{ddtapropositionidentified}[top=3pt,bottom=3pt,left=5pt,right=5pt]
\fontsize{3.65}{4.0}\selectfont
\textbf{Restore/transfer pressure, non accettata.}\par
Una candidate \texttt{transfer} \texttt{models/versions/ -> active model location : MOK} richiederebbe evidence che governi una vera conveyance; la sola presenza di uno store di versioni e di un path del modello attivo non basta.\par
\medskip
\textbf{Persistence pressure, deferita.}\par
\texttt{maintained in} governa un fatto at-rest, ma il costrutto \texttt{storedIn} non e' corrente. Inoltre \texttt{transfer-to-store != successful write/persist != at-rest association}: il caso viene registrato nel TODO cumulativo con \texttt{FR-03}/\texttt{T15}.
\end{ddtapropositionidentified}

\begin{ddtapropositionaccepted}[top=3pt,bottom=3pt,left=5pt,right=5pt]
\fontsize{4.3}{4.7}\selectfont
\texttt{--}
\end{ddtapropositionaccepted}

\begin{ddtaopenrelations}[top=3pt,bottom=3pt,left=5pt,right=5pt]
\fontsize{3.62}{3.96}\selectfont
La BA puo' essere scritta gia' a questo checkpoint anche senza una nuova BAProposition: accetta le identita' semantiche governate, preserva i binding tecnici, registra candidate decomposition e delimita le inferenze vietate. Il solo ramo governato e' \texttt{degradation > R -> recovery/restore support applicable}; l'assenza della condizione non genera un ramo \texttt{ELSE}. Nessuna \texttt{transition} viene introdotta e nessuna fine di \texttt{ME} viene inventata. La futura evidence deve cercare separatamente: (1) producer/result della valutazione post-adoption, (2) decision/avvio effettivo del rollback se esiste, (3) meccanica di restore fra version store e modello attivo, (4) persistence/write semantics.
\end{ddtaopenrelations}

\end{paracol}
%</PAGE-CONTENT:36>
\clearpage

%<PAGE-DESC:37>MR-04 > DEC-09 - Uso di evidence recente e limitata nel percorso di evoluzione del prompt
%<PAGE-CONTENT:37>
\PageHashFooter{\PageMDfive{37}}
\subsection{DEC-09 - Uso di evidence recente e limitata nel percorso di evoluzione del prompt}
{\sffamily\scriptsize\color{DDTAGray}Gerarchia: \texttt{MR-04} $\rightarrow$ \texttt{DEC-09}}\par

\setlength{\columnsep}{7mm}
\setlength{\columnseprule}{0.35pt}
\columnratio{0.67}
\begin{paracol}{2}

{\sffamily\large\bfseries\color{DDTANavy}Documentazione DDTA}\par
\vspace{0.3em}

\begin{ddtacore}[title={Decision}]
\fontsize{5.05}{5.55}\selectfont
\begin{tabularx}{\linewidth}{L{40mm}Y}
\toprule
\textbf{ID} & \texttt{DEC-09} \\
\textbf{Lifecycle} & \texttt{current} \\
\textbf{Authority} & \texttt{CURRENT\_GOVERNED} \\
\textbf{Type} & \texttt{Decision} \\
\textbf{Title} & Uso di evidence recente e limitata nel percorso di evoluzione del prompt \\
\textbf{Parent MacroRequirement} & \texttt{MR-04} \\
\textbf{Decision kind} & \texttt{SELECTABLE} \\
\bottomrule
\end{tabularx}

\textbf{Context}\par
L'evoluzione del prompt non utilizza una storia illimitata delle correzioni cliniche. La fonte governa l'uso di evidence recente e limitata per costruire il set del ciclo corrente.

\textbf{Decision}\par
Per l'evoluzione del prompt, DermaTriage usa una finestra scorrevole e limitata delle correzioni cliniche pertinenti piu' recenti.

\textbf{Consequences}\par
\texttt{FR-06} governa la selezione dell'evidence set. La finestra contiene i 20 esempi pertinenti piu' recenti; ordering, deduplication, underfill e riuso tra cicli restano non sufficientemente specificati.
\end{ddtacore}

\begin{ddtaquestions}[top=3pt,bottom=3pt,left=5pt,right=5pt]
\fontsize{4.15}{4.55}\selectfont
\begin{itemize}
  \item ``Piu' recenti'' governa la membership della finestra, non l'ordinamento con cui il result set viene restituito; nessun \texttt{orderingBasis} viene quindi assunto.
  \item Non sono specificati deduplication, underfill, reuse/overlap tra cicli ne' il criterio completo di pertinenza.
  \item La Decision non governa storage, query/read o transfer delle correzioni: il comportamento resta analizzato come black box fino a evidence tecnica successiva.
\end{itemize}
\end{ddtaquestions}

\switchcolumn

{\sffamily\large\bfseries\color{DDTABlue}Base Analysis}\par
\vspace{0.3em}

\begin{ddtareferentidentified}[top=3pt,bottom=3pt,left=5pt,right=5pt]
\fontsize{4.3}{4.7}\selectfont
\BACandidate{PromptEvolutionEvidenceSet} \BANew{}; \BACandidate{PromptEvolutionCycle}.\par
\emph{\texttt{PromptEvolutionCycle} era gia' candidato in \texttt{FR-04}; qui riappare come scope del set del ciclo corrente.}
\end{ddtareferentidentified}

\begin{ddtareferentaccepted}[top=3pt,bottom=3pt,left=5pt,right=5pt]
\fontsize{4.35}{4.75}\selectfont
\BAAccepted{DermaTriage}; \BAAccepted{PromptEvolutionPath}; \BAAccepted{ClinicalReviewResult}.
\end{ddtareferentaccepted}

\begin{ddtapropositionidentified}[top=3pt,bottom=3pt,left=5pt,right=5pt]
\fontsize{4.0}{4.38}\selectfont
\textbf{Selection pressure.}\par
La source governa una selezione recent-N da una popolazione piu' ampia di review/correzioni, ma \texttt{selection} non e' un operatore corrente R7. La Decision non viene forzata in \texttt{observe}, \texttt{transfer}, \texttt{constrain} o \texttt{decisionRule}; \texttt{FR-06} fornisce l'owner produttivo da analizzare.
\end{ddtapropositionidentified}

\begin{ddtapropositionaccepted}[top=3pt,bottom=3pt,left=5pt,right=5pt]
\fontsize{4.35}{4.75}\selectfont
\texttt{--}
\end{ddtapropositionaccepted}

\begin{ddtaopenrelations}[top=3pt,bottom=3pt,left=5pt,right=5pt]
\fontsize{3.85}{4.22}\selectfont
La struttura candidate da preservare e' una selezione locale alla futura produzione dell'evidence set: filtra le correzioni pertinenti, usa la recency per scegliere i membri e applica bound 20. Non viene creato un BAReferent \texttt{RecentPertinentClinicalCorrections}: il significato puo' essere mantenuto come membership selezionata di \texttt{ClinicalReviewResult}. La forma esatta della local structure resta da ricostruire metodologicamente; questo caso la registra come pressione DermaTriage, non come vocabulary ammesso.
\end{ddtaopenrelations}

\end{paracol}
%</PAGE-CONTENT:37>
\clearpage

%<PAGE-DESC:38>MR-04 > DEC-09 > FR-06 - Selezione dell'evidence recente per l'evoluzione del prompt
%<PAGE-CONTENT:38>
\PageHashFooter{\PageMDfive{38}}
\subsubsection{FR-06 - Selezione dell'evidence recente per l'evoluzione del prompt}
{\sffamily\scriptsize\color{DDTAGray}Gerarchia: \texttt{MR-04} $\rightarrow$ \texttt{DEC-09} $\rightarrow$ \texttt{FR-06}}\par

\setlength{\columnsep}{7mm}
\setlength{\columnseprule}{0.35pt}
\columnratio{0.67}
\begin{paracol}{2}

{\sffamily\large\bfseries\color{DDTANavy}Documentazione DDTA}\par
\vspace{0.3em}

\begin{ddtacore}[title={FunctionalRequirement}]
\fontsize{5.0}{5.5}\selectfont
\begin{tabularx}{\linewidth}{L{39mm}Y}
\toprule
\textbf{ID} & \texttt{FR-06} \\
\textbf{Lifecycle} & \texttt{current} \\
\textbf{Authority} & \texttt{CURRENT\_GOVERNED} \\
\textbf{Type} & \texttt{FunctionalRequirement} \\
\textbf{Title} & Selezione dell'evidence recente per l'evoluzione del prompt \\
\textbf{Parent Decision} & \texttt{DEC-09} \\
\textbf{Owning MR (derived)} & \texttt{MR-04} \\
\textbf{SpecializedRequirement relation} & 0 active instances \\
\bottomrule
\end{tabularx}

\textbf{Functional obligation}\par
Costruire l'evidence set del ciclo prompt usando una finestra scorrevole limitata delle correzioni cliniche pertinenti piu' recenti.

\textbf{Functional clauses / inherited \texttt{normativeClause [1..*]}}\par
Per ciascun ciclo di evoluzione del prompt, DermaTriage MUST costruire l'evidence set utilizzando una finestra scorrevole dei 20 esempi di correzione clinica pertinenti piu' recenti.

\textbf{SPO references (L2 support)}\par
\texttt{DermaTriage -- select recent bounded --> PromptEvolutionEvidenceSet}\par
\texttt{RecentPertinentClinicalCorrections -- populate --> PromptEvolutionEvidenceSet}
\end{ddtacore}

\begin{ddtaquestions}[top=3pt,bottom=3pt,left=5pt,right=5pt]
\fontsize{3.95}{4.33}\selectfont
\begin{itemize}
  \item \texttt{FR-04} attiva il ciclo quando l'evidence accumulata raggiunge 10 elementi; \texttt{FR-06} richiede una finestra di 20. La relazione fra conteggio di attivazione e popolazione disponibile alla selezione non e' specificata; se sono disponibili meno di 20 elementi, l'underfill resta \texttt{NOT SPECIFIED}.
  \item Il requisito governa la produzione dell'evidence set, non ancora la produzione di una nuova \texttt{PromptVersion}: il comportamento downstream di evoluzione/update del prompt sara' verificato nel relativo blocco documentale.
  \item Nessuna read/query o conveyance da uno storage viene inferita: se futura evidence governera' un boundary DB $\rightarrow$ funzione, quella semantica verra' analizzata separatamente.
\end{itemize}
\end{ddtaquestions}

\switchcolumn

{\sffamily\large\bfseries\color{DDTABlue}Base Analysis}\par
\vspace{0.3em}

\begin{ddtareferentidentified}[top=3pt,bottom=3pt,left=5pt,right=5pt]
\fontsize{4.15}{4.55}\selectfont
\BACandidate{PromptEvolutionCycle}; \BACandidate{PromptEvolutionEvidenceSet}.
\end{ddtareferentidentified}

\begin{ddtareferentaccepted}[top=3pt,bottom=3pt,left=5pt,right=5pt]
\fontsize{4.12}{4.52}\selectfont
\BAAccepted{DermaTriage}; \BAAccepted{PromptEvolutionPath}; \BAAccepted{ClinicalReviewResult}; \BAAccepted{PromptEvolutionEvidenceSet} \BAFirstAccepted{}.
\end{ddtareferentaccepted}

\begin{ddtapropositionidentified}[top=3pt,bottom=3pt,left=5pt,right=5pt]
\fontsize{3.72}{4.06}\selectfont
\texttt{BAP-FR06-01}\par
\texttt{semanticOperatorKey = produce}\par
\hspace*{1em}\texttt{actor -> DermaTriage}\par
\hspace*{1em}\texttt{input -> ClinicalReviewResult [0..*]}\par
\hspace*{1em}\texttt{result -> PromptEvolutionEvidenceSet}\par
\texttt{polarity = affirmative}\par
\medskip
\emph{Semantica identificata ma non consolidata nel vocabulary BA corrente: dalla popolazione disponibile di \texttt{ClinicalReviewResult}, il requisito seleziona le correzioni pertinenti piu' recenti, con bound 20, per costituire \texttt{PromptEvolutionEvidenceSet}. Nessuna struttura formale aggiuntiva viene accettata in questa pagina.}
\end{ddtapropositionidentified}

\begin{ddtapropositionaccepted}[top=3pt,bottom=3pt,left=5pt,right=5pt]
\fontsize{3.92}{4.28}\selectfont
\textbf{BAP-FR06-01} come proposition \texttt{produce}.\par
\emph{Acceptance boundary: l'input rappresenta la popolazione disponibile di \texttt{ClinicalReviewResult}; non viene introdotto un BAReferent artificiale \texttt{ClinicalReviewCollection}. L'accettazione non include una struttura formale per determinare la membership del result. ``Piu' recenti'' governa quali elementi appartengono all'evidence set e non prova ordering dell'output.}
\end{ddtapropositionaccepted}

\begin{ddtaopenrelations}[top=3pt,bottom=3pt,left=5pt,right=5pt]
\fontsize{3.62}{3.98}\selectfont
\texttt{PromptEvolutionCycle} resta candidate: il riuso di ``per ciascun ciclo'' tra \texttt{FR-04}, \texttt{DEC-09} e \texttt{FR-06} fornisce pressione di identity/scope, ma non governa ancora lifecycle, partecipanti o altri criteri sufficienti per promuoverlo come BAReferent autonomo. \texttt{PromptEvolutionEvidenceSet} e' invece il result funzionale autonomo di \texttt{FR-06}. La semantica di selezione resta un finding identificato ma non consolidato nella BA corrente: correzione e pertinenza delimitano la membership, la recency determina la scelta dei piu' recenti e il bound governato e' 20; ordering dell'output, deduplication, underfill e riuso tra cicli restano non specificati. Nessuna \texttt{observe/read} viene aggiunta: DermaTriage e' qui la black box che usa un input ampio e produce il result governato. Eventuali storage/read/transfer tecnici restano fatti separati da introdurre soltanto se governati.
\end{ddtaopenrelations}

\end{paracol}
%</PAGE-CONTENT:38>
\clearpage

% Preserve standalone FR-08 heading counter while DEC-04/DEC-05/DEC-06/DEC-07 branches are inserted ahead of it.
\setcounter{subsection}{0}
\setcounter{subsubsection}{0}
%<PAGE-DESC:39>MR-04 > DEC-11 > FR-08 - Derivazione del target di supervisione dalla priorita P-scale corretta
%<PAGE-CONTENT:39>
\PageHashFooter{\PageMDfive{39}}
\subsubsection{FR-08 - Derivazione del target di supervisione dalla priorita P-scale corretta}
{\sffamily\scriptsize\color{DDTAGray}Gerarchia: \texttt{MR-04} $\rightarrow$ \texttt{DEC-11} $\rightarrow$ \texttt{FR-08}}\par

\setlength{\columnsep}{7mm}
\setlength{\columnseprule}{0.35pt}
\columnratio{0.67}
\begin{paracol}{2}

{\sffamily\large\bfseries\color{DDTANavy}Documentazione DDTA}\par
\vspace{0.3em}
\small

\begin{ddtacore}[title={Decision + FunctionalRequirement}]
\textbf{DEC-11 - Derivazione della supervision classificatoria dalla priorita clinicamente corretta}\par
\textbf{Context}\par
L'adattamento della classificazione di urgenza necessita di un target di supervisione. La priorita P-scale clinicamente corretta e' la sorgente semantica del target \texttt{HIGH / MEDIUM / LOW}.\par

\textbf{Decision}\par
Per i casi di adattamento classificatorio con priorita P-scale clinicamente corretta, DermaTriage deriva il target secondo il mapping governato \texttt{P1/P2 -> HIGH}, \texttt{P3 -> MEDIUM}, \texttt{P4 -> LOW}.\par

\medskip
\begin{tabularx}{\linewidth}{L{38mm}Y}
\toprule
\textbf{ID} & \texttt{FR-08} \\
\textbf{Lifecycle} & \texttt{current} \\
\textbf{Authority} & \texttt{CURRENT\_GOVERNED} \\
\textbf{Type} & \texttt{FunctionalRequirement} \\
\textbf{Parent Decision} & \texttt{DEC-11} \\
\textbf{Owning MR (derived)} & \texttt{MR-04} \\
\bottomrule
\end{tabularx}

\textbf{Functional obligation}\par
Derivare il target di supervisione classificatoria dalla priorita P-scale corretta dal professionista sanitario.\par

\textbf{Functional clauses / inherited \texttt{normativeClause [1..*]}}\par
\begin{tabularx}{\linewidth}{L{52mm}Y}
\toprule
\textbf{Priorita clinicamente corretta} & \textbf{Target richiesto} \\
\midrule
\texttt{P1} & DermaTriage MUST derivare \texttt{HIGH}. \\
\texttt{P2} & DermaTriage MUST derivare \texttt{HIGH}. \\
\texttt{P3} & DermaTriage MUST derivare \texttt{MEDIUM}. \\
\texttt{P4} & DermaTriage MUST derivare \texttt{LOW}. \\
\bottomrule
\end{tabularx}

\textbf{SPO references}\par
\texttt{ClinicallyCorrectedOperationalTriagePriority -- maps to --> ClassifierUrgencyTarget}\par
\texttt{DermaTriage -- derive --> ClassifierUrgencyTarget}
\end{ddtacore}

\begin{ddtaquestions}
\small
\begin{itemize}
  \item Il comportamento per priorita corretta assente, invalida, fuori dominio o conflittuale non e' specificato da FR-08.
  \item Nessuna relazione FR-07 $\rightarrow$ FR-08 viene inferita dal solo contesto condiviso di adattamento.
\end{itemize}
\end{ddtaquestions}

\switchcolumn

{\sffamily\large\bfseries\color{DDTABlue}Base Analysis}\par
\vspace{0.4em}

\begin{ddtareferentidentified}[top=3pt,bottom=3pt,left=5pt,right=5pt]
\fontsize{5.0}{5.4}\selectfont
\texttt{--}\par
\end{ddtareferentidentified}

\begin{ddtareferentaccepted}[top=3pt,bottom=3pt,left=5pt,right=5pt]
\fontsize{5.0}{5.4}\selectfont
\BAAccepted{DermaTriage}; \BAAccepted{ClinicallyCorrectedOperationalTriagePriority} \BAFirstAccepted{}; \BAAccepted{ClassifierUrgencyTarget} \BAFirstAccepted{}.
\end{ddtareferentaccepted}

\begin{ddtapropositionidentified}[top=3pt,bottom=3pt,left=5pt,right=5pt]
\fontsize{4.55}{4.95}\selectfont
\texttt{--}\par
\emph{Nessuna BAProposition separata per il mapping: la candidate \texttt{operatorStructure.decisionRule} e' locale alla proposition \texttt{produce} accettata sotto.}
\end{ddtapropositionidentified}

\begin{ddtapropositionaccepted}[top=3pt,bottom=3pt,left=5pt,right=5pt]
\fontsize{4.05}{4.42}\selectfont
\texttt{BAP-FR08-01-01}\par
\texttt{semanticOperatorKey = produce}\par
\hspace*{1em}\texttt{actor -> DermaTriage}\par
\hspace*{1em}\texttt{input -> ClinicallyCorrectedOperationalTriagePriority}\par
\hspace*{1em}\texttt{result -> ClassifierUrgencyTarget}\par
\hspace*{1em}\texttt{operatorStructure [CANDIDATE R7 / NOT ADMITTED]}\par
\hspace*{2em}\texttt{decisionRule}\par
\hspace*{3em}\texttt{branch: ClinicallyCorrectedOperationalTriagePriority eq P1}\par
\hspace*{4em}\texttt{resultAssignment -> ClassifierUrgencyTarget = HIGH}\par
\hspace*{3em}\texttt{branch: ClinicallyCorrectedOperationalTriagePriority eq P2}\par
\hspace*{4em}\texttt{resultAssignment -> ClassifierUrgencyTarget = HIGH}\par
\hspace*{3em}\texttt{branch: ClinicallyCorrectedOperationalTriagePriority eq P3}\par
\hspace*{4em}\texttt{resultAssignment -> ClassifierUrgencyTarget = MEDIUM}\par
\hspace*{3em}\texttt{branch: ClinicallyCorrectedOperationalTriagePriority eq P4}\par
\hspace*{4em}\texttt{resultAssignment -> ClassifierUrgencyTarget = LOW}\par
\texttt{polarity = affirmative}\par
\emph{Acceptance boundary: \texttt{BAP-FR08-01-01} e' accettata come proposition \texttt{produce}. La struttura annidata e' mostrata dentro lo stesso owner come secondo positive control, ma resta candidate R7 fino alla regression Facial Access e al checkpoint metodologico. Nessun ELSE e nessun comportamento missing/invalid/out-of-domain viene inferito.}
\end{ddtapropositionaccepted}

\end{paracol}
%</PAGE-CONTENT:39>
\clearpage
%<PAGE-DESC:40>Registro Base Analysis - BAReferent e BAProposition
%<PAGE-CONTENT:40>
\PageHashFooter{\PageMDfive{40}}
\section{Registro Base Analysis}

\begin{ddtacore}[title={Sintesi BA corrente}]
\fontsize{6.2}{6.8}\selectfont
Il registro raccoglie la tracciabilita' BA corrente senza introdurre campi o stati ulteriori rispetto al modello BA usato nel case study. In R25, \texttt{FR-MR01-03-03A/B} sono label di review provvisori: servono a mostrare la riallocazione prodotta dallo split senza anticipare la decisione sugli ID canonici. Gli ID delle BAProposition preesistenti restano invariati per preservare la tracciabilita'. Il riferimento in \textbf{grassetto} identifica il primo elemento documentale nel quale il BAReferent o la BAProposition viene accettato; i riferimenti non in grassetto indicano occorrenze o candidature gia' visibili nella documentazione corrente. \textbf{Identity continuity:} lo stesso nome di BAReferent nel registro denota la stessa identita' semantica in ogni occorrenza; se due significati sono distinti devono essere rappresentati da BAReferent distinti, non disambiguati implicitamente dal contesto. Le proposition sono registrate con gli operatori, i ruoli e gli scoped modifier canonici definiti dalla guida BA di riferimento.
\end{ddtacore}

\vspace{0.2mm}
{\sffamily\bfseries\color{DDTABlue}BAReferent}\par
\vspace{0.1mm}
\fontsize{5.15}{5.5}\selectfont
\renewcommand{\arraystretch}{0.82}
\setlength{\tabcolsep}{3pt}
\begin{tabularx}{\textwidth}{L{55mm}Y}
\toprule
\textbf{BAReferent} & \textbf{Riferimenti documentali} \\
\midrule
DermaTriage & \texttt{MR-01}; \textbf{\texttt{DEC-MR01-01}}; \texttt{FR-MR01-01-01}; \texttt{DEC-MR01-03}; \texttt{FR-MR01-03-01}; \texttt{FR-MR01-03-02}; \texttt{FR-MR01-03-04}; \texttt{DEC-MR01-04}; \texttt{FR-MR01-04-01}; \texttt{MR-02}; \texttt{MR-03}; \texttt{DEC-03}; \texttt{FR-03}; \texttt{FR-12}; \texttt{DEC-15}; \texttt{FR-23}; \texttt{DEC-16}; \texttt{FR-24}; \texttt{FR-25}; \texttt{MR-04}; \texttt{DEC-04}; \texttt{FR-04}; \texttt{FR-05}; \texttt{DEC-05}; \texttt{FR-09}; \texttt{DEC-09}; \texttt{FR-06} \\
DermatologicalCase & \textbf{\texttt{MR-01}}; \texttt{FR-MR01-01-01}; \texttt{FR-MR01-03-04} \\
DermatologicalTriageEvaluation & \textbf{\texttt{MR-01}}; \texttt{DEC-MR01-01}; \texttt{FR-MR01-01-01}; \texttt{DEC-MR01-03} \\
CaseUrgency & \textbf{\texttt{MR-01}}; \texttt{FR-MR01-01-01}; \texttt{FR-MR01-03-04}; \texttt{DEC-MR01-04}; \texttt{FR-MR01-04-01} \\
OperationalTriagePriority & \textbf{\texttt{MR-01}}; \texttt{DEC-MR01-04}; \texttt{FR-MR01-04-01} \\
CaseInformation & \texttt{MR-01} \\
Patient & \texttt{MR-01} \\
LesionImage & \texttt{DEC-MR01-01}; \textbf{\texttt{FR-MR01-01-01}}; \texttt{DEC-MR01-03}; \texttt{FR-MR01-03-01}; \texttt{FR-MR01-03-02} \\
SymptomInformation & \texttt{DEC-MR01-01}; \textbf{\texttt{FR-MR01-01-01}}; \texttt{FR-MR01-03-04} \\
ImageUrgencyAnalysis & \texttt{DEC-MR01-03}; \textbf{\texttt{FR-MR01-03-01}} \\
EfficientNet-B4 & \textbf{\texttt{FR-MR01-03-01}}; \texttt{MR-04}; \texttt{FR-05}; \texttt{DEC-06}; \texttt{DEC-08}; \texttt{FR-10} \\
ImageUrgencyClassification & \textbf{\texttt{FR-MR01-03-01}}; \texttt{FR-MR01-03-04} \\
ClinicalDescriptionGeneration & \texttt{DEC-MR01-03}; \textbf{\texttt{FR-MR01-03-02}}; \texttt{FR-MR01-03-03B} \\
Qwen2-VL-7B-Instruct & \textbf{\texttt{FR-MR01-03-02}} \\
ClinicalDescription & \textbf{\texttt{FR-MR01-03-02}}; \texttt{FR-MR01-03-03B}; \texttt{FR-MR01-03-04} \\
ClinicalDescriptionContentContract & \textbf{\texttt{FR-MR01-03-02}} \\
HistoricalCaseRetrieval & \texttt{DEC-MR01-03}; \textbf{\texttt{FR-MR01-03-03B}} \\
ClinicalDescriptionEmbeddingGeneration & \textbf{\texttt{FR-MR01-03-03B}} \\
ClinicalDescriptionEmbedding & \textbf{\texttt{FR-MR01-03-03B}} \\
HistoricalCaseDescription & \textbf{\texttt{FR-MR01-03-03A}} \\
HistoricalCaseEmbeddingGeneration & \textbf{\texttt{FR-MR01-03-03A}} \\
HistoricalCaseEmbedding & \textbf{\texttt{FR-MR01-03-03A}} \\
HistoricalCaseVectorIndexing & \textbf{\texttt{FR-MR01-03-03A}} \\
HistoricalCaseVectorIndex & \textbf{\texttt{FR-MR01-03-03A}}; \texttt{FR-MR01-03-03B} \\
HistoricalCaseSimilarityRetrieval & \textbf{\texttt{FR-MR01-03-03B}} \\
ChromaDB & \textbf{\texttt{FR-MR01-03-03A}}; \texttt{FR-MR01-03-03B} \\
all-MiniLM-L6-v2 & \textbf{\texttt{FR-MR01-03-03A}}; \texttt{FR-MR01-03-03B} \\
RetrievedHistoricalCases & \textbf{\texttt{FR-MR01-03-03B}}; \texttt{FR-MR01-03-04} \\
MultiSourceTriageSynthesis & \texttt{DEC-MR01-03}; \textbf{\texttt{FR-MR01-03-03B}}; \texttt{FR-MR01-03-04} \\
BioMistral-7B & \textbf{\texttt{FR-MR01-03-04}} \\
TriageConfidence & \textbf{\texttt{FR-MR01-03-04}}; \texttt{DEC-MR01-04}; \texttt{FR-MR01-04-01} \\
AITriageSynthesis & \textbf{\texttt{FR-MR01-03-04}}; \texttt{DEC-MR01-04} \\
AITriageSynthesisContentContract & \textbf{\texttt{FR-MR01-03-04}} \\
OperationalPriorityDerivation & \textbf{\texttt{FR-MR01-04-01}} \\
AdaptationLayer & \textbf{\texttt{FR-MR01-04-01}} \\
ControlledDermaTriageAdaptation & \textbf{\texttt{MR-04}} \\
PromptEvolutionPath & \textbf{\texttt{MR-04}}; \texttt{DEC-04}; \texttt{FR-04}; \texttt{DEC-05}; \texttt{FR-13}; \texttt{FR-14}; \texttt{FR-15}; \texttt{DEC-09}; \texttt{FR-06} \\
PromptEvolutionCycle & \texttt{FR-04}; \texttt{DEC-09}; \texttt{FR-06} -- candidate; scope quantificato del ciclo, identity/lifecycle autonomi non ancora sufficientemente governati \\
PromptEvolutionEvidenceSet & \texttt{DEC-09}; \textbf{\texttt{FR-06}} \\
ClassifierAdaptationPath & \textbf{\texttt{MR-04}}; \texttt{DEC-04}; \texttt{FR-05}; \texttt{DEC-05}; \texttt{FR-13}; \texttt{FR-14}; \texttt{FR-15} \\
CandidateUrgencyClassificationAdaptation & \textbf{\texttt{DEC-06}}; \texttt{FR-09}; \texttt{DEC-07} \\
ApplicableReferenceVersion & \textbf{\texttt{DEC-06}}; \texttt{FR-09}; \texttt{DEC-07} \\
CandidateEvaluationResult & \texttt{FR-09} -- candidate decomposition; result artifact non ancora governato \\
ReferenceEvaluationResult & \texttt{FR-09}; \texttt{DEC-07} -- candidate decomposition; producing experiment/result artifact non ancora governato \\
CandidateQualificationResult & \texttt{FR-09} -- candidate decomposition; output boundary non ancora chiuso \\
ApplicableComparativeAcceptanceCriteria & \texttt{FR-09}; \textbf{\texttt{DEC-07}} \\
ApplicablePostAdoptionReference & \textbf{\texttt{DEC-08}}; \texttt{FR-10} \\
PreviousAcceptableVersionOrState & \textbf{\texttt{DEC-08}}; \texttt{FR-10} \\
PostAdoptionEvaluationResult & \texttt{DEC-08}; \texttt{FR-10} -- candidate decomposition; producer/result artifact e measurement boundary non ancora governati \\
ClinicallyCorrectedOperationalTriagePriority & \textbf{\texttt{FR-08}} \\
ClassifierUrgencyTarget & \textbf{\texttt{FR-08}} \\
SpecialistDestinationIndication & \textbf{\texttt{MR-02}} \\
ClinicalReviewResult & \textbf{\texttt{MR-03}}; \texttt{DEC-03}; \texttt{FR-03}; \texttt{FR-12}; \texttt{MR-04}; \texttt{DEC-09}; \texttt{FR-06} \\
OriginalDermaTriageOutcome & \textbf{\texttt{MR-03}}; \texttt{DEC-03}; \texttt{FR-03}; \texttt{FR-12}; \texttt{FR-23} \\
HealthcareProfessional & \textbf{\texttt{MR-03}} \\
ClinicalReviewDisposition & \textbf{\texttt{FR-12}} \\
B4 & \textbf{\texttt{DEC-15}}; \texttt{FR-23}; \texttt{DEC-16}; \texttt{FR-25} \\
B4TriageOutcome & \texttt{DEC-15}; \textbf{\texttt{FR-23}} \\
ClinicianValidation & \texttt{DEC-15}; \textbf{\texttt{FR-23}} \\
ValidatedOutcome & \texttt{DEC-15}; \textbf{\texttt{FR-23}} \\
ProtectedDermaTriageOperation & \texttt{DEC-16}; \textbf{\texttt{FR-24}} \\
B4CallAuthentication & \texttt{DEC-16}; \textbf{\texttt{FR-25}} \\
DermaTriageB4Client & \textbf{\texttt{FR-25}} \\
B4BearerJWT & \textbf{\texttt{FR-25}} \\
\bottomrule
\end{tabularx}

%</PAGE-CONTENT:40>
\clearpage
%<PAGE-DESC:41>Registro Base Analysis - BAProposition
%<PAGE-CONTENT:41>
\PageHashFooter{\PageMDfive{41}}
\section*{Registro Base Analysis - BAProposition}
\begin{ddtacore}[title={Continuita' del registro}]
\fontsize{6.0}{6.6}\selectfont
Questa pagina continua il registro BA della pagina precedente. Il \textbf{grassetto} indica la prima accettazione; il \emph{corsivo} indica proposition candidate non ancora accettate. Nessun nuovo stato BA viene introdotto dalla separazione fisica del registro su due pagine.
\end{ddtacore}

\vspace{0.35mm}
{\sffamily\bfseries\color{DDTABlue}BAProposition}\par
\vspace{0.1mm}
\fontsize{4.85}{5.18}\selectfont
\renewcommand{\arraystretch}{0.82}
\begin{tabularx}{\textwidth}{L{58mm}Y}
\toprule
\textbf{BAProposition} & \textbf{Riferimenti documentali} \\
\midrule
\emph{\texttt{BAP-MR01-01}} & \texttt{MR-01} \\
\textbf{\texttt{BAP-MR02-01}} & \texttt{MR-02} -- \texttt{produce} (\texttt{DermaTriage -> SpecialistDestinationIndication}); nessun input o routing interno inferito \\
\emph{\texttt{BAP-MR03-01}} & \texttt{MR-03} -- candidate \texttt{transfer} (\texttt{HealthcareProfessional -> DermaTriage : ClinicalReviewResult}); delivery path tecnico non ancora accettato \\
\textbf{\texttt{BAP-FR12-01}} & \texttt{FR-12} -- \texttt{constrain} (\texttt{ClinicalReviewResult}: \texttt{clinicalReviewDisposition} in \texttt{[confirmation, correction]}) \\
\emph{\texttt{BAP-DEC15-01}} & \texttt{DEC-15} -- candidate \texttt{transfer} (\texttt{DermaTriage -> B4 : B4TriageOutcome}) \\
\emph{\texttt{BAP-DEC15-02}} & \texttt{DEC-15} -- candidate \texttt{transfer} (\texttt{DermaTriage -> B4 : ClinicianValidation}) \\
\emph{\texttt{BAP-DEC15-03}} & \texttt{DEC-15} -- candidate \texttt{transfer} (\texttt{B4 -> DermaTriage : ValidatedOutcome}) \\
\textbf{\texttt{BAP-FR23-01}} & \texttt{FR-23} -- \texttt{transfer} (\texttt{DermaTriage -> B4 : B4TriageOutcome}) \\
\textbf{\texttt{BAP-FR23-02}} & \texttt{FR-23} -- \texttt{transfer} (\texttt{DermaTriage -> B4 : ClinicianValidation}) \\
\textbf{\texttt{BAP-FR23-03}} & \texttt{FR-23} -- \texttt{transfer} (\texttt{B4 -> DermaTriage : ValidatedOutcome}) \\
\emph{\texttt{BAP-DEC16-01}} & \texttt{DEC-16} -- candidate \texttt{constrain} (\texttt{ProtectedDermaTriageOperation}: \texttt{authenticationMechanism} = \texttt{X-API-Key}); consolidated by \texttt{FR-24} \\
\emph{\texttt{BAP-DEC16-02}} & \texttt{DEC-16} -- candidate \texttt{constrain} (\texttt{B4CallAuthentication}: \texttt{authenticationMechanism} = bearer JWT); consolidated by \texttt{FR-25} \\
\textbf{\texttt{BAP-FR24-01}} & \texttt{FR-24} -- \texttt{constrain} (\texttt{ProtectedDermaTriageOperation}: required \texttt{authenticationMechanism} in \texttt{[X-API-Key]}) \\
\textbf{\texttt{BAP-FR25-01}} & \texttt{FR-25} -- \texttt{realize} (\texttt{B4CallAuthentication -> DermaTriageB4Client}) \\
\textbf{\texttt{BAP-FR25-02}} & \texttt{FR-25} -- \texttt{constrain} (\texttt{B4CallAuthentication}: \texttt{authenticationMechanism} in \texttt{[bearer JWT]}) \\
\emph{\texttt{BAP-FR25-03}} & \texttt{FR-25} -- candidate \texttt{transfer} (\texttt{?? -> DermaTriageB4Client : B4BearerJWT}); token source \texttt{NOT SPECIFIED} \\
\emph{\texttt{BAP-DEC01-01}} & \texttt{DEC-MR01-01} -- \texttt{produce}, \texttt{condition: LesionImage unavailable} \\
\textbf{\texttt{BAP-FR01-01-01}} & \texttt{FR-MR01-01-01} -- \texttt{produce}, \texttt{condition: LesionImage unavailable} \\
\textbf{\texttt{BAP-FR03-01-01}} & \texttt{FR-MR01-03-01} -- \texttt{produce}, \texttt{condition: LesionImage available} \\
\textbf{\texttt{BAP-FR03-01-02}} & \texttt{FR-MR01-03-01} -- \texttt{realize} \\
\textbf{\texttt{BAP-FR03-02-01}} & \texttt{FR-MR01-03-02} -- \texttt{produce}, \texttt{condition: LesionImage available} \\
\textbf{\texttt{BAP-FR03-02-02}} & \texttt{FR-MR01-03-02} -- \texttt{realize} \\
\textbf{\texttt{BAP-FR03-02-03}} & \texttt{FR-MR01-03-02} -- \texttt{constrain} \\
\textbf{\texttt{BAP-FR03-02-04}} & \texttt{FR-MR01-03-02} -- \texttt{constrain} \\
\textbf{\texttt{BAP-FR03-03-04}} & \texttt{FR-MR01-03-03B} -- \texttt{produce} (current-case embedding), \texttt{condition: LesionImage available} \\
\textbf{\texttt{BAP-FR03-03-05}} & \texttt{FR-MR01-03-03B} -- \texttt{realize} (current-case embedding generation) \\
\textbf{\texttt{BAP-FR03-03-06}} & \texttt{FR-MR01-03-03A} -- \texttt{produce} (historical-case embedding) \\
\textbf{\texttt{BAP-FR03-03-07}} & \texttt{FR-MR01-03-03A} -- \texttt{realize} (historical-case embedding generation) \\
\textbf{\texttt{BAP-FR03-03-08}} & \texttt{FR-MR01-03-03A} -- \texttt{produce} (vector index) \\
\textbf{\texttt{BAP-FR03-03-09}} & \texttt{FR-MR01-03-03A} -- \texttt{realize} (vector indexing) \\
\textbf{\texttt{BAP-FR03-03-10}} & \texttt{FR-MR01-03-03B} -- \texttt{produce} (retrieved cases), \texttt{condition: LesionImage available} \\
\textbf{\texttt{BAP-FR03-03-11}} & \texttt{FR-MR01-03-03B} -- \texttt{realize} (similarity retrieval) \\
\textbf{\texttt{BAP-FR03-03-12}} & \texttt{FR-MR01-03-03B} -- \texttt{transfer} (Stage 2 $\rightarrow$ Stage 3), \texttt{condition: LesionImage available} \\
\textbf{\texttt{BAP-FR03-03-13}} & \texttt{FR-MR01-03-03B} -- \texttt{transfer} (Stage 3 $\rightarrow$ Stage 4), \texttt{condition: LesionImage available} \\
\textbf{\texttt{BAP-FR03-04-01}} & \texttt{FR-MR01-03-04} -- \texttt{produce}, \texttt{condition: LesionImage available} \\
\textbf{\texttt{BAP-FR03-04-02}} & \texttt{FR-MR01-03-04} -- \texttt{realize} \\
\textbf{\texttt{BAP-FR03-04-03}} & \texttt{FR-MR01-03-04} -- \texttt{constrain} (sintesi -> content contract) \\
\textbf{\texttt{BAP-FR03-04-04}} & \texttt{FR-MR01-03-04} -- \texttt{constrain} (content contract strutturato) \\
\emph{\texttt{BAP-DEC04-01}} & \texttt{DEC-MR01-04} \\
\textbf{\texttt{BAP-FR04-01-01}} & \texttt{FR-MR01-04-01} -- accepted \texttt{produce}, \texttt{condition: LesionImage available}; candidate \texttt{operatorStructure.decisionRule} mostrata inline nell'owner ma non ancora ammessa \\
\textbf{\texttt{BAP-FR04-01-02}} & \texttt{FR-MR01-04-01} -- \texttt{realize} (\texttt{OperationalPriorityDerivation -> AdaptationLayer}) \\
\textbf{\texttt{BAP-FR06-01}} & \texttt{FR-06} -- accepted \texttt{produce} (\texttt{DermaTriage}: available \texttt{ClinicalReviewResult} population $\rightarrow$ \texttt{PromptEvolutionEvidenceSet}); selection semantics retained as a non-consolidated finding on \texttt{FR-06} (correction+pertinence, recency, bound 20); no additional formal structure, output ordering, read or transfer accepted \\
\textbf{\texttt{BAP-FR08-01-01}} & \texttt{FR-08} -- accepted \texttt{produce} (\texttt{DermaTriage}: corrected P-scale $\rightarrow$ classifier target); candidate \texttt{operatorStructure.decisionRule} mostrata inline nell'owner ma non ancora ammessa \\
\bottomrule
\end{tabularx}
%</PAGE-CONTENT:41>
\clearpage
%<PAGE-DESC:42>MR-01 - Inventario temporaneo transfer, contract e decisioni BA [DA CANCELLARE]
%<PAGE-CONTENT:42>
\PageHashFooter{\PageMDfive{42}}
\section{MR-01 - Inventario temporaneo transfer, contract e decisioni BA}

\begin{ddtaopen}[title={PAGINA TEMPORANEA - AUDIT DATA-PATH MR-01 - DA CANCELLARE DOPO LA RIALLOCAZIONE DEI RISULTATI}]
\scriptsize
Questa pagina e' un \textbf{worklist di ricerca}, non introduce nuovi campi L1, nuovi stati BA o nuovi operatori. I nomi dei contract sono etichette di lavoro. Ogni riga deve essere riesaminata contro la documentazione originale: il significato supportato viene riallocato nella documentazione DDTA al livello corretto; cio' che non e' stabilito resta \texttt{NOT SPECIFIED}/gap; solo dopo si riesegue la BA. Un transfer non genera automaticamente un FunctionalRequirement: l'eventuale FR deve superare i normali gate di comportamento, ownership e utilita' downstream.

\vspace{1mm}
\fontsize{5.6}{6.4}\selectfont
\renewcommand{\arraystretch}{1.06}
\setlength{\tabcolsep}{2.4pt}
\begin{tabularx}{\textwidth}{L{12mm}L{82mm}L{150mm}}
\toprule
\textbf{ID} & \textbf{Transfer candidato} & \textbf{Information/data contract e punti da verificare} \\
\midrule
\texttt{T01} & \texttt{? -> DermaTriage : LesionImage} (direct upload) & Direct triage input contract: caller/source ?, payload completo ?, formato/encoding/dimensioni ?, case correlation ?. Endpoint corrente \texttt{POST /analyze}. \\
\texttt{T02} & \texttt{B4Platform -> DermaTriage : CaseInformation} & B4 case-information contract: schema completo ?, required/optional ?, tipi/nullability ?, consultation identity e correlation ?. Verificare se \texttt{CaseInformation} richiede split semantico. \\
\texttt{T03-A} & \texttt{B4Platform -> DermaTriage : ? document reference/list} & Consultation-document discovery contract: campi della lista ?, selection rule del documento-lesione ?, binding alla consultation ?. \\
\texttt{T03-B} & \texttt{B4Platform -> DermaTriage : LesionImage} & B4 lesion-image contract: formato/media type/dimensioni ?, lifecycle temporaneo/persistenza ?, exact consultation correlation ?. Download image source-supported. \\
\texttt{T04} & \texttt{? -> ImageUrgencyAnalysis : LesionImage} & Stage-1 image-input contract: \textbf{RGB 380x380} noto; source ?, preprocessing owner ?, normalization/crop/invalid handling ?, runtime boundary ?. \\
\texttt{T05} & \texttt{? -> ClinicalDescriptionGeneration : LesionImage} & Stage-2 image-input contract: rappresentazione/format/dimensioni/preprocessing ?, source ?, relazione con la rappresentazione T04 ?, runtime boundary ?. \\
\texttt{T07} & \texttt{ImageUrgencyAnalysis -> MultiSourceTriageSynthesis : ImageUrgencyClassification} & Stage-1 result contract: \texttt{HIGH|MEDIUM|LOW} + confidence associata; tipo/range confidence, invalid/missing semantics e concrete handoff realization da verificare. \\
\texttt{T08} & \texttt{ClinicalDescriptionGeneration -> MultiSourceTriageSynthesis : ClinicalDescription} & Verificare se Stage4 consuma lo stesso \texttt{ClinicalDescriptionContentContract} riallocato con T06 o una rappresentazione distinta; correlation al caso e realization del passaggio ?. \\
\texttt{T09} & \texttt{HistoricalCaseSimilarityRetrieval -> MultiSourceTriageSynthesis : RetrievedHistoricalCases} & Historical-cases contract: top-5 noto; per ciascun caso sono documentati \texttt{patient\_id}, \texttt{diagnosis}, \texttt{urgency\_level} e \texttt{similarity\_score}. Restano da chiarire ordering esposto come contratto, identity/correlation e underfill behavior. \\
\texttt{T10} & \texttt{? -> MultiSourceTriageSynthesis : SymptomInformation} & Symptom-input contract: source ?, schema completo, required/optional, tipi/missing/invalid semantics, same-case correlation ?. \\
\texttt{T11} & \texttt{MultiSourceTriageSynthesis -> OperationalPriorityDerivation : AITriageSynthesis / CaseUrgency + TriageConfidence} & Stage-4 output contract ora esplicito: la sintesi contiene urgency + confidence usati dal mapping. La funzione di derivazione e' realizzata correntemente dall'\emph{Adaptation Layer}; i riferimenti concreti \texttt{adaptation\_layer.py} e \texttt{map\_urgency\_to\_p\_scale()} sono preservati in \texttt{FR-MR01-04-01}. Resta da verificare se il handoff trasferisca l'intera sintesi o soltanto tali valori, oltre a confidence domain/nullability/invalid handling. Reasoning/pathology non risultano input del mapping corrente. \\
\texttt{T12} & \texttt{? -> DermaTriage : SymptomInformation} & Symptom-only input contract: esempi itching/bleeding/growth-change/pain; source, schema completo, minimum evidence set e missing-field behavior ?. Nessun BAReferent \texttt{SymptomOnlyUrgencyAnalysis} viene introdotto senza identity evidence autonoma. \\
\texttt{T13} & \texttt{DermaTriage -> ? : CaseUrgency (+ confidence ?)} & Symptom-only output contract: CaseUrgency supportata; confidence ?, destination ?, applicabilita' dello stesso mapping P1-P4 ?. \\
\texttt{T14} & \texttt{DermaTriage / OperationalPriorityDerivation -> B4Platform : triage result} & B4 result-write contract: source semantica esatta ?, payload (urgency/confidence/P-scale/reasoning/pathology/specialist/SLA) ?, correlation alla consultation ?. \\
\texttt{T15} & \texttt{DermaTriage -> LocalDiagnosisStore : persisted result} & Local persistence contract: contenuto/schema, owner/necessita' governata, case correlation, retention, update/finality semantics ?. \texttt{save\_diagnosis()} e' source-supported come realization. \textbf{TODO persistence/storage condiviso con \texttt{FR-03} e \texttt{FR-10}:} distinguere transfer verso store, write/persist riuscito e associazione at-rest; non riaprire \texttt{storedIn} finche' il blocco cumulativo non viene riesaminato. \\
\bottomrule
\end{tabularx}

\vspace{1mm}
\fontsize{6.0}{6.8}\selectfont
\textbf{Ordine di review:} T06 e' stato riesaminato e riallocato stabilmente in \texttt{FR-MR01-03-02}, nel candidato runtime \texttt{FR-MR01-03-03B} e nel registro BA; per questo non compare piu' nella worklist. Proseguire da T07. T01--T05 restano riapribili finche' l'intero MR-01 non supera il gate data-path / information-contract. Eventuali transfer di pura realization (es. query ChromaDB, process/GPU handoff) restano fuori da questo primo inventario e vengono aggiunti solo se emergono come significato materialmente analizzabile.

\vspace{1mm}
\textbf{Decisione documentale temporanea - split di \texttt{FR-MR01-03-03}:}
la pagina precedente mescolava due funzioni con lifecycle e change impact separabili: (A) preparazione della base storica usata dal retrieval e (B) retrieval runtime del caso corrente fra Stage 2 e Stage 4. R25 le espone come \texttt{FR-MR01-03-03A/B} soltanto per la review; i label non sono ancora identita' canoniche. Il parent candidato resta \texttt{DEC-MR01-03}: il runtime appartiene chiaramente allo Stage 3 della pipeline a quattro stadi, mentre per la preparazione della base resta aperto il parent test. Lo split non crea un quinto stage. Le BAProposition gia' accettate sono riallocate senza rinumerazione: \texttt{06--09} al candidato A; \texttt{04--05,10--13} al candidato B. Nessun nome nato dalla BA e' usato per costruire la prosa documentale.

\vspace{0.8mm}
\textbf{Decisione metodologica temporanea - identity continuity e scoped \texttt{condition}:}
\texttt{DermatologicalTriageEvaluation} resta un unico BAReferent introdotto in \texttt{MR-01} e riusato nei contesti no-image e image-based: il riuso dello stesso BAReferent denota la stessa identita' semantica, non due significati omonimi dipendenti dal contesto. \texttt{FR-MR01-01-01} accetta una proposition \texttt{produce} con \texttt{condition -> LesionImage unavailable}. Nel percorso image-based, \texttt{condition -> LesionImage available} viene applicata soltanto alle proposition runtime il cui fatto e' governato dal percorso corrente: produzioni di Stage 1, Stage 2, Stage 3 runtime e Stage 4, piu' i transfer runtime gia' accettati. Non viene propagata meccanicamente alle proposition \texttt{realize}, ai \texttt{constrain} o alla preparazione offline della base storica: tali fatti restano veri come struttura/realization/contract indipendentemente dalla disponibilita' dell'immagine del singolo caso.\par

\vspace{0.8mm}
\textbf{Decisione metodologica temporanea - necessity test di \texttt{correlate}:}
\texttt{BAP-MR01-02} e \texttt{BAP-DEC01-02} sono rimosse dall'insieme attivo delle BAProposition candidate. Nei due passaggi sorgente, ``informazioni disponibili sul caso dermatologico'' e ``informazioni sintomatologiche disponibili sul caso'' descrivono il contesto del caso in analisi e gli input del comportamento di triage; non e' emersa una relazione governata, autonoma e interrogabile che richieda una proposition \texttt{correlate}. La rimozione non elimina i BAReferent ne' le proposition \texttt{produce} gia' candidate. Le occorrenze lessicali di ``correlation'' nella worklist restano domande aperte di binding/identity nei data contract e \textbf{non} indicano l'operatore BA \texttt{correlate}. L'operatore verra' riesaminato soltanto se evidence successiva mostrera' una relazione indipendente non ricostruibile tramite comportamento, identity e contract gia' governati. Questa nota e' temporanea e verra' cancellata a chiusura del case study dopo il consolidamento del risultato.
\end{ddtaopen}
%</PAGE-CONTENT:42>
\clearpage
%<PAGE-DESC:43>Indice di integrita delle pagine
%<PAGE-CONTENT:43>
\PageHashFooter{\PageMDfive{43}}
\section{Indice di integrita' delle pagine}

\begin{ddtacore}[title={Regola di verifica}]
Ogni pagina possiede un MD5 del proprio contenuto canonico. Il calcolo esclude il valore MD5 stampato sulla pagina, gli header/footer e i metadati di impaginazione esterni al relativo blocco di contenuto. Per la pagina dell'indice, gli hash delle pagine precedenti fanno parte del contenuto controllato; il solo hash della pagina indice viene sostituito dal token canonico \texttt{<SELF>} durante il calcolo, evitando una dipendenza circolare.
\end{ddtacore}

\vspace{0.8mm}
\begingroup
\fontsize{4.9}{5.35}\selectfont
\renewcommand{\arraystretch}{0.78}
\setlength{\tabcolsep}{2.4pt}
\renewcommand{\IntegrityRow}[3]{#1 & #2 & {\fontsize{4.45}{4.9}\selectfont\ttfamily #3} \\
}
\begin{tabularx}{\textwidth}{L{8mm}L{82mm}Y}
\toprule
\textbf{Pag.} & \textbf{Contenuto} & \textbf{MD5 contenuto} \\
\midrule
%<PAGE-INTEGRITY-ROWS>
\IntegrityRow{1}{Frontespizio e struttura del case study}{\PageMDfive{1}}
\IntegrityRow{2}{Contesto generale del progetto - Project Problem Framing}{\PageMDfive{2}}
\IntegrityRow{3}{MR-01 - Valutazione di triage del caso dermatologico}{\PageMDfive{3}}
\IntegrityRow{4}{MR-01 > DEC-MR01-01 - Valutazione di urgenza in assenza di immagine}{\PageMDfive{4}}
\IntegrityRow{5}{MR-01 > DEC-MR01-01 > FR-MR01-01-01 - Determinazione dell'urgenza su base sintomatologica}{\PageMDfive{5}}
\IntegrityRow{6}{MR-01 > DEC-MR01-03 - Percorso image-based a quattro stadi coordinati}{\PageMDfive{6}}
\IntegrityRow{7}{MR-01 > DEC-MR01-03 > FR-MR01-03-01 - Stage 1: urgenza e confidence da immagine}{\PageMDfive{7}}
\IntegrityRow{8}{MR-01 > DEC-MR01-03 > FR-MR01-03-02 - Stage 2: descrizione clinica strutturata}{\PageMDfive{8}}
\IntegrityRow{9}{MR-01 > DEC-MR01-03 > FR-MR01-03-03A [review label] - Preparazione della base storica per il retrieval}{\PageMDfive{9}}
\IntegrityRow{10}{MR-01 > DEC-MR01-03 > FR-MR01-03-03B [review label] - Stage 3: recupero di casi storici simili}{\PageMDfive{10}}
\IntegrityRow{11}{MR-01 > DEC-MR01-03 > FR-MR01-03-04 - Stage 4: sintesi multi-source}{\PageMDfive{11}}
\IntegrityRow{12}{MR-01 > DEC-MR01-04 - Separazione tra valutazione dell'urgenza e priorita operativa}{\PageMDfive{12}}
\IntegrityRow{13}{MR-01 > DEC-MR01-04 > FR-MR01-04-01 - Mapping urgenza/confidence verso P1-P4}{\PageMDfive{13}}
\IntegrityRow{14}{MR-02 - Indicazione della destinazione specialistica}{\PageMDfive{14}}
\IntegrityRow{15}{MR-03 - Gestione della validazione clinica degli esiti}{\PageMDfive{15}}
\IntegrityRow{16}{MR-03 > DEC-03 - Separazione tra esito originario e revisione clinica}{\PageMDfive{16}}
\IntegrityRow{17}{MR-03 > DEC-03 > FR-03 - Registrazione correlata della revisione clinica}{\PageMDfive{17}}
\IntegrityRow{18}{MR-03 > DEC-03 > FR-12 - Distinzione tra conferma e correzione della revisione clinica}{\PageMDfive{18}}
\IntegrityRow{19}{MR-03 > DEC-15 - Scambio della validazione clinica con B4}{\PageMDfive{19}}
\IntegrityRow{20}{MR-03 > DEC-15 > FR-23 - Write-back e retrieval della validazione B4}{\PageMDfive{20}}
\IntegrityRow{21}{MR-03 > DEC-16 - Autenticazione delle interfacce di review}{\PageMDfive{21}}
\IntegrityRow{22}{MR-03 > DEC-16 > FR-24 - X-API-Key per operazioni amministrative DermaTriage}{\PageMDfive{22}}
\IntegrityRow{23}{MR-03 > DEC-16 > FR-25 - Bearer JWT per interazione B4}{\PageMDfive{23}}
\IntegrityRow{24}{MR-04 - Adattamento controllato sulla base della revisione clinica}{\PageMDfive{24}}
\IntegrityRow{25}{MR-04 > DEC-04 - Attivazione dell'adattamento su accumulo di evidence clinica}{\PageMDfive{25}}
\IntegrityRow{26}{MR-04 > DEC-04 > FR-04 - Attivazione del percorso di evoluzione del prompt}{\PageMDfive{26}}
\IntegrityRow{27}{MR-04 > DEC-04 > FR-05 - Attivazione del percorso di adattamento classificatorio}{\PageMDfive{27}}
\IntegrityRow{28}{MR-04 > DEC-05 - Separazione dei percorsi di adattamento per capability}{\PageMDfive{28}}
\IntegrityRow{29}{MR-04 > DEC-05 > FR-13 - Indipendenza delle condizioni di attivazione tra percorsi di adattamento}{\PageMDfive{29}}
\IntegrityRow{30}{MR-04 > DEC-05 > FR-14 - Indipendenza della qualificazione dell'evidence tra percorsi di adattamento}{\PageMDfive{30}}
\IntegrityRow{31}{MR-04 > DEC-05 > FR-15 - Indipendenza dei risultati di qualificazione e accettazione tra percorsi di adattamento}{\PageMDfive{31}}
\IntegrityRow{32}{MR-04 > DEC-06 - Accettazione comparativa dell'adattamento della classificazione di urgenza}{\PageMDfive{32}}
\IntegrityRow{33}{MR-04 > DEC-06 > FR-09 - Qualificazione comparativa del candidato di adattamento classificatorio}{\PageMDfive{33}}
\IntegrityRow{34}{MR-04 > DEC-07 - Prioritizzazione asimmetrica delle metriche di qualita}{\PageMDfive{34}}
\IntegrityRow{35}{MR-04 > DEC-08 - Reversibilita dell'adattamento della classificazione in caso di degrado}{\PageMDfive{35}}
\IntegrityRow{36}{MR-04 > DEC-08 > FR-10 - Revoca di un adattamento classificatorio degradante}{\PageMDfive{36}}
\IntegrityRow{37}{MR-04 > DEC-09 - Uso di evidence recente e limitata nel percorso di evoluzione del prompt}{\PageMDfive{37}}
\IntegrityRow{38}{MR-04 > DEC-09 > FR-06 - Selezione dell'evidence recente per l'evoluzione del prompt}{\PageMDfive{38}}
\IntegrityRow{39}{MR-04 > DEC-11 > FR-08 - Derivazione del target di supervisione dalla priorita P-scale corretta}{\PageMDfive{39}}
\IntegrityRow{40}{Registro Base Analysis - BAReferent e BAProposition}{\PageMDfive{40}}
\IntegrityRow{41}{Registro Base Analysis - BAProposition}{\PageMDfive{41}}
\IntegrityRow{42}{MR-01 - Inventario temporaneo transfer, contract e decisioni BA [DA CANCELLARE]}{\PageMDfive{42}}
\IntegrityRow{43}{Indice di integrita delle pagine}{\PageMDfive{43}}
%</PAGE-INTEGRITY-ROWS>
\bottomrule
\end{tabularx}
\endgroup

\vspace{0.8mm}
{\fontsize{4.8}{5.2}\selectfont\color{DDTAGray}Quando una pagina viene modificata intenzionalmente, cambia il suo MD5. Una variazione inattesa segnala una modifica da riesaminare prima di proseguire.}
%</PAGE-CONTENT:43>

\end{document}
